% !TEX root = nc-cubic.tex
% Supporting source files live in papers/support/papers/.
\makeatletter
\def\input@path{{papers/support/}}
\makeatother
\documentclass[a4paper,11pt]{ltjsarticle}
\usepackage[margin=25mm,headheight=16pt]{geometry}
\usepackage{amsmath,amssymb,amsthm,mathtools,bm,mathrsfs}
\usepackage{booktabs,longtable,array,enumitem}
\usepackage{xcolor}
\usepackage[unicode,colorlinks=true,linkcolor=blue!40!black,citecolor=blue!40!black,urlcolor=blue!45!black]{hyperref}
\usepackage{fancyhdr}
\usepackage{microtype}
\setlength{\parindent}{1em}
\setlength{\parskip}{0.25em}
\setlist{itemsep=0.2em,topsep=0.4em}
\allowdisplaybreaks[2]
\pagestyle{fancy}
\fancyhf{}
\fancyhead[L]{\small 3次超曲面のPoisson構造と非可換化}
\fancyhead[R]{\small 2026-11-09 / v6}
\fancyfoot[C]{\thepage}
\newcommand{\CC}{\mathbb C}
\newcommand{\ZZ}{\mathbb Z}
\newcommand{\PP}{\mathbb P}
\newcommand{\OO}{\mathcal O}
\newcommand{\cK}{\mathcal K}
\newcommand{\cP}{\mathcal P}
\newcommand{\cD}{\mathcal D}
\newcommand{\cC}{\mathcal C}
\newcommand{\cE}{\mathcal E}
\newcommand{\cF}{\mathcal F}
\newcommand{\cM}{\mathcal M}
\newcommand{\fm}{\mathfrak m}
\newcommand{\Perf}{\operatorname{Perf}}
\newcommand{\Ku}{\operatorname{Ku}}
\newcommand{\MF}{\operatorname{MF}}
\newcommand{\HH}{\operatorname{HH}}
\newcommand{\QGr}{\operatorname{QGr}}
\newcommand{\qgr}{\operatorname{qgr}}
\newcommand{\Dsg}{D_{\mathrm{sg}}}
\newcommand{\RHom}{\operatorname{RHom}}
\newcommand{\REnd}{\operatorname{REnd}}
\newcommand{\Hom}{\operatorname{Hom}}
\newcommand{\Ext}{\operatorname{Ext}}
\newcommand{\Ind}{\operatorname{Ind}}
\newcommand{\Def}{\operatorname{Def}}
\newcommand{\thick}{\operatorname{thick}}
\newcommand{\Spec}{\operatorname{Spec}}
\newcommand{\Spf}{\operatorname{Spf}}
\newcommand{\Id}{\operatorname{Id}}
\newcommand{\sym}{\operatorname{sym}}
\newcommand{\Cone}{\operatorname{Cone}}
\newcommand{\rank}{\operatorname{rank}}
\newcommand{\Pic}{\operatorname{Pic}}
\newcommand{\otL}{\mathbin{\otimes^{\mathbf L}}}
\newcommand{\DocDate}{2026年11月9日}
\newcommand{\DocVersion}{論文分割版 v6}
\theoremstyle{plain}
\newtheorem{theorem}{定理}[section]
\newtheorem{proposition}[theorem]{命題}
\newtheorem{lemma}[theorem]{補題}
\newtheorem{corollary}[theorem]{系}
\theoremstyle{definition}
\newtheorem{definition}[theorem]{定義}
\newtheorem{problem}[theorem]{問題}
\newtheorem{example}[theorem]{例}
\theoremstyle{remark}
\newtheorem{remark}[theorem]{注意}
\renewcommand{\proofname}{証明}
\hypersetup{pdfauthor={},pdfsubject={3次超曲面のPoisson構造・非可換量子化・rotation付き残余圏}}

% Every displayed mathematical formula carries a numerical equation number.
\numberwithin{equation}{section}
\hypersetup{pdftitle={滑らかな三次超曲面のPoisson構造：正規形・退化軌跡・代数的量子化}}
\title{滑らかな三次超曲面のPoisson構造\\
\large 正規形・退化軌跡・代数的量子化}
\author{}
\date{2026年11月14日・論文A（投稿準備稿 v61）}
\fancyhead[L]{\small 滑らかな三次超曲面のPoisson構造}
\fancyhead[R]{\small 2026-11-14 / A-v61}
\begin{document}
\maketitle
\begin{abstract}
複素数体上の滑らかな射影三次超曲面
$X\subset\mathbb P^{n+1}$、$n\ge3$の正則Poisson構造を研究する。
大域双ベクトル場の定数三重ベクトル表示とHessianの線形系を用い、
全ての非零Poisson構造を三変数Jacobianブロックの直和に帰着する。
この正規形から、全ての許される退化階数に対するBondal予想の次元下界と、
退化スキームの方程式を導く。
三次四様体の一般階数四の構造は二つの滑らかな平面三次の分離型に限られ、
全て純楕円型の対数symplectic構造となる。
また各Poisson構造を、原点を含む一変数Zariski開集合上の
Koszul Artin--Schelter正則代数と、その正則中心三次商に量子化する。
中心元は、滑らかな三変数制限に対する普遍Hesse公式と、
古典Casimir核および普遍交換子行列の階数を用いて代数的に構成する。
この構成は三変数制限が特異または零の場合も含み、
固定した変数分解の全滑らかな係数空間上に中心元を選べる。
階数二の場合には、supportの三平面を全て動かし、三変数制限の型によらず、
全座標群に共変な中心元を、普遍Hessianの二次補正まで指定して零切断近傍に選ぶ。
証明にはGrassmann多様体上の純中心束のコホモロジー消滅と、
有限中心核のNoetherian完備化を用いる。
従って$\gcd(n+2,3)=1$なら、被約な階数二のPoisson対の全モジュライと
変形直線の積の零切断近傍へ量子化族を降ろせる。
特に三次三様体では、全ての非零Poisson対を同時に扱う。
一般次元でも三次Veroneseは全階数二層へ降下する。
全ての非零階数の被約層で、中心スカラーgerbeの位数は$\gcd(n+2,3)$であり、
これが次数一を保つ線束補正による降下の障害を測る。
中心変数が残らない満support層でも、ブロックの順序を底に加えずに、
指定二次jetを持つ共変な中心元を零切断近傍に選べる。
この場合は係数空間のアフィン性により有限jetを直接持ち上げ、
三次Veroneseを被約Poissonモジュライへ降ろす。
特に三次四様体の階数四の全Poisson対を同時に扱う。
制限が恒等的零の層では、PBW対称化した三次式自体が
全パラメータで全ブロック群に共変な正則中心元となる。
ただし$n\ge4$では、この標準中心直線を全support上の選択と一致させることは
二次jetの相対コホモロジー障害によって妨げられる。
節点・尖点層では特異点直線によって次数重みを除き、
中心変数の個数によらず、Poisson対の各層と変形直線の積の零切断近傍へ量子化族を降ろす。
また、固定した正規形のPoisson線形族を同時に量子化する。
さらに、固定超曲面上のPoissonスキームとその自己同型群による商、
超曲面自身を動かす相対モジュライを区別し、
Fermat超曲面の成分を集合分割で記述するとともに、
相対モジュライの既約成分数が$\lfloor(n+2)/3\rfloor$であることを示す。
付録では、偏極付きの量子超曲面からの二次ambientの復元と、
関連する一般的な存在判定を整理する。
\end{abstract}
\tableofcontents

\section{序論}\label{sec:paperA-introduction}

三次超曲面のPoisson構造は、その可換幾何と次数付き非可換変形を結ぶ。
本稿の中心結果は、任意次元の滑らかな三次超曲面の全Poisson構造を
三変数Jacobianブロックへ還元する正規形定理である。
この構造定理から退化軌跡とモジュライの成分を記述し、
各Poisson構造を正則中心三次商による代数族として量子化する。
以下では$X=V(F)\subset\mathbb P(V)$、$\dim V=N=n+2\ge5$とし、
基礎体は$\mathbb C$とする。
量子化では周囲の二次代数を$A$（族では$\mathcal A$）、
中心商を$B=A/(\Omega)$（族では$\mathcal B$）と書き、
ambient復元の付録でもこの記法を保つ。

\subsection*{正規形と退化軌跡}
三次式$F$を固定すると、大域双ベクトル場は
\begin{equation}\label{eq:intro-bivector-representation}
 \bigwedge\nolimits^3V\xrightarrow{\sim}H^0(X,\bigwedge\nolimits^2T_X),
 \qquad P\longmapsto\pi_P=\iota_{dF}P|_X
\end{equation}
と一意に表示される（命題\ref{prop:pg-represent}）。
Jacobi恒等式を$F$のHessianに関する線形条件へ変換し、
基点を持たない二次形式の系に対応するLie代数を解析することにより、
次の正規形を得る（定理\ref{thm:pg-normal}）。
任意の非零Poisson構造に対し、
$N=3k+m$および$V=W_1\oplus\cdots\oplus W_k\oplus U$、
$\dim W_s=3$となる分解を選んで
\begin{equation}\label{eq:intro-block-normal-form}
 F=\sum_{s=1}^k f_s(x_s,u),\qquad
 P=\sum_{s=1}^k
 \partial_{x_{s1}}\wedge\partial_{x_{s2}}\wedge\partial_{x_{s3}}
\end{equation}
と書ける。$u$はPoisson中心の座標であり、一般Poisson階数は$2k$である。
証明の核心は補題\ref{lem:pg-basepoint-free-lie}である。
基点自由な四項恒等式の形式系を持つ完全・中心なしLie代数について、
可解根基を排除し、全単純因子が$\mathfrak{sl}_2$となることを示す。
補題\ref{lem:pg-quadratic-lie}が三重ベクトルからこの状況への還元を与える。
個々の座標や純中心三次項のブロックへの配分の一意性は必要としない。

$n=3$では$k=1$となり、既知のGrassmann記述を回復する。
一般の滑らかな$F$でも非零構造の一般階数は二であるが、
特殊な$F$には複数ブロックの構造が現れる（系\ref{cor:pg-generic}）。
従って、一般の超曲面上の構造と、固定した特殊超曲面のPoissonスキームの
一般点とでは、一般階数が異なり得る。

$D_{2j}(\pi)=V(\pi^{j+1})$を階数が$2j$以下となる退化スキームとする。
正規形から、全ての$2j<n$について
\begin{equation}\label{eq:intro-bondal}
 \dim D_{2j}(\pi)\ge2j+1
\end{equation}
を得る（定理\ref{thm:pg-bondal}）。
$j<k$では、$j+1$個のブロック上の平面曲線の射影joinを構成し、
Euler方向を除いた双ベクトルの階数を計算する。
このjoinは次元$2j+1$で$D_{2j}(\pi)$に含まれる。
三変数制限$f_s(x_s,0)$の滑らかさはこの次元下界には不要である。

退化スキームの生成イデアルは命題\ref{prop:pg-deg-ideal}で与える。
各$f_s(x_s,0)$が滑らかな平面三次を定める場合には、
有限平坦critical coverと零点の完全交叉性・次数をさらに記述する
（補題\ref{lem:pg-critical-cover}、命題\ref{prop:pg-zero-ci}）。
この追加仮定を、正規形とBondal予想の定理の仮定から区別する。

三次四様体では、一般階数四の構造を持つことと
$F=f(x_0,x_1,x_2)+g(y_0,y_1,y_2)$という分離表示を持つことが同値である
（系\ref{cor:pg-cubic-fourfold-elliptic}）。
この場合、退化因子は二本の楕円曲線のjoinであり、
正規性、零点スキームおよびmodularベクトル場を直接計算して
purely elliptic log symplecticであることを示す
（命題\ref{prop:pg-split-elliptic}）。
Bricalli--Favale--Pirola\cite[Theorem~3.4]{PBricalliFavalePirola}の
Thom--Sebastiani分解の判定を合わせると、一般階数四の構造の存在は
周囲の$\PP^5$におけるHessianの階数が三以下となる軌跡が線形平面を含むこととも同値である。
Hessian行列式の可約性だけでは十分でない
（例\ref{ex:pg-reducible-hessian-no-rank-four}）。

\subsection*{代数的量子化と中心元}
定理\ref{thm:pg-quantization}は、固定した任意のPoisson構造に対し、
$0$を含むZariski開集合$U\subset\mathbb A^1_\hbar$上に
二次代数$A_U$と正則中心三次元$\Omega_U$を構成する。
各次数成分は底上で平坦であり、
$A_U$の閉ファイバーは$N$生成元、Hilbert級数$(1-t)^{-N}$、
大域次元$N$のKoszul AS正則代数となる。
中心商$B_U=A_U/(\Omega_U)$は$\operatorname{Sym}(V^*)/(F)$の平坦量子化で、
各閉ファイバーはAS Gorenstein、Hilbert級数$(1-t^3)/(1-t)^N$を持つ。
関係式は、各ブロックについて巡回的に
\begin{equation}\label{eq:intro-quantum-relations}
 [x_{s2},x_{s3}]
 =\hbar\,\operatorname{sym}(\partial_{x_{s1}}f_s)
\end{equation}
であり、異なるブロックと中心変数の間の交換子は零である。

中心三次元の構成は、以下の有限次元の議論による。
まず滑らかな三変数制限を持つブロックについて、
Hesse座標で全混合係数を許す厳密な多項式中心公式を与える
（定理\ref{thm:pg-fermat-mixed-universal}、
\ref{thm:pg-universal-hesse-mixed}）。
次に、\emph{全体の}三次式$f$が滑らかであれば、
三変数制限の特異性によらず古典Casimir核は
\begin{equation}\label{eq:intro-classical-kernel-all}
 \ker\left(S_3\xrightarrow{p\mapsto(\{p,x_i\}_f)_{i=1}^3}
 (S_4)^3\right)=\mathbb C f\oplus\mathbb C[u]_3
\end{equation}
となる。ここで$S=\mathbb C[x_1,x_2,x_3,u]$である。
滑らかな三変数制限を持つ係数の稠密開集合での中心公式から、
普遍量子交換子行列の所定の小行列式が恒等的に消える。
これと古典核の一定次元性を合わせ、Cramer公式で
$\hbar=0$において正則な有理式の中心元を得る
（定理\ref{thm:pg-universal-determinant}）。
中心項を配分して各ブロック全体を滑らかにできることは
補題\ref{lem:pg-central-splitting-smooth}で示す。
このため量子化定理には形式的中心元の存在を仮定する必要がない。

固定した変数分解の全滑らかな係数空間についても、
零切断全体を含む一つの開集合上に偶な中心元を選べる
（定理\ref{thm:pg-det-global-coefficient-lift}）。
ここでは$\hbar$で割った交換子の核を定階数のベクトル束とし、
純中心部分による商線束をfactorialな係数環上で自明化する。
異なる規格化の間の遷移は一般に
$\Omega'=a\Omega+c$（$a$は単元、$c$は純中心三次式）である。
同じ古典極限を持つ中心元でも、自然な周囲代数の同一視で
同じ中心イデアルを定めるとは限らない
（例\ref{ex:pg-det-distinct-normalized-ideals}）。
従って、係数空間上の存在とモジュライへの同変な降下は区別する。

同変な選択については、まず
$G_0=(\mathrm{GL}_3\times\mathrm{GL}_m)\times\mathbb G_m$
という変数分解を保つ群に対し、零切断全体を含む不変開上で
中心元を選べることを証明する
（定理\ref{thm:pg-equivariant-levi-center}）。
証明には既知のReynolds平均を用い、
中心元だけでなく、その定義域と古典極限も同変に保つ。
さらに$g=f(x,0)$が滑らかな部分、および節点・尖点の各軌道上では、
$g$の方向微分が張る二次式の部分束を同変に分裂し、
混合二次項を用いてactive座標の中心を選ぶ。
この操作により$X=Mx+Cu,\ U=Nu$という
中心変数による平行移動を含む群に対しても同変な中心元を得る
（定理\ref{thm:pg-equivariant-parabolic-center}）。
特異な二軌道では、基準形の安定化群と同変な左逆を明示し、
軌道上へ正則に移す（補題\ref{lem:pg-equivariant-singular-retractions}）。
Poisson係数と変形変数を分けると、変形変数を固定したまま
対応する商スタック上に中心商の族が降りる
（系\ref{cor:pg-equivariant-central-quotient}）。
また$g=0$ならJacobian括弧は中心係数環上でactive変数についてaffine-linearとなり、
PBW対称化の随伴同変性から
\begin{equation}\label{eq:intro-zero-restriction-center}
 \Omega_{f,t}=\sym(f),\qquad
 [x_i,\sym(p)]=t\,\sym(\{x_i,p\}_f)
\end{equation}
を得る。この場合は全$t\in\mathbb A^1$でPBW性と全ブロック群への共変性が成り立ち、
例外パラメータの除去も純中心係数の変更も不要である
（定理\ref{thm:pg-equivariant-zero-center}）。
これらの層ごとの構成とは別に、階数二の全supportにわたる共変な選択を
定理\ref{thm:pg-intrinsic-rank-two-center}で証明する。
その証明は中心調整を経由せず、可約・非被約・零の三変数制限も含む。
$g$が滑らかな場合、全ブロック群$G$による商は、
階数二のPoisson対の被約層内の対応する開部分上の
$\mathbb G_m$-gerbeである
（命題\ref{prop:pg-equivariant-poisson-gerbe}）。
$\gcd(m,3)=1$なら指標を用いて次数線形化を補正でき、
次数付き代数の族がこの被約層へ降りる
（系\ref{cor:pg-equivariant-character-descent}）。
特に三次三様体では$m=2$なので、十七次元の被約Poissonモジュライの稠密開上に
この族を得る。
ここで十七は底の次元であり、固定した非零変形パラメータでの
非可換代数の同型類の像の次元を主張するものではない。
零制限の層はsupport平面が$X$に含まれる条件であり、$m\ge3$のときに限り非空となる。
この層でも同じgerbeと指標補正を用い、$\gcd(m,3)=1$なら
全変形直線上の代数族をPoissonモジュライへ降ろせる
（命題\ref{prop:pg-equivariant-zero-stratum}）。

節点・尖点の各層は階数二層内でそれぞれ余次元$1,2$である。
ここでは一意な特異点の直線$\ell_g\subset W$が正則に動く。
引戻しの線形化で$K_{\mathrm{sc}}$が$\ell_g$に重み$-1$で作用するため、
\begin{equation}\label{eq:intro-singular-point-descent}
 \mathcal A_d^{\mathrm{sp}}=\mathcal A_d\otimes\ell_g^{\otimes d},
 \qquad
 \mathcal B_d^{\mathrm{sp}}=\mathcal B_d\otimes\ell_g^{\otimes d}
\end{equation}
は重み零になる。
従ってこの二層では$\gcd(m,3)=1$という条件を要せず、
各層と変形直線の積の零切断近傍へ実際の次数付き代数族が降りる
（命題\ref{prop:pg-equivariant-singular-strata}）。
三次三様体ではその底の次元はそれぞれ$16,15$である。
これらの明示的な層ごとの選択が互いに一致するとは限らない。
型の境界をまたぐ一つの族は、次の内在的な構成によって選ぶ。

$C_3^\times(V)\subset\bigwedge^3V$を非零分解可能三重ベクトルの錐とし、
$\mathcal T_N$を全三次式が滑らかな係数空間とする。
$\mathcal S_N=\mathcal T_N\times C_3^\times(V)$では
support $W_P$が三平面の全体を動く。
古典Casimir核は自然に
\begin{equation}\label{eq:intro-intrinsic-classical-kernel}
 \mathcal K_0=\OO_{\mathcal S_N}F\oplus
 \operatorname{Sym}^3(V/W_P)^*
\end{equation}
と分裂する。Grassmann多様体上のBorel--Weil--Bott計算により
$H^1(\mathcal S_N,\mathcal K_0)=0$であるため、
普遍二次Hessian補正を初期値として$\hbar$進の中心核へ順次持ち上げられる。
Stiefel開上の正則関数の延長によって係数を有限生成の不変式環に置き、
有限加群の完備化を$\hbar^3$で還元し、指定二次jetを持つ\emph{代数的な}中心元を得る。
最後に還元的群$\mathrm{GL}(V)\times\mathbb G_m$で平均する。
この方法で全$\mathcal S_N\times\{0\}$を含む不変開上の共変な中心元と、
次数別局所自由な正則中心商族を構成する
（定理\ref{thm:pg-intrinsic-rank-two-center}）。
形式的級数の全体を代数的な式と同一視するのではなく、
中心性と指定された有限jetを持つ有限核の元を選ぶ点が要となる。

さらに$\gcd(N,3)=1$なら、$Nr-3s=1$を満たす整数を使った
指標$\theta(L,\lambda)=(\det L)^r\lambda^s$で次数重みを補正し、
被約な階数二のPoisson対の\emph{全}モジュライと変形直線の積の
零切断近傍へ代数族を降ろす
（系\ref{cor:pg-intrinsic-rank-two-descent}）。
三次三様体では$N=5$なので、十七次元の非零Poisson対の全被約モジュライが底となる。
これは変形原点の近傍における存在であり、全変形直線への延長や一意性は主張しない。
一般の$N$でも、次数$3d$部分を$\lambda^d$で捩る三次Veronese
$\mathcal A^{[3]}_d=\mathcal A_{3d}\otimes\lambda^d$とその中心商は、
全階数二層の零切断近傍へ降下する
（系\ref{cor:pg-intrinsic-third-veronese-descent}）。
その古典偏極は各ファイバー上で$\OO(3)$である。
一方、$\operatorname{Pic}(\mathcal S_N)=0$と同変可逆層の重みを計算すると、
中心スカラーgerbeの位数は正確に$\gcd(N,3)$となる
（命題\ref{prop:pg-intrinsic-gerbe-period}）。
従って$3\mid N$では、指標に限らず、可逆層の次数冪による補正によって
元の次数一の族を降ろすことはできない。

零制限層の標準中心直線との一致には、これとは別の障害がある。
$N\ge6$では、この層上で$\mathscr H(P,F)$は非零の純中心切断になり得るが、
全support上の純中心束の$P$重み二の大域切断はBott計算によって零である。
従って$\Omega^{\mathrm z}=\sym(F)$への全層での一致は、
任意の大域的な偶中心元について、同変性を要求しなくても二次jetで不可能となる
（命題\ref{prop:pg-intrinsic-zero-matching-obstruction}）。
これは所定の周囲代数における中心直線の比較であり、中心商の抽象的な非同型は意味しない。
全係数での補空間の同変選択は実際には不可能である
（命題\ref{prop:pg-equivariant-global-centering-obstruction}）が、
この障害は中心核そのものの持上げを妨げない。

固定した正規形では、独立なパラメータ$\eta_1,\ldots,\eta_k$を導入し、
各ブロックの中心元を加えることで
\begin{equation}\label{eq:intro-independent-center}
 \Omega_{\boldsymbol\eta}=\sum_{s=1}^k\Omega_{f_s,\eta_s}
\end{equation}
を得る。ブロック間の追加補正は不要である。
$\eta_s=\lambda_s\hbar$により、$\sum_s\lambda_s\pi_s$という
Poisson線形族の全方向を一つの相対代数族で量子化できる
（定理\ref{thm:pg-all-dimensions}、\ref{thm:pg-ad-all-poisson}）。
滑らかな三変数制限では高々$\eta_s^4$の多項式公式を用い、
零制限では上のPBW対称化を、非零の特異な制限を含む一般の場合には
原点で正則な有理式を用いる。
また、定数三重ベクトルとHessianの縮約
\begin{equation}\label{eq:intro-intrinsic-hessian}
 \mathscr H(P,F)=\frac16P^{ijk}P^{abc}F_{ia}F_{jb}F_{kc}
\end{equation}
により、中心元の二次補正を内在的に表す
（命題\ref{prop:pg-universal-casimir-jet}、
\ref{prop:pg-ad-intrinsic-hessian}）。
二次jetの普遍恒等式とHesse表示での全次数中心公式は、異なる主張である。

中心変数が残らない$N=3k$の満support層では、固定した正規形を離れても
一つの共変な中心元を選べる
（定理\ref{thm:pg-intrinsic-full-support-center}）。
まず、三重ベクトルから復元される順序のない$k$個のブロックを
有限\'etale被覆として相対化する
（補題\ref{lem:pg-intrinsic-block-cover}）。
満supportの三重ベクトルの軌道は
$\mathrm{GL}_N/(\mathrm{SL}_3^k\rtimes\mathfrak S_k)$であり、
既知のアフィン商の定理によってアフィンである。
適合する滑らかな三次式の係数空間も、ベクトル束の判別式主開としてアフィンになる。
この底上で割った交換子の核を階数$k$の束とし、
$\sym(F-\hbar^2\mathscr H(P,F)/12)$を法$\hbar^3$の切断として直接持ち上げる。
還元群と偶性による平均はこのjetを保ち、
全零切断を含む不変開上に代数的な中心元を与える。
ブロックごとの次数によるfiltrationから正則性を示すので、
テンソル積代数全体のdomain性を仮定しない。
得られる三次Veroneseと中心商は、実際の被約Poisson層へ降下する
（系\ref{cor:pg-full-support-veronese-descent}）。
満supportの底の次元は$2k$であり、三次四様体の階数四では四次元となる。
これはPoisson対の底の次元であり、固定した非零$\hbar$での
非可換同型類の像の次元ではない。

また、中心スカラーgerbeの位数$\gcd(N,3)$は階数二に限らない。
命題\ref{prop:pg-intrinsic-all-rank-gerbe}では、各非零階数$2k$の全被約層で
これを証明する。$3\mid N$の場合、Fermat点の安定化群内の二つの元の交換子が
中心の三乗根となることから、全ての同変可逆層の中心重みが三の倍数だと分かる。
この証明は底のPicard群の消滅を必要としない。
従って満support層の三次Veronese降下は、
元の次数一の族の可逆層による補正とは区別される。

この量子化は各構造、固定した正規形の線形族、および固定した変数分解の
係数空間上での構成を含み、上記の範囲では座標変換との同変性も備える。
次数付き代数を載せる商スタックでは、次数$d$に重み$d$で作用する
スカラーの$\mathbb G_m$をまず保持し、指標補正・特異点直線による捻り、または三次Veroneseによって
上記の範囲においてこれを除く。
中心変数を残す一般の複数ブロック層での同時選択や、
全閉点での$\operatorname{qgr}B_\hbar$のsmoothnessは結論に含めない。

\subsection*{固定・相対モジュライとFermatの成分}
固定した$X$上のJacobi零スキームと、その自己同型群による商を
\begin{equation}\label{eq:intro-poisson-space}
 \operatorname{Pois}(X)=
 \{\pi\in H^0(X,\bigwedge\nolimits^2T_X):[\pi,\pi]=0\},
\end{equation}
\begin{equation}\label{eq:intro-poisson-fixed-stack}
 \mathfrak{Pois}(X)=[\operatorname{Pois}(X)/\operatorname{Aut}(X)],
 \quad \operatorname{Aut}(X,\pi)=
 \{\phi\in\operatorname{Aut}(X):\phi_*\pi=\pi\}
\end{equation}
と書く。前者をPoisson構造のspace、後者を固定モジュライstackと呼ぶ。
Poisson構造のスカラーは割らず、零構造も含める。
超曲面自身を動かす相対モジュライは
\begin{equation}\label{eq:intro-poisson-relative-stack}
 \mathfrak{Pois}^{\mathrm{rel}}_n=[\mathscr P_n/\mathrm{PGL}_{n+2}]
\end{equation}
である。$\mathscr P_n$は滑らかな三次超曲面のパラメータ空間上の
相対Jacobi零スキームである。方程式の代表を使う表示では
$(F,P)\sim(cF,c^{-1}P)$とし、$\pi_P$を不変に保つ。
命題\ref{prop:pg-fixed-stack-algebraic}、
定理\ref{thm:pg-relative-stack-algebraic}により、
両モジュライは有限型Deligne--Mumford代数的スタックとなる。

固定した一般の滑らかなcubicではspaceも固定stackも既約だが、
特殊な複素構造では既約成分が増える。
定理\ref{thm:pg-fermat-relative-count}ではFermat超曲面
\begin{equation}\label{eq:intro-fermat-equation}
 X_N=V(F_N)\subset\mathbb P^{N-1},\qquad F_N=\sum_{i=1}^Nz_i^3
\end{equation}
について、各部分の大きさが三以上の非順序集合分割
$\mathcal B=\{B_1,\ldots,B_k\}$と
$V_B=\langle e_i:i\in B\rangle$を用い、
\begin{equation}\label{eq:intro-fermat-components}
 Z_{\mathcal B}=\sum_{B\in\mathcal B}\widehat{\operatorname{Gr}}(3,V_B),
 \qquad \operatorname{Pois}(X_N)_{\mathrm{red}}
 =\bigcup_{\mathcal B}Z_{\mathcal B}
\end{equation}
と記述する。$\widehat{\operatorname{Gr}}$は原点を含むPl\"ucker錐であり、
$Z_{\mathcal B}$がspaceの相異なる既約成分である。
$k$個の部分を持つ成分の一般Poisson階数は$2k$となる。
従ってspaceの成分数は
\begin{equation}\label{eq:intro-fermat-space-count}
 \begin{aligned}
 \#\operatorname{Irr}(\operatorname{Pois}(X_N))
 &=\sum_{k=1}^{\lfloor N/3\rfloor}
 \left\{\begin{smallmatrix}N\\k\end{smallmatrix}\right\}_{\ge3},\\
 \sum_{N\ge0}\frac{t^N}{N!}
 \left\{\begin{smallmatrix}N\\k\end{smallmatrix}\right\}_{\ge3}
 &=\frac{(e^t-1-t-t^2/2)^k}{k!}.
 \end{aligned}
\end{equation}
一方、$\operatorname{Aut}(X_N)=(\mu_3^N/\mu_3)\rtimes\mathfrak S_N$は
同じ部分の大きさを持つ集合分割を同一視する。
固定stackの成分数は、各部分が三以上の整数分割数$p_{\ge3}(N)$となる。
相対モジュライでは一般Poisson階数ごとに一つの成分があり、
\begin{equation}\label{eq:intro-poisson-stack-counts}
 \begin{aligned}
 \#\operatorname{Irr}(\mathfrak{Pois}(X_N))
 &=p_{\ge3}(N)=[q^N]\prod_{b\ge3}(1-q^b)^{-1},\\
 \#\operatorname{Irr}(\mathfrak{Pois}^{\mathrm{rel}}_n)
 &=\left\lfloor\frac{n+2}{3}\right\rfloor.
 \end{aligned}
\end{equation}
例えば次のように、三種類の成分数は異なる。
\begin{equation}\label{eq:intro-poisson-count-table}
 \begin{array}{c|rrrrr}
 n&3&4&6&7&10\\\hline
 \#\operatorname{Irr}(\operatorname{Pois}(X_N))&1&11&92&491&60105\\
 \#\operatorname{Irr}(\mathfrak{Pois}(X_N))&1&2&3&4&9\\
 \#\operatorname{Irr}(\mathfrak{Pois}^{\mathrm{rel}}_n)&1&2&2&3&4
 \end{array}
\end{equation}
ここではFermatが固定超曲面の成分数を最大化するとは主張しない。
また、滑らかなcubic fivefoldの固定Poissonスキームに、
節点型・尖点型の三変数制限を持つ成分が現れる具体例を
定理\ref{ne:main}に与える。
その成分性は、直接Schouten線形化でも検証できる整数小行列式で示す。

\subsection*{ambientの復元}
三次以上の一つの斉次関係で二次代数を割っても、次数一・二の積は変わらない。
この事実から量子超曲面の二次包絡としてambientを復元し、
商写像の三次核から中心方程式の直線を得る
（定理\ref{thm:pg-ambient-recovery-general}）。
左右noetherianなAS Gorenstein性と所定のinjective dimensionの下では、
偏極付き三つ組からsection algebraを復元して同じ構成を適用できる
（定理\ref{thm:pg-ambient-recovery-polarized}）。
三次三様体については、Serre双対性・二変数中心・中心上の自由性などを仮定して、
構造対象付き導来圏から偶数Veroneseを経由する復元も示す
（系\ref{cor:pg-ambient-derived-threefold}）。
一般次元でSerre関手から得られるVeroneseが元の代数を決定するかは別の問題である。

\subsection*{先行研究との関係と各定理の位置づけ}
以下では、証明に用いる既知の入力と、本稿で証明する追加部分を
定理ごとに区別する。既知の中心元の存在や低次元の分類自体は、
本稿の新規成果として数えない。

\paragraph{正規形：射影的持上げ、半単純核計算、根基排除。}
射影空間上のPoisson構造の斉次divergence-free持上げは
Bondal\cite[\S4, pp.~11--12]{PBondal93}による。
本稿の命題\ref{prop:pg-represent}とその後のJacobiatorの議論では、
三次超曲面上での消滅をEuler--Koszul表示と低次数syzygyの消滅により
周囲の錐上での消滅へ移す。
半単純Lie代数について、式\eqref{eq:pg-lie-H}の核の記述は
Derdzinski--Gal\cite[(1.4), Theorem~B]{PDerdzinskiGal}に含まれ、
Cartan三形式からの単純ブロックの内在性も同論文Lemma~7.1(a)にある。
定理\ref{thm:pg-normal}で追加されるのは、
滑らかな三次式のHessianが基点を持たないことを使い、
\emph{半単純性を仮定する前に}可解根基を排除する段階と、
これを全Poisson構造の三変数ブロック分解へ結びつける段階である。
この根基排除は補題\ref{lem:pg-basepoint-free-lie}として独立に述べ、
用いる表現論の仮定と参照箇所をそこでまとめる。

\paragraph{低次元の既知分類と高次元の幾何学的帰結。}
三次三様体の分類はLoray--Pereira--Touzet
\cite[Theorems~7.1, 9.1]{LPTPoissonCubic}による。
Mayanskiy\cite[Theorem~1]{PMayanskiy}はこの帰属を明記し、
Jacobi方程式によるGrassmannスキームの記述とPoisson cohomologyの計算を与える。
本稿の追加は、高次元で現れる複数ブロックも含む必要十分な正規形、
固定した特殊超曲面の成分、および相対モジュライの成分の記述である。
Bondal予想については、Gualtieri--Pym
\cite[Theorem~7.13]{BondalGualtieriPym}のFano四次元の場合と、
Druel--Pereira--Pym--Touzet
\cite[Corollary~1.3]{BondalDruelPereiraPymTouzet}のFano五次元の場合は既知である。
Aza\cite[Theorem~1.3]{BondalAza}の結果は、隣接する退化軌跡の
真包含の下で少なくとも一方が次元下界を満たすという選言である。
本稿の定理\ref{thm:pg-bondal}は三次超曲面というクラスについて
\emph{任意次元の全退化階数}を同時に扱い、下界を担うjoinを具体的に構成する。

\paragraph{一ブロックの量子化：既知の存在定理と相対的規格化。}
三生成元同次代数のAS正則性・3-Calabi--Yau性の判定には
Mori--Smith\cite[Theorems~1.4, 4.2]{PMoriSmith}を用いる。
このうちTheorem~1.4はArtin--Tate--Van den Berghの判定の再掲である。
低次potentialによるPBW変形はBerger--Taillefer
\cite[Theorem~1.1(i), p.~242]{PBergerTaillefer}の結果であり、
命題\ref{prop:pg-PBW}では中心係数環とactive次数の濾過に即して適用する。

さらに、三変数Calabi--Yau最高次potentialに低次項を加えた代数に
非スカラー中心元が存在することは、Etingof--Ginzburg
\cite[Theorem~3.3.2(ii), p.~1385]{EG}が既に証明している。
その証明は同論文\S8.3, pp.~1410--1411で、形式的量子化、Rees構成と特殊化を用いる。
またProposition~8.1.1(i), pp.~1407--1408は、正の重み$(a,b,c)$、$d=a+b+c$と
$\deg\Phi_j\le d$を満たす$\Phi=xyz-yxz+\sum_{j\ge1}\hbar^j\Phi_j$に対し、
古典剰余$(\Phi_1)^{\mathrm{ab}}$を持つ次数$d$以下の形式中心元を与える。
この項には孤立特異点の仮定はなく、その仮定は同命題(ii), (iii)で追加される。
証明には相対係数環と還元群の作用に対応した形式性も既に現れる。
従って、指定した古典剰余の形式中心の存在や、相対的・同変な形式構成自体も
本稿の新規性の根拠とはしない。
従って、一ブロックに中心元が存在することや、一般点から特殊な係数へ
中心元を延長するという原理自体を本稿独自の結果とはしない。
本稿の定理\ref{thm:pg-universal-determinant}と
定理\ref{thm:pg-det-global-coefficient-lift}が与えるのは、
\emph{指定した古典Casimir $f$}を極限に持つ中心元の、
独立な中心混合係数を含む係数環上の選択と、$\hbar=0$における正則性である。
その証明では、古典核の一定次元性と普遍交換子行列の小行列式を用い、
三変数制限が特異・零の場合も同じ有限線形代数で扱う。
中心元の存在、左右正則性、および中心商の平坦性は別々に証明する。
零制限については、定理\ref{thm:pg-equivariant-zero-center}で
PBW対称化が随伴作用と可換する既知の原理
\cite[Corollaries~13.7--13.8, p.~75]{PEtingofLie}を用いて、
全パラメータの中心元を直接構成する。
この原理自体を新規とせず、中心混合係数を動かす全ブロック群への共変性と
support平面を含む層上の族への適用を明示する。

\paragraph{明示中心公式と全ブロックへの組立て。}
純Sklyanin中心元の基本公式は既知である
\cite{SmithTate,DeLaetLeBruyn}。
Etingof--Ginzburg\cite[\S9, p.~1414]{EG}には、
一変数低次項を許す中心元の明示公式も収められている。
一般Hesse表示の混合二次項は、Hesse係数を可逆にすると
中心平行移動でこの表示へ移せる
（注意\ref{rem:pg-hesse-eg-translation}）。
従って、混合係数を含むという事実だけを公式の新規性の根拠とはしない。
定理\ref{thm:pg-universal-hesse-mixed}では固定した係数環のまま、
この平行移動による極を導入しない規格化と普遍恒等式を与える。
有限の次数四線形代数が証明する中心性と、
全次数のPBW性を明確に区別する。

節点型についても、Chekhov--Mazzocco--Rubtsov
\cite[Lemma~4.6, 式(4.39), p.~22]{PCMR}が、
一般化Etingof--Ginzburg代数の$\gamma=0$の場合の中心三次元を明示している。
同じaccepted versionのTable~5, p.~41には、$xyz+x^3+y^3$に
対角二次項・一般一次項を加えた古典式との対応がある。
注意\ref{rem:pg-cmr-nodal-comparison}では、一般の二次混合項をこの表示へ移す
安定化群同変な左逆、対称Jacobian関係とのパラメータ対応、
古典極限を$f$にする規格化と純中心補正を明示する。
節点型の中心元の存在と基準座標での公式を本稿の新規成果とはしない。

全次元の構成は、この一ブロックの入力を正規形定理と組み合わせるものである。
定理\ref{thm:pg-quantization}では中心項の配分、ブロックの組立て、
最高active記号による中心正則性を証明し、
定理\ref{thm:pg-all-dimensions}と\ref{thm:pg-ad-all-poisson}では
独立な多パラメータを保持する。
同変性については定理\ref{thm:pg-equivariant-levi-center}、
\ref{thm:pg-equivariant-parabolic-center}、\ref{thm:pg-equivariant-zero-center}で
所定の群への同変な選択を、系\ref{cor:pg-equivariant-character-descent}と
命題\ref{prop:pg-equivariant-zero-stratum}、\ref{prop:pg-equivariant-singular-strata}で
所定の条件下の降下を証明する。
節点・尖点軌道への移送には既知の軌道商の理論
\cite[Propositions~7.12, 7.17]{PMilneGroups}を、
特異点直線の捻りにはgerbe上の重み零の層の既知の降下
\cite[Proposition~4.10, p.~1487]{PBerghSchnurerGerbes}を用いる。
これらの一般原理自体を新規性の根拠とはしない。
定理\ref{thm:pg-intrinsic-rank-two-center}では、既知のBorel--Weil--Bott消滅、
不変式環の有限生成、Noetherian完備化の完全性を、
全support上のCasimir核と有限量子交換子行列に適用する。
追加されるのは、三変数制限の型を限定しない共変な中心元の選択と、
系\ref{cor:pg-intrinsic-rank-two-descent}の範囲での全階数二層への降下である。
ここでは、点ごとの非スカラー中心元や所定の形式中心の存在に加えて、
全supportを動かす底の上で古典剰余を固定した代数的な大域切断を選び、
全$\mathrm{GL}_N\times\mathbf G_m$への共変性を課している。
今回確認した上記文献の箇所から、この同時選択と降下が直接与えられるわけではない。
また、重み零の降下、Veronese同値、および底の可逆層による次数別の捩りは
既知の一般原理である\cite{PBerghSchnurerGerbes,PSmithVeronese,PVdBNoncommP1Bundles}。
本稿ではそれらを全supportの族へ適用し、さらに階数二の同変Picard群と、
全非零階数層の安定化群を用いてgerbeの位数を決定する。
零制限の標準中心直線への非延長も別に計算する。
満supportの構成では、既知の還元群によるアフィン商の定理
\cite[\S2, Proposition~2.2, Theorem~2.1]{PArzhantsevMatsushima}と
有限射影加群の切断の持上げを用いる。
ブロックの点ごとの内在性は補題\ref{ne:blocks}に既に含まれ、
補題\ref{lem:pg-intrinsic-block-cover}はその族としての表示である。
定理\ref{thm:pg-intrinsic-full-support-center}で加わる内容は、
順序を選ばない全満support層における指定二次jet付きの代数的な同時選択と、
系\ref{cor:pg-full-support-veronese-descent}の中心商族である。
ただし、関連文献全体に対する新規性の最終確認は残る。
一般の形式的量子化とunimodularityに関する
Kontsevich・Dolgushevの理論\cite{PKontsevich,PDolgushev}は背景であり、
本稿の有限Cramer構成の証明入力ではない。

\paragraph{分離型四様体と持上げ障害。}
Bricalli--Favale--Pirola\cite[Theorem~BおよびTheorem~3.4]{PBricalliFavalePirola}は、
滑らかな三次式のThom--Sebastiani分解をHessian階数軌跡内の線形空間で特徴付けた。
系\ref{cor:pg-cubic-fourfold-elliptic}の第二項と第三項の同値はこの既知定理であり、
第一項との同値は本稿のPoisson正規形を用いる。
例\ref{ex:pg-reducible-hessian-no-rank-four}では、
定数symmetrizerの階数三の冪等元が存在しないことにより、
Hessian行列式の可約性だけでは十分でないことを具体化する。
この線形代数の議論は、特定の座標表示で混合項がないことを調べるだけでなく、
任意の線形座標変換による$(3,3)$分解を排除する。

分離型三次四様体の直接計算
（命題\ref{prop:pg-split-elliptic}）と、
非連結な中心の高次消失を保つ持上げ障害
（命題\ref{prop:pg-triple-lifting-obstruction}、
系\ref{cor:pg-split-lifting-obstruction}）を独立に証明する。
持上げ命題は、全連結成分が正次元の滑らかな中心と任意の消失次数$r\ge2$に対し、
超曲面と同じ次数の切断の障害を$\mathbb C^{\pi_0(C)}/\mathbb C$として記述する。
文献との比較は、本文を確認したPym
\cite{PPymElliptic}のarXiv:1507.05668v1に限定する。
同版のDefinition~4.2, p.~17、Lemma~5.5, pp.~28--29、
Proposition~5.7, p.~30との関係は
注意\ref{rem:pg-split-elliptic-lift-obstruction}で詳述する。
刊行版・関連訂正文の本文との最終照合は残るため、
刊行済み定理の訂正を確定したとは述べない。

\paragraph{付録の位置づけ。}
二次包絡によるambient復元は低次数の核の議論であり、
偏極付き圏のsection algebra復元はArtin--Zhangの理論
\cite{PAZAmbient}の適用として記す。
Serre関手からの復元には中心・自由性・滑らかさ等の追加仮定が必要である。
完全交叉の存在判定はBr\"uckmann型消滅の既知の帰結
\cite[Lemma~4.1]{FuHwangCI}であり、必要な低twist消滅は
補題\ref{lem:ci-low-twist-vanishing}で独立にも証明する。
これらは応用と背景を整理するための付録とし、中心的な新規性とは位置づけない。

\subsection*{証明の依存関係と本文の構成}
主要な依存関係は次の通りである。表の各行で、追加仮定を要する帰結は
それを用いない主定理と区別する。
\begin{center}
\small
\begin{tabular}{p{.14\linewidth}p{.36\linewidth}p{.38\linewidth}}
\toprule
段階 & 入力と証明 & 結論と範囲 \\
\midrule
正規形 & Euler--Koszul表示、Hessianの基点自由性、
補題\ref{lem:pg-basepoint-free-lie} &
定理\ref{thm:pg-normal}。全三次式の滑らかさを仮定 \\
幾何学的帰結 & 正規形、射影化後の階数、joinと階数層 &
全退化階数の下界、固定・相対成分。critical coverには三変数制限の追加仮定 \\
一ブロック & 命題\ref{prop:pg-PBW}、普遍Hesse恒等式、古典核とCramer公式 &
指定した古典極限を持つ相対中心元。有限次数検算と全次数PBW性を分離 \\
全ブロック & 中心項の配分、中心元の和、最高記号の正則性 &
定理\ref{thm:pg-quantization}、\ref{thm:pg-ad-all-poisson}。商の平坦性まで証明 \\
層ごとの構成 & 不変開、Reynolds作用素、軌道上のcentering、PBW対称化、重みの補正 &
零制限では全パラメータ、節点・尖点では全$m$の層への降下 \\
階数二の全体 & 純中心束のBott消滅、有限核の完備化、全座標群の平均 &
型をまたぐ中心商族。三次Veroneseは全次元で降下し、次数一の障害は位数$\gcd(N,3)$ \\
満supportの全体 & 有限\'etaleブロック被覆、還元安定化群のアフィン商、有限jetの持上げ &
順序を選ばない共変中心商族と三次Veronese降下。全階数層のgerbeの位数は別に安定化群で計算 \\
\bottomrule
\end{tabular}
\end{center}

第\ref{sec:pg-extension}--\ref{sec:pg-normal}節で正規形を証明し、
第\ref{sec:pg-leaves}--\ref{sec:pg-bondal}節で幾何学的帰結を扱う。
第\ref{sec:pg-poisson-stacks}--\ref{ne:examples}節に
固定・相対モジュライと非楕円成分を置く。
量子化は第\ref{sec:pg-quantization-foundations}節の一ブロックの入力から始め、
第\ref{sec:pg-fermat-mixed-universal}--\ref{sec:pg-universal-determinant}節で
具体的・普遍的な中心元を構成する。
第\ref{sec:pg-quantization}--\ref{sec:pg-all-dimensions}節で全ブロックへ組み立て、
第\ref{sec:pg-equivariant-central-lifts}節で層ごとの同変化と降下を扱い、
第\ref{sec:pg-intrinsic-rank-two-center}節で階数二の全supportへ拡張する。
第\ref{sec:pg-full-support-quantization}節では全階数層のgerbeと、
中心変数が残らない全満support層の同時量子化を扱う。
その後に具体例を置き、付録に既知中心公式、ambient復元、完全交叉の背景を収める。

\section{大域的双ベクトル場と射影空間への延長}\label{sec:pg-extension}
本節から\S\ref{sec:pg-special}までは $\CC$ 上で議論する。
$V$ を次元 $N=n+2\geq5$ のベクトル空間、$S=\operatorname{Sym}(V^*)$ とし、
$F\in S_3$ が滑らかな超曲面 $X=V(F)\subset\PP(V)=\PP^{N-1}$ を定めるとする。
$\PP(V)$ は $V$ の一次元部分空間をパラメータ化する。
座標を $z_1,\ldots,z_N$、Euler 場を $E=\sum_i z_i\partial_{z_i}$ と書き、
$F_i=\partial_iF$、$F_{ij}=\partial_i\partial_jF$ と略記する。
縮約の符号は
\[
\iota_\eta(v_1\wedge\cdots\wedge v_r)
 =\sum_{a=1}^r(-1)^{a-1}\eta(v_a)
 v_1\wedge\cdots\wedge\widehat v_a\wedge\cdots\wedge v_r
\]
とする。Schouten 括弧に関して $[\pi,\pi]=0$ を満たす
$\pi\in H^0(X,\bigwedge^2T_X)$ を Poisson 構造と呼ぶ。
射影空間の場 $\pi$、その制限 $\pi_X$、アフィン錐上の代表 $q$ を区別する。
射影スキームに関する斉次イデアルの等式は、不要イデアルで飽和したイデアル層の等式を意味する。
滑らかさから $F$ は既約である。実際、正次数の二因子があればその射影零点は交わり、交点は $X$ の特異点となる。

\subsection{Gradient の Koszul 複体}
\begin{lemma}\label{lem:pg-gradient}
$F_1,\ldots,F_N$ は $S$ の正則列である。
特に $\sum_i a_iF_i=0$ で、各 $a_i$ が定数または一次式ならば全ての $a_i$ は零である。
\end{lemma}
\begin{proof}
Euler 等式 $3F=\sum_i z_iF_i$ と射影的な滑らかさにより、gradient の共通零点はアフィン原点だけである。
従ってそのイデアルの高さは $N$ である。
$S$ は Cohen--Macaulay なので、$N$ 個の生成元は正則列となる。
第一 syzygy は $F_je_i-F_ie_j$ という Koszul 関係で生成され、その係数次数は二である。
よって係数次数零または一の syzygy は存在しない。
\end{proof}
\begin{lemma}\label{lem:pg-euler}
$\mathbb P=\PP(V)$ とすると
\[
H^0(\mathbb P,\textstyle\bigwedge^2 T_{\mathbb P}(-3))
=H^1(\mathbb P,\textstyle\bigwedge^2 T_{\mathbb P}(-3))=0,
\qquad
H^0(\mathbb P,\textstyle\bigwedge^3 T_{\mathbb P}(-3))\simeq\bigwedge^3V.
\]
\end{lemma}
\begin{proof}
$\dim\mathbb P\geq4$ における線束の中間コホモロジーの消滅\cite[Chapter III]{PHartshorne}を使う。
Euler 列とその外積は
\begin{align*}
0&\longrightarrow\OO(-3)\longrightarrow V\otimes\OO(-2)
 \longrightarrow T_{\mathbb P}(-3)\longrightarrow0,\\
0&\longrightarrow T_{\mathbb P}(-3)\longrightarrow\bigwedge^2V\otimes\OO(-1)
 \longrightarrow\bigwedge^2T_{\mathbb P}(-3)\longrightarrow0,\\
0&\longrightarrow\bigwedge^2T_{\mathbb P}(-3)\longrightarrow\bigwedge^3V\otimes\OO
 \longrightarrow\bigwedge^3T_{\mathbb P}(-3)\longrightarrow0.
\end{align*}
初めの二列が消滅を与え、最後の列の大域切断が同型を与える。
\end{proof}
\begin{proposition}\label{prop:pg-represent}
$F$ を固定すると線形同型
\[
\bigwedge^3V\xrightarrow{\sim} H^0(X,\bigwedge^2T_X),
\qquad P\longmapsto(\iota_{dF}P)|_X
\]
がある。右辺は二次係数の場を射影化して定める。
任意の双ベクトル場は $\PP(V)$ 上に一意に延長される。
\end{proposition}
\begin{proof}
補題\ref{lem:pg-euler}と制限列から
$H^0(\mathbb P,\bigwedge^2T_{\mathbb P})\simeq
H^0(X,(\bigwedge^2T_{\mathbb P})|_X)$ である。
従って延長は存在して一意であり、外積 Euler 列により
$B\in S_2\otimes\bigwedge^2V$ で代表できる。
接性により、$\iota_{dF}B$ の射影接束への像は $X$ 上零である。
Euler 列の核と $H^0(X,\OO_X(3))$ の多項式への持上げを用いると
\[
\iota_{dF}B=FA+hE,\qquad A\in S_1\otimes V,\quad h\in S_3.
\]
再び $dF$ で縮約すれば $0=FA(F)+3Fh$、従って $h=-A(F)/3$ である。
そこで $q=B-\frac13E\wedge A$ とおくと $\iota_{dF}q=0$ となる。
補題\ref{lem:pg-gradient}の Koszul 複体の次数部分
\[
S_0\otimes\bigwedge^3V\xrightarrow{\iota_{dF}}
S_2\otimes\bigwedge^2V\xrightarrow{\iota_{dF}}S_4\otimes V
\]
は完全で、左からの寄与は負次数なので零である。
よって $q=\iota_{dF}P$ を満たす定数 $P$ が一意に存在する。

代表の選択にも依存しない。
射影化が零である二次係数の場は $E\wedge Z$（$Z$ は線形）と書ける。
これが $F$ を Casimir にもつなら $3FZ=Z(F)E$ である。
既約性と各座標が $F$ と互いに素であることから $F\mid Z(F)$、
従って $Z(F)=cF$、$Z=(c/3)E$、$E\wedge Z=0$ となる。
逆に任意の $P$ について $\iota_{dF}^2P=0$ なので、射影化は $X$ に接する。
\end{proof}

\subsection{Jacobi 恒等式の延長}
\begin{proposition}\label{prop:pg-jacobi-extension}
$q=\iota_{dF}P$ の $X$ 上の射影化が Poisson ならば $[q,q]=0$ である。
従って射影空間上の一意な延長も Poisson である。
\end{proposition}
\begin{proof}
標準体積形式の divergence を
\[
\delta(a\partial_{i_1}\wedge\cdots\wedge\partial_{i_r})
=\sum_{b=1}^r(-1)^{b-1}\partial_{i_b}a\,
\partial_{i_1}\wedge\cdots\wedge\widehat{\partial_{i_b}}\wedge\cdots\wedge\partial_{i_r}
\]
と定める。交代テンソルと Hessian の縮約なので $\delta q=0$。
積の微分則から $\delta$ は Schouten 括弧の導分であり、
$J=[q,q]$ は $\delta J=0$ を満たす。
$F$ が Casimir であることから $\iota_{dF}J=0$ でもある。
$J$ の射影化は $X$ 上零なので、補題\ref{lem:pg-euler}より
\[
J=FT+E\wedge B_2,\qquad T\in\bigwedge^3V,
\quad B_2\in S_2\otimes\bigwedge^2V.
\]
定義から $\delta(FT)=\iota_{dF}T$、
$\delta(E\wedge B_2)=NB_2-E\wedge\delta B_2$。
$\delta J=0$ に $E$ を外積すると
$E\wedge B_2=-N^{-1}E\wedge\iota_{dF}T$ となる。
これを代入して $dF$ で縮約すると
\[
0=(1-3/N)F\,\iota_{dF}T.
\]
$N\geq5$ より $\iota_{dF}T=0$。
gradient の定数係数での独立性から $T=0$、従って $J=0$ である。
\end{proof}

\section{Hessian による Poisson 構造の正規形}\label{sec:pg-normal}
\subsection{Jacobi 条件の線形化}
$P\in\bigwedge^3V$、$\alpha,\beta\in V^*$ に対し、
$\gamma(P(\alpha,\beta))=P(\alpha,\beta,\gamma)$ により
$P(\alpha,\beta)\in V$ を定める。次の線形空間を導入する：
\begin{align}
K(P)=\{H\in\operatorname{Sym}^2V^*\mid{}
&H(P(\alpha,\beta),P(\gamma,\eta))
 +H(P(\beta,\gamma),P(\alpha,\eta))\notag\\
&+H(P(\gamma,\alpha),P(\beta,\eta))=0
\quad(\alpha,\beta,\gamma,\eta\in V^*)\}.\label{eq:pg-K}
\end{align}
\begin{lemma}\label{lem:pg-linearize}
$q=\iota_{dF}P$ が Poisson であることと、全 $p\in V$ について
$\operatorname{Hess}_F(p)\in K(P)$ であることは同値である。
また $F$ の滑らかさの仮定の下で、Hessian の線形系は $\PP(V)$ 上基点をもたない。
\end{lemma}
\begin{proof}
$P^{ija}=P(dz_i,dz_j,dz_a)$ とすれば $q^{ij}=\sum_aP^{ija}F_a$ である。
Jacobiator の係数は非零定数倍を除き
\[
\sum_aF_aL_a^{ijk},\qquad
L_a^{ijk}=\sum_{r,b}(P^{ira}P^{jkb}+P^{jra}P^{kib}+P^{kra}P^{ijb})F_{br}.
\]
$L_a^{ijk}$ は一次式なので、補題\ref{lem:pg-gradient}より
Jacobi 条件は全 $L_a^{ijk}=0$ と同値である。
これは\eqref{eq:pg-K}への座標共変ベクトルの代入と符号を除いて一致する。
多重線形性が第一の主張を与える。
三次微分の対称性により
$\operatorname{Hess}_F(p)(w,w)=2dF_w(p)$。
非零の $w$ が全 Hessian の零点なら $dF_w=0$、Euler 等式より $F(w)=0$ となり、
$[w]$ が $X$ の特異点となるので矛盾する。
\end{proof}

\subsection{基点をもたない二次形式と三重ベクトル}
\begin{lemma}\label{lem:pg-quadratic-lie}
$K(P)$ が $\PP(V)$ 上基点をもたなければ、直和をなす三次元部分空間
$W_1,\ldots,W_k\subset V$ と $0\neq P_s\in\bigwedge^3W_s$ が存在して
$P=\sum_sP_s$ と書ける。$P=0$ では $k=0$ とする。
\end{lemma}
\begin{proof}
$P\neq0$ とし、その support を
$W=\operatorname{span}\{P(\alpha,\beta)\mid\alpha,\beta\in V^*\}$ とおく。
これは $P\in\bigwedge^3W$ を満たす最小の部分空間である。
$K(P)$ の $\PP(W)$ への制限も基点をもたない。
標数零の Bertini 定理により、それぞれの一般の二次超曲面は滑らかであり、
従って $V$ 上および $W$ 上での非退化条件は非空開条件である。
両者を満たす $B\in K(P)$ を選ぶ。

$B^\flat:V\to V^*$ により $[x,y]=P(B^\flat x,B^\flat y)$ と定める。
$B([x,y],z)=P(B^\flat x,B^\flat y,B^\flat z)$ は交代的なので $B$ は不変である。
Jacobiator を $B$ で任意の $t$ と対にすると、\eqref{eq:pg-K}の $H=B$ が得られる。
よってこれは Lie 代数 $\mathfrak g$ である。
$[\mathfrak g,\mathfrak g]=W$、$Z(\mathfrak g)=W^\perp$ であり、$B|_W$ の非退化性から
\[
\mathfrak g=W\oplus Z(\mathfrak g),\qquad [W,W]=W,\qquad Z(W)=0.
\]
さらに任意の $H\in K(P)$ について
\begin{equation}
H([a,b],[c,d])+H([b,c],[a,d])+H([c,a],[b,d])=0.\label{eq:pg-lie-H}
\end{equation}

$\mathfrak r=\operatorname{rad}(W)\neq0$ と仮定する。
その導来列の最後の非零項から極小の非零 $W$-イデアル $I$ を取る。
$I$ は可換で、既約な $W$-加群である。
複素有限次元 Lie 代数の既約表現では可解根基はスカラーで作用する
\cite[Proposition 16.9]{PEtingofLie}。
$W=[W,W]$ のため全作用素の trace は零であり、$\mathfrak r$ は $I$ に零で作用する。
従って $I$ は $W/\mathfrak r$ の既約加群であり、$Z(W)=0$ から自明表現ではない。
最高ウェイトベクトル $0\neq v\in I$ と単純根 $\alpha$ を
$\lambda(h_\alpha)>0$ となるように選び、
$e_\alpha,f_\alpha,h_\alpha$ を $W$ に持ち上げる。
$w=\lambda(h_\alpha)^{-1}[f_\alpha,v]\in I$ とおくと
\[
[e_\alpha,w]=v,\quad [e_\alpha,v]=0,\quad
[h_\alpha,v]=\lambda(h_\alpha)v,\quad[v,w]=0.
\]
商の $\mathfrak{sl}_2$ 関係が作用するのは、根基の作用が零だからである。
\eqref{eq:pg-lie-H}に $(a,b,c,d)=(h_\alpha,v,e_\alpha,w)$ を代入すると
$\lambda(h_\alpha)H(v,v)=0$ が全 $H$ に対して成立する。
これは基点自由性に反する。従って $W$ は半単純である。

階数二以上の単純因子があるとする。
その最高根 $\theta$ は単純根ではない
（既約かつ階数二以上の根系では最高根のsupportは全ての単純根を含む）。
正根でありながら単純でない根は二つの正根の和となるので
\cite[Lemma~21.14および\S32.2]{PEtingofLie}、
$\theta=\alpha+\beta$ と書ける。
対応する根ベクトル$a=e_\alpha$、$b=e_\beta$を選ぶと
根空間の標準的交換子関係から
$v=[a,b]\neq0$ は最高根空間の元となる。
最高根性により$\theta+\alpha,\theta+\beta$は根でなく、
従って $[a,v]=[b,v]=0$ である。
$\theta(h)\neq0$ となる Cartan 元を選び、
\eqref{eq:pg-lie-H}に $(a,b,h,v)$ を代入すると
$\theta(h)H(v,v)=0$ となり、再び基点を得て矛盾する。
従って $W\simeq\mathfrak{sl}_2^{\oplus k}$ である。
用いた可解根基の表現定理は\cite[Proposition~16.9]{PEtingofLie}、
根分解・最高ウェイトと最高根の事実は
\cite[\S19.4, Lemma~21.14, Theorem~25.17, \S32.2]{PEtingofLie}を参照。

異なる単純イデアルは $B$ に関して直交する。
実際、$x=[a,b]$ が一方に、$y$ が他方に属せば
$B(x,y)=B(a,[b,y])=0$ である。中心も $W$ に直交する。
よって Cartan 三形式 $B([x,y],z)$ は各 $\mathfrak{sl}_2$ 因子上の三形式の和であり、
$B$ で添字を上げれば所要の $P=\sum_sP_s$ を得る。
\end{proof}

\subsection{三次式の分解と一般階数}
\begin{theorem}[Poisson 正規形]\label{thm:pg-normal}\label{pthm:normal-form}
任意の非零 Poisson 構造に対して、$N=3k+m$ となる $k\geq1$ と線形座標
$x_s=(x_{s1},x_{s2},x_{s3})$、$u=(u_1,\ldots,u_m)$ が存在し、
\begin{equation}
P=\sum_{s=1}^k\partial_{x_{s1}}\wedge\partial_{x_{s2}}\wedge\partial_{x_{s3}},
\qquad F=\sum_{s=1}^k f_s(x_s,u)\label{eq:pg-normal}
\end{equation}
となる。括弧は
\begin{equation}
\{x_{s1},x_{s2}\}=(f_s)_{x_{s3}},\quad
\{x_{s2},x_{s3}\}=(f_s)_{x_{s1}},\quad
\{x_{s3},x_{s1}\}=(f_s)_{x_{s2}}\label{eq:pg-bracket}
\end{equation}
と、異なるブロック間および $\{u_\alpha,-\}$ が零であることで定まる。
逆に\eqref{eq:pg-normal}を満たす滑らかな $F$ はこの括弧によって Poisson 構造を定める。
$\pi$ と $\pi_X$ の一般階数はともに $2k$ である。
\end{theorem}
\begin{proof}
補題\ref{lem:pg-linearize}、\ref{lem:pg-quadratic-lie}を適用し、各三次元因子の基底の行列式を調整する。
異なる因子からそれぞれ $a,b$ と $c,d$ を取ると、
\eqref{eq:pg-lie-H}は $H([a,b],[c,d])=0$ となる。
$[W_s,W_s]=W_s$ より、全 Hessian の $W_s\times W_t$ 成分は $s\neq t$ のとき零である。
従って $\partial_{x_{si}}\partial_{x_{tj}}F=0$。
$F$ の各単項式は高々一つの active ブロックを含み、中心変数のみの項をいずれかに配分すれば所要の分解となる。
逆向きは三変数 Jacobian 括弧の Jacobi 恒等式
$\nabla f\cdot(\nabla\times\nabla f)=0$ と、異なるブロックが互いの係数を微分しないことから従う。

各 active gradient は恒等的に零ではない。
そうなら $F$ は対応する三変数に依存せず、その座標平面に特異点をもつ。
従ってアフィン空間の一般階数は $2k$。
非零な二次偏微分は既約な三次式 $F$ で割り切れないので、$X$ の一般点でも全ブロックが階数二をもつ。
階数 $2r$ の交代双ベクトルを直線 $L$ で割ると、
$L$ がその像に入らなければ階数は $2r$、入れば $2r-2$ である。
これは像上の Darboux 基底により分かる。
ここで $L=\CC E$。
$m>0$ なら一般点で $E$ は中心方向成分をもち、像に入らない。
$m=0$ なら第 $s$ ブロックの像は $\ker df_s$ で、Euler 方向が入る条件は $3f_s=0$。
$N\geq5$ の下で $m=0$ なら $k\geq2$ であり、全 $f_s=0$ は $X$ の真の閉集合である。
よって一般点では射影化で階数は低下しない。
\end{proof}
座標、ブロックの順序、中心三次式の配分の一意性は主張しない。
また $F$ を $aF$ に替えると命題\ref{prop:pg-represent}の同型は $a$ 倍される。

\begin{corollary}\label{cor:pg-generic}
固定した滑らかな $X$ 上の一般階数二の Poisson 構造の射影化は
$\operatorname{Gr}(3,V)\subset\PP(\bigwedge^3V)$ である。
一般の滑らかな三次超曲面上では、全ての非零 Poisson 構造が一般階数二をもつ。
\end{corollary}
\begin{proof}
$k=1$ と非零 $P$ の分解可能性が同値なので第一の主張を得る。
第二の主張には $k\geq2$ の族の次元を評価する。
直和をなす順序付き $W_1,\ldots,W_k$ の空間は
$\operatorname{Gr}(3,V)^k$ の開部分で、次元 $3k(N-3)$。
この組を固定したときの三次式の空間の次元は
\[
\binom{m+2}{3}+k\left(10+6m+3\binom{m+1}{2}\right).
\]
従って全三次式の空間における余次元は少なくとも
\begin{equation}
\binom{N+2}{3}-\binom{m+2}{3}
-k\left(10+6m+3\binom{m+1}{2}\right)-3k(N-3)
=\frac{3k}{2}\bigl((3k-5)m+3k(k-1)\bigr)>0.\label{eq:pg-codim}
\end{equation}
可能な $k$ は有限個なので一般の三次式はこれらを避ける。
\end{proof}
\begin{remark}
\eqref{eq:pg-codim}は余次元の下界である。
$(n,k)=(6,2),(7,2),(7,3)$ ではそれぞれ $24,27,81$ となる。
正確な余次元や正規形の一意性を結論するものではない。
三次元の場合の Grassmann 記述は\cite[Theorem 1]{PMayanskiy}の既知の結果である。
\end{remark}

\section{Casimir、対数形式および symplectic leaf}\label{sec:pg-leaves}
以下、\eqref{eq:pg-normal}を固定する。
多項式 $G$ が Casimir であるとは $\{G,S\}=0$ を意味する。
射影超曲面が Poisson であるとは、そのイデアル層が括弧で保たれることをいう。

\subsection{Poisson 超曲面}
\begin{proposition}\label{prop:pg-casimir}
既約な斉次式 $G\in S$ に対し、$V(G)\subset\PP(V)$ が Poisson であることと、
$G$ がアフィン錐上の Casimir であることは同値である。
一次 Casimir の空間は $\operatorname{span}(u_1,\ldots,u_m)$ である。
従って Poisson 超平面のパラメータ空間は $m>0$ なら $\PP^{m-1}$、$m=0$ なら空である。
\end{proposition}
\begin{proof}
Casimir なら $(G)$ は括弧で保たれる。
逆向きは $G$ が $F$ の定数倍の場合は既知なので、$\gcd(F,G)=1$ とする。
$d=\deg G$ とおくと、射影的接性と Euler 列より
\[
\iota_{dG}q=GA+hE,\qquad A\in S_1\otimes V,\quad h\in S_d.
\]
$dG$ で縮約すれば $h=-A(G)/d$。
$dF$ で縮約し $\iota_{dF}q=0$ を使うと
$GA(F)-(3/d)FA(G)=0$ である。
互いに素であることと次数から $A(F)=cF$、$A(G)=cdG/3$。
$A-(c/3)E$ は gradient の一次 syzygy なので零である。
従って $\iota_{dG}q=0$。
一次式 $\ell$ については gradient の定数係数での独立性から
$\iota_{d\ell}q=0\iff\iota_{d\ell}P=0$ であり、正規形では
$\ell$ が全 active 変数に依存しないことと同値である。
\end{proof}
各 $u_\alpha,f_s,F$ とその多項式結合は Casimir であるが、
多項式中心全体がそれらで生成されるとは主張しない。

\subsection{最高外積の対数表示}
\begin{proposition}\label{prop:pg-log}
体積形式を固定すると、$\pi^k$ に対応する微分形式は非零定数倍を除き
\begin{equation}
\beta=\iota_E(du_1\wedge\cdots\wedge du_m\wedge df_1\wedge\cdots\wedge df_k)
\in H^0(\PP^{N-1},\Omega^{m+k-1}(N)).\label{eq:pg-log-main}
\end{equation}
$m>0$ では
\[
\beta=\left(\prod_{\alpha=1}^m u_\alpha\prod_{s=1}^kf_s\right)
\bigwedge_{\alpha=2}^m d\log(u_\alpha/u_1)
\wedge\bigwedge_{s=1}^k d\log(f_s/u_1^3),
\]
$m=0$ では
\[
\beta=3\left(\prod_{s=1}^kf_s\right)
\bigwedge_{s=2}^k d\log(f_s/f_1).
\]
$X$ 上の $\pi_X^k$ に対応する形式は
\begin{equation}
\alpha=\iota_E(du_1\wedge\cdots\wedge du_m\wedge df_1\wedge\cdots\wedge df_{k-1})|_X
\in H^0(X,\Omega_X^{m+k-2}(N-3)).\label{eq:pg-log-X}
\end{equation}
\end{proposition}
\begin{proof}
第 $s$ ブロックの場を $q_s$ と書く。
$q_s^2=0$ より $q^k=k!q_1\wedge\cdots\wedge q_k$。
アフィン体積形式への縮約は、非零定数倍を除き
$du_1\wedge\cdots\wedge du_m\wedge df_1\wedge\cdots\wedge df_k$ である。
$df_s$ の中心方向成分は先頭の外積で消える。
さらに $E$ で縮約すれば射影化に対応する水平形式となり、
全スケーリングの重み $m+3k=N$ が twist を与える。
$m>0$ では $\iota_Ed\log u_1=1$、各比の対数微分の縮約は零であり、外積を展開して所要の式を得る。
$m=0$ では $\iota_Ed\log f_1=3$ を使う。
計算は分母が消えない開集合で行うが、左辺が大域的正則なので大域表示として解釈できる。

最後に $f_k$ を $F=\sum_sf_s$ に替えても微分の外積は変わらない。
$X$ 上で法線方向 $dF$ を除き adjunction を適用すると\eqref{eq:pg-log-X}を得る。
局所的に $dF$ が座標微分となる表示で縮約を行えばよく、法線束 $\OO_X(3)$ の遷移関数が twist を三減らす。
形式次数は $m+k-2=n-2k$ となる。
\end{proof}
対数表示の因子の次数の総和は $m+3k=N$ であり、
$V(u_1\cdots u_mf_1\cdots f_{k-1}F)$ を反標準因子にできる。
これは最高外積の表示であって、$k>1$ では $\pi^k$ だけから $\pi$ は復元できず、
この因子の全体が退化軌跡であるとも限らない。

\subsection{葉とその閉包}
\begin{proposition}\label{prop:pg-first-integrals}
有理写像
\[
\Phi:\PP^{N-1}\dashrightarrow\PP(1^m,3^k),\qquad
[x,u]\longmapsto[u_1:\cdots:u_m:f_1:\cdots:f_k]
\]
は一般点で独立な第一積分を与える。
一般 fiber の滑らかな最大階数部分の連結成分は symplectic leaf である。
$X$ 上の底は $m>0$ なら $\PP(1^m,3^{k-1})$、$m=0$ なら $\PP^{k-2}$ であり、次元は $n-2k$ である。
\end{proposition}
\begin{proof}
全 active gradient が非零の開集合では、$du_\alpha$ を法とした $df_s$ は独立なブロックに非零成分をもつ。
従って $(u,f_1,\ldots,f_k)$ の微分の階数は $m+k$。
$F=0$ への制限と射影化は、それぞれ独立な条件 $\sum_sf_s=0$ とスケーリングを一つ除く。
Casimir 性から Hamilton 分布は $\ker d\Phi$ に含まれ、
定理\ref{thm:pg-normal}によって両者の次元は $2k$ なので一致する。
最大階数部分での逆二形式は Jacobi 恒等式により閉じている。
よって滑らかな fiber の連結成分が葉となる。
不定点では $u=0$ かつ全 $f_s=0$。
active gradient が一つでも零ならアフィン階数は高々 $2k-2$、
全て非零なら Euler 方向が像に入るため射影階数は $2k-2$ である。
よって最大階数の葉はこの不定点集合を横切らない。
\end{proof}
\begin{proposition}\label{prop:pg-leaf-closure}
各 $f_{s,0}=f_s(x_s,0)$ が滑らかな平面三次曲線を定めるとする。
$m>0$ なら一般の葉の閉包は $\PP^{3k}$ 内の型 $(3^k)$ の既約な完全交叉であり、
$m=0$ なら $\PP^{3k-1}$ 内の型 $(3^{k-1})$ の既約な完全交叉である。
いずれも次元は $2k$ である。
\end{proposition}
\begin{proof}
$m>0$ では底の一般点で $u_\alpha=c_\alpha u_1$、$c_1=1$ とし、$t=u_1$ とおく。
候補は $f_s(x_s,ct)-d_st^3=0$（$1\leq s\leq k$）であり、$X$ 上では $\sum_sd_s=0$。
$t=0$ の部分は独立な式 $f_{s,0}=0$ で定まり次元 $2k-1$。
$t\neq0$ では $t=1$ と正規化すると $k$ 個のアフィン曲面の積となる。
従って候補は純次元 $2k$ の完全交叉で、無限遠に余分な成分をもたない。
標数零の generic smoothness により、一般の底点上のアフィン積は滑らかである。
各因子の射影閉包は無限遠でも $\nabla f_{s,0}\neq0$ により滑らかであり、
連結な射影超曲面なので既約である。
よってアフィン fiber は既約。
完全交叉は Cohen--Macaulay で埋め込み成分をもたず、全最小成分がこの滑らかな開集合と交わる。
従って候補は被約かつ既約で、葉の閉包である。

$m=0$ では $f_1\neq0$ の一般底点に対して $f_s=c_sf_1$（$s\geq2$）と書ける。
$f_1=0$ の部分は全 $f_s=0$ で定まり次元 $2k-1$。
開部分は $\{f_1=1\}\times\prod_{s=2}^k\{f_s=c_s\}$ の同時 $\mu_3$ 作用による像である。
$c_s\neq0$ の一般点では、各因子は滑らかな三次曲面
$f_s(x_s)=c_st^3$ のアフィン開集合なので既約である。
全 active gradient が非零なので方程式の微分は独立であり、fiber は滑らか。
前段と同じ純次元性・Cohen--Macaulay 性により結論を得る。
\end{proof}
\begin{remark}\label{rem:pg-zero-restriction-leaf}
平面三次制限の滑らかさを仮定しなくても命題\ref{prop:pg-first-integrals}は成立する。
ただし上の方程式から閉包を得るには、$m>0$ では $t$ に関して飽和し、必要なら fiber の既約成分を選ぶ。
例えば $k=1$、$f_{1,0}=0$ では $F(x,ct)=tQ_c(x,t)$ となる。
$Q_\alpha(x)=F_{u_\alpha}(x,0)$ は $\PP_x^2$ 上共通零点をもたない。
そうした点は $X$ の特異点となるからである。
Bertini により一般の $\sum_\alpha c_\alpha Q_\alpha$ は滑らかな平面二次曲線を定める。
従って $Q_c(x,t)$ の階数は少なくとも三で、$t$ で割り切れず既約である。
葉の閉包は固定平面 $t=0$ を除いた residual quadric $Q_c=0$ となる。
\end{remark}

\section{全退化スキームの方程式}\label{sec:pg-degeneracy}
\subsection{任意のブロック表示に対する生成元}
$q_s=\iota_{df_s}(\partial_{x_{s1}}\wedge\partial_{x_{s2}}\wedge\partial_{x_{s3}})$ とし、
\[
A_{s,i}=\frac{\partial f_s}{\partial x_{si}},\qquad
p_s=\sum_i x_{si}A_{s,i}=3f_s-\sum_\alpha u_\alpha\frac{\partial f_s}{\partial u_\alpha}
\]
とおく。$D_{2j-2}(\pi)=Z(\pi^j)$ をスキームとして考える。
\begin{proposition}\label{prop:pg-deg-ideal}
$1\leq j\leq k$ に対し、$D_{2j-2}(\pi)$ のイデアル層は、全ての
$I\subset\{1,\ldots,k\}$、$|I|=j$ について
\begin{equation}
 u_\alpha\prod_{s\in I}A_{s,i_s},\qquad
 p_t\prod_{s\in I\setminus\{t\}}A_{s,i_s}\ (t\in I),\qquad
 x_{t,l}\prod_{s\in I}A_{s,i_s}\ (t\notin I)\label{eq:pg-deg-generators}
\end{equation}
を生成元とする。各添字は可能な範囲を動く。
$X$ 上のイデアルはさらに $F$ を加えて得られる。
各生成元の次数は $2j+1$ である。
\end{proposition}
\begin{proof}
錐の原点以外では $E$ の生成する直線を局所的な直和因子に選べる。
その商束の外積から $E\wedge-$ は分裂単射となるので、
射影化の零点イデアルは $E\wedge q^j$ の係数イデアルに等しい。
$q_s^2=0$ より $q^j=j!\sum_{|I|=j}\bigwedge_{s\in I}q_s$。
$E$ の中心成分は第一の型を与える。
$t\in I$ では $(\sum_lx_{tl}\partial_{x_{tl}})\wedge q_t
=p_t\partial_{x_{t1}}\wedge\partial_{x_{t2}}\wedge\partial_{x_{t3}}$ なので第二の型となる。
$t\notin I$ は第三の型を与える。
ブロック別外積次数は、中心に一、選択したブロックに三、選択しないブロックに一、という異なる型をもつ。
残りの選択ブロックは二であるため、異なる型の間に相殺はない。
最後に $T_X\subset T_{\PP(V)}|_X$ は局所的に分裂するので、外積への包含は切断の零点イデアルを変えない。
\end{proof}

\subsection{有限平坦 critical cover}
本小節以降、特に断らない限り
\begin{equation}
f_{s,0}=f_s(x_s,0)\text{ が滑らかな平面三次曲線 }E_s\text{ を定める}\label{eq:pg-smooth-block}
\end{equation}
と仮定する。$E_s$ は $u=0$、$x_t=0$（$t\neq s$）により $X$ 内に入り、互いに交わらない。
\begin{lemma}\label{lem:pg-critical-cover}
$m>0$ なら
\[
K_s=V(A_{s,1},A_{s,2},A_{s,3})\subset\PP(x_s,u)=\PP^{m+2},
\qquad\Delta=V(A_{s,i}\mid s,i)
\]
は $u=0$ と交わらない。
中心座標への射影は $K_s\to\PP_u^{m-1}$ を有限平坦八次被覆、
$\Delta\to\PP_u^{m-1}$ を有限平坦 $8^k$ 次被覆とし、
$\Delta\simeq K_1\times_{\PP_u^{m-1}}\cdots\times_{\PP_u^{m-1}}K_k$。
$m\geq0$ の場合を含め、$\Gamma_I=V(A_{s,i}\mid s\in I,1\leq i\leq3)$ は
非空なら型 $(2^{3|I|})$ の完全交叉で、
\[
\dim\Gamma_I=N-1-3|I|,\qquad\deg\Gamma_I=8^{|I|}.
\]
特に $m>0$ では
$\dim\Delta=m-1$、$\deg\Delta=8^k$、$\omega_\Delta\simeq\OO_\Delta(3k-m)$。
$m=0$ では $\Delta=\varnothing$ とする。
\end{lemma}
\begin{proof}
$R=\CC[u_1,\ldots,u_m]$、$Q_s=R[x_s]/(A_{s,1},A_{s,2},A_{s,3})$ とする。
$Q_s/(u)$ は滑らかな三次式の Jacobian 環で、Hilbert 級数 $(1+t)^3$ をもつ。
その八個の斉次基底を持ち上げれば、次数による帰納法から $Q_s$ は有限 $R$-加群となる。
三方程式による次元下界と有限性より $\dim Q_s=m$。
従って方程式は正則列で $Q_s$ は Cohen--Macaulay である。
$u_1,\ldots,u_m$ は斉次パラメータ系なので $Q_s$-正則列である。
持上げからの全射
\[
R\oplus R(-1)^{\oplus3}\oplus R(-2)^{\oplus3}\oplus R(-3)\longrightarrow Q_s
\]
は同型である。実際、$u$ の Koszul 完全性により
$\operatorname{Tor}_1^R(Q_s,R/(u))=0$ であり、全射の核を $(u)$ で割ると零となる。
次数付き Nakayama が核を零にする。
複数ブロックの環は $R$ 上のテンソル積なので階数 $8^{|I|}$ の有限自由加群である。
残りの active 変数を加えれば、$\Gamma_I$ の高さと完全交叉性が従う。
$m=0$ かつ $I$ が全体のときだけ射影スキームは空である。

$u=0$ では各 gradient の共通零点は原点だけなので、$K_s$ と $\Delta$ は射影的にその部分と交わらない。
$u_\alpha\neq0$ で全ブロックを同時に $u_\alpha=1$ と正規化すれば、有限自由性が被覆の有限平坦性と fiber product 表示を与える。
重なりでも同時正規化するため貼り合わさる。次数と双対化層は Hilbert 級数および随伴公式から従う。
\end{proof}

\subsection{退化集合と零点スキーム}
空でない $L\subset\{1,\ldots,k\}$ に対し
\[
J_L=V(u,x_t\ (t\notin L),f_{s,0}\ (s\in L))
\]
を $E_s$（$s\in L$）の射影的 join とする。
独立な錐の積の射影化なので、被約既約で次元 $2|L|-1$ である。
\begin{proposition}\label{prop:pg-deg-support}
$0\leq r<k$ について
\[
|D_{2r}(\pi)|=\bigcup_{|I|=k-r}|\Gamma_I|\ \cup\ \bigcup_{|L|=r+1}J_L.
\]
$X$ 上では右辺を $X$ と交わらせる。
零点についてはスキームとして
\[
D_0(\pi)=\left(\coprod_sE_s\right)\sqcup\Delta,
\qquad D_0(\pi_X)=\left(\coprod_sE_s\right)\sqcup Z,
\qquad Z=\Delta\cap X.
\]
\end{proposition}
\begin{proof}
非零 active gradient をもつブロック数を $\ell$ とするとアフィン階数は $2\ell$。
$u\neq0$ では Euler 方向が像に入らず、射影化で階数は変わらない。
$u=0$ では\eqref{eq:pg-smooth-block}により gradient が零であることと $x_s=0$ が同値である。
非零ブロックの像は $\ker df_{s,0}$ なので、Euler 方向が像に入る条件は各 $3f_{s,0}=0$。
この場合に限り射影階数は $2\ell-2$ となる。これが台の表示を与える。

スキーム構造には命題\ref{prop:pg-deg-ideal}を用いる。
$u_\alpha\neq0$ では $u_\alpha A_{s,i}$ が全 active 偏微分を生成するので零点は $\Delta$ である。
$E_s$ の点ではある $A_{s,i}$ が単元となり、生成元は局所的に $u,x_t$（$t\neq s$）および $p_s$ を生成する。
これらを法として $p_s=3f_{s,0}$ で、他の生成元は消える。
従って零点イデアルは厳密に $(u,x_t\ (t\neq s),f_{s,0})$。
各 $E_s$ は被約であり、$\Delta$ と交わらないので所要の分解となる。
\end{proof}
\begin{remark}
台の表示を高次の退化スキームの被約分解と取り違えてはならない。
例えば $m=0$ の $\Gamma_I$ は座標線形空間を台とするが、Jacobian イデアルは一般に非被約な厚みを定める。
\end{remark}

\subsection{完全交叉性と一般の滑らかさ}
\begin{proposition}\label{prop:pg-zero-ci}
$m=1$ なら $Z=\varnothing$。
$m\geq2$ なら一般性を仮定せず $Z$ は型 $(2^{3k},3)$ の完全交叉で、
\[
\dim Z=m-2,\qquad \deg Z=3\cdot8^k,\qquad
\omega_Z\simeq\OO_Z(3k+3-m).
\]
滑らかな $f_{s,0}$ を固定して他の係数を一般に選ぶと、$K_s,\Delta$ および $m\geq2$ の場合の $Z$ は滑らかである。
これらが正次元ならば、滑らかな場合には既約である。
\end{proposition}
\begin{proof}
$\Delta$ の斉次環を $Q$ とする。
$Q$ は $R=\CC[u]$ 上有限自由な Cohen--Macaulay 環で、各最小素イデアルは $R$ の零イデアルに縮約する。
実際、各成分は純次元性により次元 $m$ であり、有限性から像も次元 $m$ である。
ある成分上 $F=0$ とし、その関数体を $L$ とする。
$L/\CC(u)$ は有限分離拡大なので、$du_1,\ldots,du_m$ は $\Omega^1_{L/\CC}$ の基底である。
全 active 偏微分が零なので
$0=dF=\sum_\alpha F_{u_\alpha}du_\alpha$、従って全中心偏微分も零となる。
この成分は $u$-空間を支配し原点以外の点をもつため、$X$ の滑らかさに矛盾する。
よって $F$ はどの最小素点にも入らず、Cohen--Macaulay 性により非零因子である。
これが $m\geq2$ の完全交叉性と数値を与える。
$m=1$ では交点があれば $u_1\neq0$ と $3F=u_1F_{u_1}$ により全偏微分が零となり矛盾する。

一般性を確認する。$u$ を含む係数を動かすパラメータ空間では、
$u_\alpha\neq0$ の点で項 $x_{si}u_\alpha^2$ の係数が $A_{s,i}$ の値を独立に変える。
従って全 critical incidence は滑らかである。
$u=0$ には解がなく、incidence はパラメータ空間と射影空間の積の閉部分なので、その射影は proper である。
全 fiber は非空な有限被覆なので射影は支配的であり、標数零の generic smoothness から一般の $K_s,\Delta$ は滑らかとなる。
そのような $\Delta$ を固定し $F$ に中心三次式 $c(u)$ を加えても active 偏微分は変わらない。
切断 $c(u)|_\Delta$ は基点をもたず、Bertini から一般の $Z$ は滑らかとなる。
$X$ の滑らかさも非空開条件である。
実際 $\sum_sf_{s,0}+\sum_\alpha u_\alpha^3$ が滑らかな例を与える。
従ってこれらの開条件を同時に課せる。
正次元の射影完全交叉の Koszul resolution と線束の中間コホモロジー消滅から $H^0(\OO)=\CC$。
よって連結であり、滑らかならば既約である。
\end{proof}

\section{六次元および七次元の幾何}\label{sec:pg-six-seven}
\subsection{可能な階数}
定理\ref{thm:pg-normal}により次の表を得る。
\begin{center}
\begin{tabular}{ccccc}\toprule
$n$ & $k$ & 一般階数 & $m$ & Poisson 超平面\\\midrule
6&1&2&5&$\PP^4$\\
6&2&4&2&$\PP^1$\\
7&1&2&6&$\PP^5$\\
7&2&4&3&$\PP^2$\\
7&3&6&0&なし\\\bottomrule
\end{tabular}
\end{center}
全双ベクトル場の空間の次元は $\binom83=56$、$\binom93=84$。
階数二の構造の射影化は $\operatorname{Gr}(3,8)$、$\operatorname{Gr}(3,9)$ で、次元は十五、十八である。

\subsection{階数二の零点：一般型および Calabi--Yau 完全交叉}
$k=1$ とし、active 座標を $x,y,z$ と書く。$m=n-1$ である。
$C=X\cap\{u=0\}$、$\Gamma=V(F_x,F_y,F_z)$、$Y=\Gamma\cap X$ とする。
\begin{proposition}\label{prop:pg-rank2}
固定した任意の滑らかな $X$ に対し、一般の三次元部分空間を active ブロックに選ぶと、
$C,\Gamma,Y$ は全て滑らかで、スキームとして
$D_0(\pi)=C\sqcup\Gamma$、$D_0(\pi_X)=C\sqcup Y$ である。
さらに
\[
\Gamma:(2,2,2),\quad\deg\Gamma=8,\quad K_\Gamma=(4-n)H,
\qquad
Y:(2,2,2,3),\quad\deg Y=24,\quad K_Y=(7-n)H.
\]
\end{proposition}
\begin{proof}
全偏微分の張る二次式の線形系は射影空間上基点をもたない。
Bertini を三回適用すると一般の三線形結合の零点 $\Gamma$ は滑らかとなる。
同じ線形系を滑らかな $X$ に制限すれば $Y$ も滑らかとなる。
一般の平面切断 $C$ の滑らかさも Bertini による。
これらは active 部分空間の Grassmann 多様体上の非空開条件なので同時に満たせる。
命題\ref{prop:pg-deg-support}、\ref{prop:pg-zero-ci}と随伴公式を適用する。
\end{proof}
$n=6$ では $\Gamma$ は $-K_\Gamma=2H$ の次数八の Fano 四次元多様体、
$Y$ は $K_Y=H$、$K_Y^3=24$ の一般型三次元多様体となる。
$n=7$ では $\Gamma$ は $-K_\Gamma=3H$ の Fano 五次元多様体、
$Y$ は次数二十四の Calabi--Yau 四次元多様体である。
ここで Calabi--Yau とは $K_Y\simeq\OO_Y$ および $H^i(Y,\OO_Y)=0$（$0<i<4$）を意味し、
Lefschetz 定理から単連結性も従う。
$h^{p,q}=\dim H^q(\Omega^p)$、$p_g=h^0(\omega)$、$e$ を位相的 Euler 数と書くと
\begin{align*}
n=6:&\quad p_g(Y)=8,\quad h^{1,1}(Y)=1,\quad h^{2,1}(Y)=246,\quad e(Y)=-504,\\
n=7:&\quad h^{1,1}(Y)=1,\quad h^{2,1}(Y)=0,\quad h^{3,1}(Y)=183,\quad h^{2,2}(Y)=780,\quad e(Y)=1152.
\end{align*}
計算は\S\ref{sec:pg-numerical}で与える。

\subsection{階数四と中心変数}
$F=f(a,u)+g(b,u)$、$a,b\in\CC^3$ とし、
$A_i=f_{a_i}$、$B_j=g_{b_j}$、$\rho_f=\sum_i a_iA_i$、$\rho_g=\sum_jb_jB_j$ とおく。
命題\ref{prop:pg-deg-ideal}は
\begin{align*}
I_{D_2(\pi)}&=(u_\alpha A_iB_j,\rho_gA_i,\rho_fB_j),\\
I_{D_0(\pi)}&=(u_\alpha A_i,u_\alpha B_j,\rho_f,\rho_g,b_jA_i,a_iB_j)
\end{align*}
を与える。
$f_0,g_0$ が滑らかな平面三次曲線 $E_f,E_g$ を定めるとし、
$\Gamma_f=V(A_1,A_2,A_3)$、$\Gamma_g=V(B_1,B_2,B_3)$、
$J=V(u,f_0,g_0)$ とおく。
$J$ は二曲線の三次元 join であり、
\[
|D_2(\pi)|=|\Gamma_f|\cup|\Gamma_g|\cup J,
\qquad |D_2(\pi_X)|=|\Sigma_f|\cup|\Sigma_g|\cup J,
\quad\Sigma_f=\Gamma_f\cap X,
\]
および同様の $\Sigma_g$ を得る。
\begin{proposition}\label{prop:pg-cubic-surface-family}
$\Sigma_f,\Sigma_g$ は型 $(2,2,2,3)$、次元 $m+1$、次数二十四の完全交叉であり、
一般のデータでは既約である。
$K_f=V(A_1,A_2,A_3)\subset\PP(a,u)$ および同様の $K_g$ を底として、
それぞれの双有理モデルは平坦な三次曲面族をなす。
相対 Jacobian イデアルで定まる fiber の特異スキームは $Z=\Delta\cap X$ である。
一般 fiber は滑らかだが全 fiber が滑らかではない。
\end{proposition}
\begin{proof}
$\Gamma_f$ は $K_f$ と反対側の $\PP_b^2$ の join であり、その斉次環は
$\CC[a,u]/(A)$ に変数 $b$ を加えた環である。
$F$ の最高 $b$ 次数部分 $g_0(b)$ は非零で係数は $\CC$ に属する。
従って随伴次数環での $g_0$ の乗法は単射であり、元の環で $F$ は非零因子となる。
補題\ref{lem:pg-critical-cover}から完全交叉性と数値を得る。

一般のデータを取り $K_f$ を滑らかとする。
join の頂点の解消として、直線をパラメータ化する
$\mathbb P_f=\PP_{K_f}(\OO_{K_f}(-1)\oplus\OO_{K_f}^{\oplus3})\to\Gamma_f$ を取る。
底の代表 $(a,u)$ を使う表示では $[t:b]\mapsto[ta:tu:b]$ であり、strict transform は
\begin{equation}
t^3f(a,u)+g(b,tu)=0.\label{eq:pg-surface-fiber}
\end{equation}
代表変更に伴い $t$ が逆に変わるので式は貼り合わさる。
$t=0$ に全成分は含まれず、余分な例外成分はない。
各 fiber の式は $t=0$ で非零の $g_0(b)$ となり、常に三次曲面を定める。
相対超曲面の完全列は fiber への制限後も左端が単射なので、族は平坦である。
$t=0$ では $g_0=0$ が滑らかなので相対特異点はない。
$t=1$ では相対特異イデアルは $(F,B_1,B_2,B_3)$ で、底の $(A_1,A_2,A_3)$ を合わせると正確に $Z$ となる。
$\Delta\to K_f$ は有限で $\dim Z=m-2$。
$K_f$ は一般に滑らかな連結完全交叉、従って既約で次元 $m-1$ なので $Z$ は底を支配しない。
一般 fiber は滑らかな三次曲面で幾何学的に既約。
平坦性により垂直な既約成分はなく、全空間とその像 $\Sigma_f$ は既約である。
$g$ 側も同様であり、$\deg Z=192>0$ から特異 fiber は存在する。
\end{proof}
一般のデータについての数値は次の通りである。
\begin{center}
\begin{tabular}{lll}\toprule
 & $n=6,m=2$ & $n=7,m=3$\\\midrule
$K_f,K_g$ & 次数8、種数5の曲線 & 次数8の K3 曲面\\
$\dim\Sigma_f$ & 3 & 4\\
$\Delta$ & 次数64、種数129の曲線 & 次数64、$K_\Delta=3H$ の曲面\\
$Z$ & 192個の被約点 & 次数192、種数577の滑らかな曲線\\
$D_0(\pi_X)$ & $E_f\sqcup E_g\sqcup Z$ & $E_f\sqcup E_g\sqcup Z$\\\bottomrule
\end{tabular}
\end{center}
$n=7$ の $\Delta$ では
$K_\Delta^2=576$、$c_2(\Delta)=768$、$\chi(\OO_\Delta)=112$、$p_g(\Delta)=111$、$h^1(\OO_\Delta)=0$。
両次元の一般葉の閉包は $\PP^6$ 内の型 $(3,3)$ の四次元完全交叉であり、
$X$ 上の底はそれぞれ $\PP(1,1,3)$、$\PP(1,1,1,3)$。
対数表示の反標準因子の次数は $(1,1,3,3)$、$(1,1,1,3,3)$ である。

\subsection{七次元の階数六}
ここでは $m=0$、$F=f(a)+g(b)+h(c)$、$a,b,c\in\CC^3$ である。
$X$ の滑らかさは各 $f,g,h$ が滑らかな平面三次曲線を定めることと同値である。
一方が特異なら対応する座標平面に $X$ の特異点が生じ、
逆に各 gradient の共通零点が原点だけなら全 gradient についても同じだからである。
曲線を $E_f,E_g,E_h$、二つずつの join を $J_{fg},J_{fh},J_{gh}$、
三曲線の join を $J_3=V(f,g,h)$ とする。
それぞれの次元は三、五である。
一次 Casimir と Poisson 超平面は存在しないが、$f,g,h$ の線形結合は Casimir である。
最高外積は
\begin{align*}
\pi^3&\longleftrightarrow\iota_E(df\wedge dg\wedge dh)\\
&=3(f\,dg\wedge dh-g\,df\wedge dh+h\,df\wedge dg)\\
&=3fgh\,d\log(g/f)\wedge d\log(h/f)\in H^0(\PP^8,\Omega^2(9)).
\end{align*}
$[f:g:h]:\PP^8\dashrightarrow\PP^2$ は $X$ 上で直線 $f+g+h=0$ に値を取る。
$f/g$ は有理 Casimir で、一般六次元葉の閉包は $\PP^8$ 内の型 $(3,3)$ の完全交叉。
$X$ 上の最高外積に対応する一形式は
$f\,dg-g\,df\in H^0(X,\Omega_X^1(6))$ である。
退化の台は
\begin{align*}
|D_4(\pi)|&=J_3\cup V(a)\cup V(b)\cup V(c),\\
|D_2(\pi)|&=J_{fg}\cup J_{fh}\cup J_{gh}\cup\PP_a^2\cup\PP_b^2\cup\PP_c^2,\\
D_0(\pi)&=E_f\sqcup E_g\sqcup E_h.
\end{align*}
$X$ 上では第一行の三つの $\PP^5$ が滑らかな三次四次元多様体となり、
第二行の三つの平面はそれぞれ曲線に制限され、既に join に含まれる。
従って $|D_2(\pi_X)|=J_{fg}\cup J_{fh}\cup J_{gh}$。

\subsection{非被約構造と primary 分解}
$A_i=f_{a_i}$、$B_j=g_{b_j}$、$C_l=h_{c_l}$ とし、
$\mathfrak a=(A_1,A_2,A_3)$、$\mathfrak b=(B_1,B_2,B_3)$、$\mathfrak c=(C_1,C_2,C_3)$ と書く。
\begin{proposition}\label{prop:pg-D4-primary}
周囲では $I_{D_4(\pi)}=(fB_jC_l,gA_iC_l,hA_iB_j)$。
$X$ 上ではイデアル層として
\begin{equation}
I_{D_4(\pi_X)}=I_{J_3}\cap\mathfrak a\cap\mathfrak b\cap\mathfrak c.\label{eq:pg-D4-primary}
\end{equation}
$J_3$ は被約成分で、他の三成分は三次四次元多様体を台とする primary な厚化である。
その一般点での横断長は八である。
\end{proposition}
\begin{proof}
生成元は命題\ref{prop:pg-deg-ideal}の $j=3$ である。
punctured affine cone の $A_i\neq0$ の開集合で $f=-g-h$ とおけば
$I_{D_4}=(g\mathfrak c,h\mathfrak b)$。
この開集合から $b,c$ 空間への写像は $a_i$ 方向の偏微分が単元なので滑らか、従って平坦である。
独立変数環 $\CC[b,c]$ では
\[
(g\mathfrak c,h\mathfrak b)=(g,h)\cap\mathfrak b\cap\mathfrak c.
\]
Euler 等式から $g\in\mathfrak b$、$h\in\mathfrak c$ なので一方の包含は明らか。
逆に $gr+hs$ が右辺に属せば、$\mathfrak b$ を法として非零多項式 $h(c)$ の乗法は単射だから $s\in\mathfrak b$。
同様に $r\in\mathfrak c$ となる。
平坦基底変換は有限交叉と可換し、全点ではいずれかの gradient 成分が単元となるので、これらを貼り合わせて\eqref{eq:pg-D4-primary}を得る。

例えば $\CC[b]/\mathfrak b$ は Hilbert 級数 $(1+t)^3$、長さ八の局所 Artin 環である。
$\mathfrak b$ による錐の商は
$\CC[a,c]/(f+h)\otimes_\CC\CC[b]/\mathfrak b$。
第一因子は滑らかな射影三次四次元多様体の錐の整域である。
第二因子の冪零根基の冪で filter すると、被約商で非零の元は非零因子であると分かる。
従って primary で、一般点の横断長は八。
$J_3$ は独立な三曲線の錐の積の射影化なので被約既約である。
\end{proof}
\begin{proposition}\label{prop:pg-D2-primary}
周囲では
\[
I_{D_2(\pi)}=(fB_j,gA_i,fC_l,hA_i,gC_l,hB_j,
 c_lA_iB_j,b_jA_iC_l,a_iB_jC_l).
\]
$X$ 上では
\begin{equation}
I_{D_2(\pi_X)}=I_{J_{fg}}\cap I_{J_{fh}}\cap I_{J_{gh}}
\cap(\mathfrak a+\mathfrak b)\cap(\mathfrak a+\mathfrak c)\cap(\mathfrak b+\mathfrak c).
\label{eq:pg-D2-primary}
\end{equation}
最後の三項はそれぞれ $E_h,E_g,E_f$ を台とする埋め込まれた primary 成分であり、各厚化の横断長は六十四である。
\end{proposition}
\begin{proof}
生成元は $j=2$ の公式による。
$A_i\neq0$ 上で $F=0$ を使うと
$I_{D_2}=(g,h,c_lB_j,b_jC_l)$。
$\mathfrak m_b=(b_1,b_2,b_3)$、$\mathfrak m_c=(c_1,c_2,c_3)$ とおくと、独立変数環で
\begin{align}
(\mathfrak m_c,g)\cap(\mathfrak m_b,h)&=(g,h,\mathfrak m_b\mathfrak m_c),\label{eq:pg-D2-reduced}\\
(g,h,\mathfrak m_b\mathfrak m_c)\cap(\mathfrak b+\mathfrak c)
&=(g,h,\mathfrak m_c\mathfrak b+\mathfrak m_b\mathfrak c).\label{eq:pg-D2-intersection}
\end{align}
第一式は $(g,h)$ で割って各因子を定数部分と正次数部分に分解すれば分かる。
第二式は $b,c$ 二重次数で確認する。
両次数が正なら $\mathfrak b+\mathfrak c$ のその成分は
$\mathfrak m_c\mathfrak b+\mathfrak m_b\mathfrak c$ に等しく、一方の次数が零なら第一式のイデアルは $(g)$ または $(h)$ であり、Euler 等式により $\mathfrak b+\mathfrak c$ に入る。
平坦基底変換により局所的に
$I_{D_2}=(\mathfrak m_c,g)\cap(\mathfrak m_b,h)\cap(\mathfrak b+\mathfrak c)$ となる。
初めの二項はそこでの $J_{fg},J_{fh}$ のイデアルである。
他の gradient 開集合でも同様に計算し、錐から $X$ へ降ろせば\eqref{eq:pg-D2-primary}を得る。

$\mathfrak b+\mathfrak c$ の商は
$\CC[a]/(f)\otimes\CC[b]/\mathfrak b\otimes\CC[c]/\mathfrak c$。
前命題と同じ filtration により primary で、横断長は $8\cdot8=64$ となる。
この成分は冗長ではない。
$b_jc_l$ は\eqref{eq:pg-D2-reduced}の右辺に属すが二重次数 $(1,1)$ をもち、$\mathfrak b+\mathfrak c$ に入らない。
台 $E_f$ は二つの join に含まれるので埋め込み成分である。他も同様。
\end{proof}
\begin{remark}\label{rem:pg-length49}
六十四は明示した primary 厚化の横断長であり、スキームと被約化との差の長さではない。
$E_f$ の一般点では
$\ker(\OO_{D_2(\pi_X)}\to\OO_{D_2(\pi_X)_{\mathrm{red}}})$ の横断長は四十九である。
実際、$Q_g=\CC[b]/\mathfrak b$、$Q_h=\CC[c]/\mathfrak c$ の極大イデアルを
$\mathfrak n_g,\mathfrak n_h$ とすると、\eqref{eq:pg-D2-reduced}、\eqref{eq:pg-D2-intersection}により核は横断方向で $\mathfrak n_g\otimes\mathfrak n_h$ に等しい。
各 Jacobian 環は長さ八で定数部分が一次元なので、長さは $7^2=49$ となる。
\end{remark}

\section{代数的な次数付き量子化}\label{sec:pg-quantization}
本節は形式的完備化だけでなく、原点を含む代数的な開集合上の量子化を構成する。
中心元の形式的存在には formality と Calabi--Yau 双対性を用い、有限次数の線形代数で代数化する。

\subsection{定義と主張}
\begin{definition}
連結次数付き可換代数 $C$ に次数零の Poisson 括弧があるとする。
$0\in U\subset\mathbb A^1_\hbar$、$K=\OO(U)$ 上の次数付き量子化とは、
$\deg\hbar=0$、$\mathcal C_0=K$、各 $\mathcal C_d$ が有限自由な
次数付き代数 $\mathcal C$ と $\mathcal C/(\hbar)\simeq C$ であって、
\[
\frac{[\widetilde a,\widetilde b]}\hbar\bmod\hbar=\{a,b\}
\]
を満たすものをいう。特殊 fiber の可換性と平坦性により持上げの選択に依存しない。
\end{definition}
連結次数付き代数 $A$ が Koszul であるとは、自明加群 $\CC=A/A_{\geq1}$ が線形自由分解をもつことをいう。
左右の injective dimension が $d$ で、左右それぞれ
\[
\Ext_A^i(\CC,A)=
\begin{cases}\CC(\ell)&i=d,\\0&i\neq d\end{cases}
\]
ならば、次元 $d$、AS parameter $\ell$ の AS-Gorenstein 代数という。
さらに大域次元 $d$ と有限 GK 次元をもてば AS-regular という。
次数移動は $M(\ell)_j=M_{j+\ell}$ とする。
以下の Ext は全内部次数を含み、injective/global dimension は通常の左右加群圏で取る。

\begin{theorem}[代数的量子化]\label{thm:pg-quantization}\label{pthm:quantization}
任意の Poisson 構造 $\pi_X$ に対し、$0\in U\subset\mathbb A^1_\hbar$、
$K=\OO(U)$ 上の二次関係式で定まる次数付き代数 $A$、および三次の中心元 $\Omega$ が存在して次を満たす。
\begin{enumerate}
\item $A$ は $S$ 上の $q$ の量子化であり、全ての閉点 $\lambda\in U$ で
$A_\lambda$ は Koszul AS-regular、
$\operatorname{gldim}A_\lambda=N$、$\operatorname{Hilb}_{A_\lambda}(t)=(1-t)^{-N}$、AS parameter $N$ である。
\item $\Omega\bmod\hbar=F$ で、各 $\Omega_\lambda$ は左右正則元である。
\item $B=A/(\Omega)$ は $S/(F)$ の量子化であり、各 $B_\lambda$ は AS-Gorenstein で、
\[
\operatorname{Hilb}_{B_\lambda}(t)=\frac{1-t^3}{(1-t)^N},\qquad
\operatorname{injdim}B_\lambda=N-1,\qquad
\text{AS parameter}=N-3=n-1.
\]
\end{enumerate}
\end{theorem}

\subsection{ブロック代数と PBW 性}
$R=\CC[u_1,\ldots,u_m]$ とし、$\sym_R$ は中心係数を外へ出して active 変数の語を平均する対称化である。
一ブロック $f(x_1,x_2,x_3,u)$ に対し
\begin{equation}
A_f=R[\hbar]\langle x_1,x_2,x_3\rangle/
\bigl([x_2,x_3]-\hbar\sym_R(f_{x_1}),\ [x_3,x_1]-\hbar\sym_R(f_{x_2}),\
[x_1,x_2]-\hbar\sym_R(f_{x_3})\bigr).\label{eq:pg-block-algebra}
\end{equation}
これは相対 cyclic potential
\begin{equation}
W=x_1x_2x_3-x_2x_1x_3-\hbar\sym_R(f)\label{eq:pg-potential}
\end{equation}
の active cyclic derivatives で定まる。
cyclic word は自由代数を交換子の線形包で割った類であり、
$\partial_i^{\mathrm{cyc}}(a_1\cdots a_r)
=\sum_{a_j=x_i}a_{j+1}\cdots a_ra_1\cdots a_{j-1}$ とする。

\begin{proposition}\label{prop:pg-PBW}
有限個のブロックに共通な $0\in U_{\mathrm{reg}}\subset\mathbb A^1$ を選び、以下を満たせる。
$g_s(x_s)=f_s(x_s,0)$ とおくと、全 $\lambda\in U_{\mathrm{reg}}$ で
\[
S_{s,\lambda}=\CC\langle x_{s1},x_{s2},x_{s3}\rangle/
([x_{s2},x_{s3}]-\lambda\sym((g_s)_{x_{s1}}),\text{循環置換})
\]
は三次元 Koszul AS-regular domain であり、3-Calabi--Yau である。
その Hilbert 級数は $(1-t)^{-3}$。
active 次数で filter し、$u_\alpha$ の filtration degree を零にすると
\[
\operatorname{gr}_xA_{f_s,\lambda}\simeq S_{s,\lambda}\otimes_\CC R.
\]
族全体にも対応する同型があり、$A_{f_s}|_{U_{\mathrm{reg}}}$ は
$R_{\mathrm{reg}}=\OO(U_{\mathrm{reg}})[u]$ 上自由である。
\end{proposition}
\begin{proof}
一つの $g$ について関係式の列は $M_\lambda x$ で表される。ただし
\[
M_\lambda=
\begin{pmatrix}
-\frac\lambda2g_{11}&-x_3-\frac\lambda2g_{12}&x_2-\frac\lambda2g_{13}\\
x_3-\frac\lambda2g_{21}&-\frac\lambda2g_{22}&-x_1-\frac\lambda2g_{23}\\
-x_2-\frac\lambda2g_{31}&x_1-\frac\lambda2g_{32}&-\frac\lambda2g_{33}
\end{pmatrix}.
\]
反対称成分が三つの交換子を与えるので関係式は独立である。
potential 由来の表示は\cite[Corollary 2.7]{PMoriSmith}の意味で standard であり、
同論文 Theorem 1.4 によって AS-regular 性は全 $\PP^2$ 上で $\rank M_\lambda\geq2$ と同値である。
階数一以下の incidence は $\mathbb A^1\times\PP^2$ の閉集合で、その底への像も閉じている。
$M_0$ は全 $\PP^2$ 上で階数二の交代行列なので、像を除けば原点を含む開集合を得る。
有限個のブロックについて交わりを取る。
3-CY 性と domain 性は\cite[Corollary 2.9]{PMoriSmith}による。
三生成元三二次関係式の三次元 AS-regular 代数では、最小分解を双対化すると最後の shift は三となり、
\[
0\longrightarrow S_{s,\lambda}(-3)\longrightarrow S_{s,\lambda}(-2)^3
\longrightarrow S_{s,\lambda}(-1)^3\longrightarrow S_{s,\lambda}
\longrightarrow\CC\longrightarrow0.
\]
これが Koszul 性と Hilbert 級数を与える。

$\lambda$ を固定する。$u=a$ への特殊化では、\eqref{eq:pg-potential}は
$g$ の三次 potential に active 次数二以下の項を加えたものである。
\cite[Theorem 1.1(i)]{PBergerTaillefer}により、その代数は $S_{s,\lambda}$ の PBW 変形である。
$S_{s,\lambda}$ の斉次基底を語の線形結合として持ち上げる。
最高 active 次数の関係式から次数帰納法で、これらは $A_{f_s,\lambda}$ を $R$ 上生成する。
線形関係があれば全 $u=a$ で PBW 基底間の関係となり、係数多項式が全点で零だから元の関係も零となる。
従って所要の PBW 性と自由性を得る。

族の斉次 $S_{s,U_{\mathrm{reg}}}$ の各次数部分は底上有限生成で、全 fiber の次元が一定である。
底環は $\CC[\hbar]$ の局所化で PID なので、各部分は有限自由である。
torsion があればその台で fiber 次元が増えるからである。
その各次数の基底を自由代数へ持ち上げ、同じ帰納法と全 $(\lambda,a)$ への特殊化を使えば、族上の自由性も従う。
\end{proof}
\begin{remark}\label{rem:pg-regular-all}
$g_s$ が滑らかな平面三次曲線を定めるとき、
\cite[Theorem 4.2(1)]{PMoriSmith}により $S_{s,\lambda}$ の AS-regular 性のために除く $\lambda$ はない。
$g_s=0$ の場合も常に多項式環である。
中心元の代数化・正則性のためには、以下でさらに開集合を縮める。
ambient の AS-regular 性は、その中心超曲面商の圏論的滑らかさとは別条件である。
\end{remark}
以下
$A_{\mathrm{reg}}=A_{f_1}\otimes_{R_{\mathrm{reg}}}\cdots\otimes_{R_{\mathrm{reg}}}A_{f_k}$
とする。異なるブロックは可換で、各全次数部分は底上有限自由であり、
$A_{\mathrm{reg}}/(\hbar)\simeq S$。関係式の一次項は\eqref{eq:pg-bracket}である。

\subsection{形式的な中心元}
この小節では $\widehat A_f$、$\widehat R=R[[\hbar]]$ という $\hbar$ 進完備化を用い、
$\widehat\otimes$、双加群分解および $\RHom$ も完備化の意味で解釈する。
内部次数を忘れた $d$-Calabi--Yau 性とは、包絡代数上 perfect で、包絡代数への導来双対が代数自身の $[-d]$ であることをいう。
以下では内部 shift も記録する。
\begin{lemma}\label{lem:pg-formal-center}
任意のブロックに対して、三次の中心元
$\widehat\Omega_f\in Z(\widehat A_f)$ で
$\widehat\Omega_f\bmod\hbar=f$ を満たすものが存在する。
$f(x,0)$ の滑らかさや孤立特異点は仮定しない。
\end{lemma}
\begin{proof}
\cite[Proposition 8.1.1(i)の証明]{EG}の方法を相対係数へ適用するため、
まず絶対 CY 性を確認する。
$A=\widehat A_f$、$r_i=\partial_i^{\mathrm{cyc}}W$ と書く。
相対完備自由双加群複体
\begin{equation}
0\longrightarrow P_3\xrightarrow{d_3}P_2\xrightarrow{d_2}P_1
\xrightarrow{d_1}P_0\xrightarrow\mu A\longrightarrow0\label{eq:pg-potential-resolution}
\end{equation}
を、
\[
P_0=A\widehat\otimes_{\widehat R}A,\quad
P_1=\bigoplus_{i=1}^3Ae_iA,\quad
P_2=\bigoplus_{i=1}^3Av_iA,\quad P_3=AwA
\]
で定める。$\deg e_i=1$、$\deg v_i=2$、$\deg w=3$ であり、
\[
d_1(e_i)=x_i\otimes1-1\otimes x_i,\quad
 d_3(w)=\sum_i(x_iv_i-v_ix_i),\quad
 d_2(v_i)=\sum_j\sum_{r_i\ni ax_jb}ae_jb.
\]
最後の和は語の各出現を係数付きで数える。
cyclic derivative の恒等式 $\sum_i[x_i,r_i]=0$ を非可換微分すると $d_2d_3=0$。
また $d_1d_2(v_i)=r_i\otimes1-1\otimes r_i=0$。
$\hbar=0$ では対角の相対 Koszul 分解である。
cycle を法 $\hbar$ で boundary として持ち上げ、差を $\hbar$ で割る操作を反復すると、
平坦性と完備性により完全性が従う。左端の核も全 $\hbar^rP_3$ に入り零である。

$\partial_j^{\mathrm{nc}}r_i=\sum_{r_i\ni ax_jb}a\otimes b$ とし
$\tau(a\otimes b)=b\otimes a$ と書く。
cyclic word の二切断位置を交換すれば
$\tau(\partial_j^{\mathrm{nc}}r_i)=\partial_i^{\mathrm{nc}}r_j$。
これが中央微分の自己双対性を与え、両端も符号を除き双対である。
逆順の双対微分は $(-d_1,d_2,-d_3)$ となり、各項に $(1,-1,-1,1)$ を掛けると元の複体に一致する。
従って $D_{\mathrm{rel}}=A\widehat\otimes_{\widehat R}A^{\mathrm{op}}$ に対し
\[
\Ext_{D_{\mathrm{rel}}}^q(A,D_{\mathrm{rel}})=
\begin{cases}A(3)&q=3,\\0&q\neq3.\end{cases}
\]
automorphism による twist はない。
絶対包絡代数 $D_{\mathrm{abs}}=A\widehat\otimes_{\CC[[\hbar]]}A^{\mathrm{op}}$ は、PBW 基底により
中心係数 $\CC[u^L,u^R][[\hbar]]$ 上平坦である。
従って $u_\alpha^L-u_\alpha^R$ は中心正則列で、その商は $D_{\mathrm{rel}}$。
その Koszul 分解と change of rings により
\[
\Ext_{D_{\mathrm{abs}}}^q(A,D_{\mathrm{abs}})=
\begin{cases}A(m+3)&q=m+3,\\0&q\neq m+3.\end{cases}
\]
相対分解の各項は $D_{\mathrm{abs}}$ 上 perfect なので、絶対的にも $(m+3)$-CY である。

対称化は完備化後に次数を保つ $\widehat R$-加群同型
$R[x_1,x_2,x_3][[\hbar]]\simeq A$ を与える。
積を移せば相対 Hochschild cochain による形式変形となる。
相対 polydifferential cochain の包含は、対角の Koszul 分解による HKR 計算から quasi-isomorphism である。
Kontsevich の affine formality のグラフ公式\cite[\S4.6, \S6]{PKontsevich}で active 変数の微分だけを用いれば、$R$ 上の formality を得る。
全変数の同時スケーリングの重み零の部分に制限する。
bar--Koszul 比較を次数付きで取れるので HKR 同型と包含はこの部分でも quasi-isomorphism であり、
Kontsevich 写像の各 Taylor 成分も重みを保つ。
$\hbar$ 倍して完備化した pronilpotent dg Lie 代数の Maurer--Cartan 同値を使えば、
形式 Poisson 構造と gauge transformation をともに重み零で選べる。
従って
\[
\Pi=\hbar\pi_f+O(\hbar^2),\qquad
\pi_f=\sum_{\mathrm{cyc}}f_{x_1}\partial_{x_2}\wedge\partial_{x_3}
\]
という相対形式 Poisson 構造の star product が $A$ と次数付き $R$-線形同型をもつ。
$\Pi$ は $\partial_{u_\alpha}$ 成分をもたないため、絶対グラフ公式で現れる微分も active 変数だけであり、
同じ $\Pi$ が絶対 Kontsevich class の代表となる。

\cite[Theorem 3, Proposition 1]{PDolgushev}を適用する。
台は滑らかな affine space で標準束は自明であり、双対化双加群が自明であることは上で確認した。
よって標準体積形式
$\omega_0=dx_1\wedge dx_2\wedge dx_3\wedge du_1\wedge\cdots\wedge du_m$ に対し、
単元 $\rho\in R[x][[\hbar]]^\times$ で $L_\Pi(\rho\omega_0)=0$ となるものが存在する。
全次数 $m+3$ の成分を取る。
$\Pi$ の重みは零であり、$\rho$ の次数零成分は $c(\hbar)\in\CC[[\hbar]]^\times$ だから
$L_\Pi\omega_0=0$ となる。
ここで単元性は $\rho\bmod\hbar$ が多項式環の単元、従って非零定数であることによる。

$\nu=dx_1\wedge dx_2\wedge dx_3$ とすると相対一形式
$\alpha=\iota_\Pi\nu$ は $d_R\alpha=0$ を満たす。
相対 polynomial de Rham 複体の active Euler homotopy によって
$\alpha=d_R\psi$、$\psi=\hbar f+O(\hbar^2)$、$\deg\psi=3$ を得る。
積分定数を $R[[\hbar]]$ で調整して最初の係数を $f$ にする。
$\Pi$ は $\psi$ の相対 Jacobian 場なので $\psi/\hbar$ は Casimir である。
twisted formality の第一 Taylor 成分はこの零次 Poisson cocycle を零次 Hochschild cocycle、すなわち中心元に送る。
gauge 同型で $A$ に戻したものが $\widehat\Omega_f$。
全写像が次数を保ち、法 $\hbar$ で関数上の HKR 写像は恒等なので、その次数は三、residue は $f$ である。
\end{proof}
\begin{remark}
\cite[Proposition 8.1.1]{EG}の孤立特異点の追加仮定は中心全体・cohomologyを扱う (ii),(iii) に課される。
上の証明は中心元一つの存在を扱い、中心全体が一つの potential で生成されるとは主張しない。
相対 formality および絶対 CY 双対性という外部入力を明示しており、計算実験だけから中心元を推測する議論ではない。
\end{remark}

\subsection{中心元の代数化と正則中心拡大}
\begin{lemma}\label{lem:pg-algebraize-center}
$0\in U_0\subset\mathbb A^1_\hbar$ 上の次数付き代数 $D$ の各次数部分が有限自由であり、
$D$ が次数一の元 $y_1,\ldots,y_r$ と中心元で生成されるとする。
完備化に次数 $d$ の中心元 $\widehat\eta$ があれば、
ある $0\in U\subset U_0$ 上で中心元 $\eta\in D_d$ を
$\eta\bmod\hbar=\widehat\eta\bmod\hbar$ となるように選べる。
\end{lemma}
\begin{proof}
$O=\CC[\hbar]_{(\hbar)}$ として有限自由加群間の写像
$(D_O)_d\to(D_O)_{d+1}^{\oplus r}$、$a\mapsto([a,y_i])_i$ の核を $L$ とする。
完備化の平坦性から完備化後の核は $L\otimes_O\CC[[\hbar]]$。
$L$ の有限個の生成元 $l_j$ を用いて
$\widehat\eta=\sum_ja_j(\hbar)l_j$ と書けば
$\eta_O=\sum_ja_j(0)l_j$ は同じ residue をもつ。
係数の分母は原点で消えないので、それらの零点を除けばよい。
\end{proof}
\begin{lemma}\label{lem:pg-central-linear}
$D$ を連結次数付き代数、$z\in D_1$ を中心正則元とする。
$C=D/(z)$ が次元 $d$、AS parameter $d$ の Koszul AS-regular なら、
$D$ は次元 $d+1$、AS parameter $d+1$ の Koszul AS-regular である。
\end{lemma}
\begin{proof}
$0\to D(-1)\xrightarrow{z}D\to C\to0$ より、
$\operatorname{Tor}^D_q(C,\CC)$ は $q=0,1$ でそれぞれ $\CC,\CC(-1)$、他は零である。
change-of-rings spectral sequence
\[
E^2_{p,q}=\operatorname{Tor}^C_p(\operatorname{Tor}^D_q(C,\CC),\CC)
\Longrightarrow\operatorname{Tor}^D_{p+q}(\CC,\CC)
\]
の各項は内部次数 $p+q$ に集中する。
微分は内部次数を保ち全ホモロジー次数を一減らすため全て零であり、
$D$ は Koszul、$j>d+1$ で $\operatorname{Tor}^D_j(\CC,\CC)=0$、$j=d+1$ では非零である。
$D^!$ を二次双対とすると、Koszul 双加群複体の項は $D\otimes(D_i^!)^*\otimes D$。
全内部次数を固定し右の $D$ 次数で filter すれば有限 filtration となり、
随伴複体の正ホモロジーは左 Koszul 分解で消える。
よって増大双加群複体は長さ $d+1$ の完全列であり、右加群として分裂完全なので
任意の左加群とテンソルしても完全である。従って通常の大域次元も $d+1$。
さらに $\RHom_D(C,D)\simeq C(1)[-1]$ と導来随伴から
\[
\Ext_D^{i+1}(\CC,D)\simeq\Ext_C^i(\CC,C)(1).
\]
右加群でも同様で、左右 AS 条件と parameter を得る。
$\operatorname{Hilb}_D=\operatorname{Hilb}_C/(1-t)$ なので GK 次元も有限である。
\end{proof}
\begin{lemma}\label{lem:pg-central-cubic}
$A$ を次元 $N$、AS parameter $N$ の AS-regular 代数、$\eta\in A_3$ を中心正則元とする。
$B=A/(\eta)$ は左右 injective dimension $N-1$、AS parameter $N-3$ の AS-Gorenstein 代数で、
$\operatorname{Hilb}_B=(1-t^3)\operatorname{Hilb}_A$ である。
\end{lemma}
\begin{proof}
$0\to A(-3)\xrightarrow{\eta}A\to B\to0$ が Hilbert 級数を与える。
導来随伴により任意の左 $B$-加群 $M$ について
$\Ext_B^i(M,B)\simeq\Ext_A^{i+1}(M,A)(-3)$。
右辺は $i\geq N$ で零、$M=\CC$、$i=N-1$ で $\CC(N-3)$ となる。
よって injective dimension はちょうど $N-1$ で AS 条件も成立する。右側も同様。
\end{proof}

\subsection{代数的量子化定理の証明}
\begin{proof}[定理\ref{thm:pg-quantization}の証明]
零構造は $A=S[\hbar]$、$\Omega=F$ でよい。
非零構造に対して $A_{\mathrm{reg}}$ を構成する。
PBW 基底により各 fiber で $u_1,\ldots,u_m$ は中心正則列で、
\[
A_\lambda/(u_1,\ldots,u_m)\simeq\bigotimes_{s=1}^kS_{s,\lambda}.
\]
右辺の自明加群の分解、双加群分解およびその双対は各因子のもののテンソル積で得られる。
従って次元 $3k$、AS parameter $3k$ の Koszul AS-regular 代数である。
補題\ref{lem:pg-central-linear}を $m$ 回適用して $A_\lambda$ の次元 $N=3k+m$ と AS parameter $N$ を得る。
Hilbert 級数は PBW 基底から $(1-t)^{-N}$ である。

補題\ref{lem:pg-formal-center}と\ref{lem:pg-algebraize-center}を各ブロックに適用し、開集合を共通に縮めると
$\Omega_s\in Z(A_{f_s}|_U)_3$、$\Omega_s\bmod\hbar=f_s$ を得る。
異なるブロックは可換なので $\Omega=\sum_s\Omega_s$ は $A$ の中心元で、residue は $F$。

少なくとも一つの $f_s$ は active 変数に依存する。
そうでなければ $F$ は中心変数だけに依存し、active 座標点に特異点をもつ。
その $s$ を固定し、PBW 基底での $\Omega_s$ の active 成分のうち原点で非零の係数が消えないようさらに縮める。
各 $\lambda$ で $\Omega_{s,\lambda}$ の active 次数は正となる。
最高次数が $\lambda$ によらず一定である必要はない。
$\lambda$ を固定し $x_s$ 次数だけで filter すると
\[
\operatorname{gr}_sA_\lambda\simeq
(S_{s,\lambda}\otimes_\CC R)\otimes_R A_{\mathrm{rest},\lambda},
\qquad A_{\mathrm{rest},\lambda}=\bigotimes_{t\neq s,R}A_{f_t,\lambda}.
\]
他の $\Omega_t$ は filtration degree 零なので、$\Omega_\lambda$ の最高記号は
$\Omega_{s,\lambda}$ の非零最高記号 $a_s\in S_{s,\lambda}\otimes R$ である。
$S_{s,\lambda}$ は domain だから多項式拡大も domain で、$a_s$ の左右乗法は単射。
$R$-自由な $A_{\mathrm{rest},\lambda}$ とテンソルしても単射なので、最高次数を比較して
$\Omega_\lambda$ の左右正則性が従う。非可換 domain のテンソル積が domain であるとは仮定していない。

従って各 fiber に
$0\to A_\lambda(-3)\xrightarrow{\Omega_\lambda}A_\lambda\to B_\lambda\to0$
があり、所要の Hilbert 級数を得る。
各次数の余核は PID $\OO(U)$ 上で全 fiber 次元が一定なので有限自由である。
法 $\hbar$ での商は $S\to S/(F)$ で、交換子の一次項もその商の括弧と一致する。
補題\ref{lem:pg-central-cubic}が残りの AS-Gorenstein の主張を与える。
\end{proof}
二次関係式は各生成元対に一つで、反対称成分により独立なので個数は $\binom N2$ である。
六次元では生成元八、関係式二十八、大域次元八、商の injective dimension 七、AS parameter 五。
七次元ではそれぞれ九、三十六、九、八、六となる。
\begin{remark}\label{rem:pg-quantization-scope}
定理\ref{thm:pg-quantization}は、個々の Poisson 構造に対する量子化の存在定理である。
代数の同型分類、中心全体の記述、$\qgr B_\lambda$ の滑らかさは含まない。
とくに本定理だけを根拠に全パラメータの Kuznetsov 成分の smooth proper 性や可換復元を結論することはできない。
\end{remark}

\section{特殊なブロック制限と Fermat 型}\label{sec:pg-special}
\subsection{零の三変数制限}
\begin{proposition}\label{prop:pg-zero-restriction}
$X$ が滑らかで、あるブロックの制限 $f_s(x_s,0)$ が恒等的に零ならば $m\geq3$ である。
$m=3$ ではこのような滑らかな例が存在する。
\end{proposition}
\begin{proof}
対応する座標平面 $\PP_s^2$ 上では全 active 偏微分が零であり、
残る偏微分は高々 $m$ 個の二次式 $F_{u_\alpha}(x_s,0)$ である。
$m\leq2$ なら三変数多項式環でこれらのイデアルの高さは高々二なので、非零の共通零点をもつ。
そこで $F$ の全偏微分が消え、滑らかさに反する。

存在について、九変数 $a_1,a_2,a_3,b_1,b_2,b_3,u_1,u_2,u_3$ の三次式
\begin{equation}\label{eq:pg-zero-restriction-example}
F=\sum_{i=1}^3(u_i a_i^2+u_i^3)+b_1^3+b_2^3+b_3^3
\end{equation}
を考える。全偏微分が零なら $2u_ia_i=0$ および $a_i^2+3u_i^2=0$ から
$a_i=u_i=0$、さらに $3b_i^2=0$ から $b_i=0$ となる。従って射影超曲面は滑らかである。
これは $k=2,m=3$ の正規形で、$a$ ブロックの $u=0$ 制限は零である。
\end{proof}
この例にも正規形定理\ref{thm:pg-normal}、全退化イデアルの命題\ref{prop:pg-deg-ideal}、
量子化定理\ref{thm:pg-quantization}を適用できる。
一方、平面楕円曲線を用いた零点の非交和表示の仮定は満たさないため、
退化スキームは一般式\eqref{eq:pg-deg-generators}から計算する必要がある。

\subsection{Fermat 型の中心元}
\begin{lemma}\label{lem:pg-fermat-center}
結合代数において
\[
[x,y]=\hbar z^2,\qquad [y,z]=\hbar x^2,\qquad [z,x]=\hbar y^2
\]
が成り立つとする。$S=x^3+y^3+z^3$、$T=\sym(xyz)$ とおくと $S-2\hbar^2T$ は中心元である。
ここで $\sym(xyz)$ は六個の語の平均である。
\end{lemma}
\begin{proof}
関係式と積の交換子を展開すると
\[
[S,x]=\hbar(zy^2z-yz^2y)=\hbar[zy,yz]=\hbar^2[zy,x^2]
\]
を得る。一方、$6T$ の六個の語を関係式で並べ替えると
$2T=xzy+zyx+\hbar x^3$ となる。従って $2[T,x]=[zy,x^2]$ であり、
$[S-2\hbar^2T,x]=0$ を得る。循環置換により $y,z$ との交換子も零である。
\end{proof}
特に
\begin{equation}\label{eq:pg-fermat}
F=\frac13\left(\sum_{\alpha=1}^m u_\alpha^3+
\sum_{s=1}^k\sum_{i=1}^3x_{si}^3\right)
\end{equation}
に対する量子化では、各ブロックで $[x_{s1},x_{s2}]=\hbar x_{s3}^2$ とその循環置換を課し、
異なるブロックと中心変数は可換とする。このとき
\begin{equation}\label{eq:pg-fermat-center}
\Omega=\frac13\left[\sum_{\alpha=1}^m u_\alpha^3+
\sum_{s=1}^k\left\{\sum_{i=1}^3x_{si}^3-2\hbar^2\sym(x_{s1}x_{s2}x_{s3})\right\}\right]
\end{equation}
は全パラメータで中心元であり、$\Omega|_{\hbar=0}=F$ を満たす。
定理\ref{thm:pg-quantization}の適切な開集合上で中心正則超曲面の族を与える。

\subsection{Fermat 型の非被約な零点}
\begin{proposition}\label{prop:pg-fermat-thickening}
\eqref{eq:pg-fermat}で $m>0$ ならば $\Delta=V(x_{si}^2\mid s,i)$ は
中心線形空間 $\PP_u^{m-1}$ を台とする横断長 $8^k$ の厚化である。
$m\geq2$ ならば $Z=\Delta\cap X$ は
$V(\sum_\alpha u_\alpha^3)\subset\PP_u^{m-1}$ の同じ横断厚化となる。
$m=0$ の場合には $\Delta=\varnothing$ である。
\end{proposition}
\begin{proof}
active 偏微分は $A_{s,i}=x_{si}^2$ なので、その商環は
\[
\CC[u]\otimes_\CC\bigotimes_{s=1}^k
\CC[x_{s1},x_{s2},x_{s3}]/(x_{s1}^2,x_{s2}^2,x_{s3}^2)
\]
である。各 Artin 因子の長さは八であり、この商では全 $x_{si}^3$ が零なので、
$F$ の像は $\frac13\sum_\alpha u_\alpha^3$ となる。これが厚化の記述を与える。
$m=0$ では偏微分の共通零点はアフィン原点だけであり、射影化は空である。
\end{proof}
六次元・階数四では、$Z$ は中心 $\PP^1$ 上の三点の各々に長さ六十四の厚みをもつ。
七次元・階数四では、中心 $\PP^2$ 上の Fermat 楕円曲線に横断長六十四の厚みをもつ。
次数はともに百九十二で、後者の算術種数は五百七十七である。
従って一般の場合との違いは次数や算術種数には現れず、被約性と滑らかさに現れる。

\section{完全交叉の数値不変量}\label{sec:pg-numerical}
本節は\S\ref{sec:pg-six-seven}で用いた数値を導出する。
滑らかな完全交叉 $W=(d_1,\ldots,d_c)\subset\PP^M$ に対し、
$H=c_1(\OO_W(1))$ とおくと
\begin{equation}\label{eq:pg-ci-formulas}
\omega_W\simeq\OO_W\left(\sum_i d_i-M-1\right),\qquad
c(T_W)=\frac{(1+H)^{M+1}}{\prod_i(1+d_iH)},\qquad
\int_W H^{M-c}=\prod_i d_i.
\end{equation}
これは順に随伴公式、接束の完全列と Euler 列、完全交叉の次数公式である。
また Koszul 分解から
$H^i(W,\OO_W(t))=0$（$0<i<\dim W$、$t\in\ZZ$）を得る。
以下の滑らかさを要する Hodge 数の主張は、一般のデータ、または滑らかさを個別に確認したデータに対して適用する。

\subsection{階数二の零点成分}
$n=6$ で $Y=(2,2,2,3)\subset\PP^7$ の場合、
\[
K_Y=H,\qquad c(T_Y)=\frac{(1+H)^8}{(1+2H)^3(1+3H)}
=1-H+7H^2-21H^3.
\]
従って $K_Y^3=24$、$e(Y)=-21\cdot24=-504$ である。
全方程式が二次以上なので $p_g(Y)=h^0(\OO_Y(1))=8$。
Lefschetz 定理と Hodge 分解により
$h^{1,1}=1$、$h^{1,0}=h^{2,0}=0$ であり、
\[
-504=e(Y)=4-2p_g(Y)-2h^{2,1}(Y)
\]
から $h^{2,1}(Y)=246$ を得る。

$n=7$ で $Y=(2,2,2,3)\subset\PP^8$ の場合は
\[
K_Y=0,\qquad c(T_Y)=\frac{(1+H)^9}{(1+2H)^3(1+3H)}
=1+6H^2-14H^3+48H^4,
\]
従って $e(Y)=48\cdot24=1152$ である。
Koszul 分解は $H^i(\OO_Y)=0$（$0<i<4$）を、
Lefschetz 定理は $h^{1,1}=1$、$h^{2,1}=0$ を与える。
$h^{3,1}$ を求めるため normal sequence を使う。
次数二の方程式は三個、次数三の方程式はそれらの一次式倍と独立な三次式一個だから
\[
h^0(\OO_Y(2))=\binom{10}{2}-3=42,\qquad
h^0(\OO_Y(3))=\binom{11}{3}-3\cdot9-1=137.
\]
Euler 列を $Y$ に制限し上記の消滅を使うと
$h^0(T_{\PP^8}|_Y)=9\cdot9-1=80$、$H^1(T_{\PP^8}|_Y)=0$ である。
$K_Y\simeq\OO_Y$ による $T_Y\simeq\Omega_Y^3$、$H^3(\OO_Y)=0$ と Hodge 対称性から
$H^0(T_Y)=0$ である。従って
\[
0\longrightarrow T_Y\longrightarrow T_{\PP^8}|_Y
\longrightarrow\OO_Y(2)^{\oplus3}\oplus\OO_Y(3)\longrightarrow0
\]
は $h^{3,1}(Y)=h^1(T_Y)=3\cdot42+137-80=183$ を与える。
最後に Hodge 対称性と Poincar\'e 双対性により
$1152=e(Y)=6+2h^{3,1}(Y)+h^{2,2}(Y)$、従って $h^{2,2}(Y)=780$ である。

\subsection{階数四の critical covers と零点}
$m=2$ の $K_f,K_g$ は型 $(2,2,2)$ の $\PP^4$ 内の曲線である。
従って次数八、$K_{K_f}=H$、$2g(K_f)-2=8$ で、種数五を得る。
$m=3$ では $\PP^5$ 内の同じ型の曲面であり、標準束が自明で $H^1(\OO)=0$ なので K3 曲面である。

$n=6,k=2,m=2$ の場合、$\Delta=(2^6)\subset\PP^7$ は次数六十四の曲線で $K_\Delta=4H$。
従って $2g(\Delta)-2=4\cdot64=256$、$g(\Delta)=129$ である。
$Z=(2^6,3)$ は長さ $3\cdot64=192$ の完全交叉で、一般のデータでは滑らか、すなわち百九十二個の被約点となる。

$n=7,k=2,m=3$ の場合、$\Delta=(2^6)\subset\PP^8$ は次数六十四の曲面で、
\[
c(T_\Delta)=\frac{(1+H)^9}{(1+2H)^6}=1-3H+12H^2.
\]
よって
\[
K_\Delta^2=9\cdot64=576,\qquad c_2(\Delta)=12\cdot64=768,\qquad
\chi(\OO_\Delta)=\frac{576+768}{12}=112.
\]
$H^1(\OO_\Delta)=0$ なので $p_g(\Delta)=111$ であり、これは
$h^0(\OO_\Delta(3))=\binom{11}{3}-6\cdot9=111$ からも確認できる。
最後に $Z=(2^6,3)\subset\PP^8$ は次数百九十二の曲線で $K_Z=6H$ だから
\[
2g(Z)-2=6\cdot192=1152,\qquad g(Z)=577.
\]
この式は特殊な非被約完全交叉でも算術種数の公式として保たれる。

% 2026-11-11 / Bondal's conjecture for smooth cubic hypersurfaces
\section{Poisson正規形から従うBondal予想}
\label{sec:pg-bondal}

Bondal予想はFano多様体の退化軌跡
\[
 D_{2j}(\pi)=\{p:\operatorname{rank}(\pi_p)\le2j\}
 =V(\pi^{j+1})
\]
が、$0\le 2j<\dim X$ の各$j$に対して次元$2j+1$以上の
既約成分を持つと主張する
\cite{BondalGualtieriPym,BondalDruelPereiraPymTouzet}。
ここでは$\S\ref{sec:pg-normal}$の正規形から、cubicに対して
全ての退化階数を一度に扱う。高次元の一般のFano多様体に対する主張ではない。

\begin{theorem}[滑らかなcubicに対するBondal予想]
\label{thm:pg-bondal}
$X=V(F)\subset\mathbb P^{n+1}_{\mathbb C}$を$n\ge3$の滑らかな三次超曲面、
$\pi$を任意の正則Poisson構造とする。このとき$0\le 2j<n$の各$j$について
\[
 \boxed{\quad \dim D_{2j}(\pi)\ge 2j+1.\quad}
\]
より強く、$\pi\ne0$の一般階数を$2k$とすると、任意の$0\le j<k$に対して
次元がちょうど$2j+1$の既約なjoin部分多様体
$J_{j+1}\subset D_{2j}(\pi)$を構成できる。
従って$D_{2j}(\pi)$は次元$2j+1$以上の既約成分を持つ。
\end{theorem}

\begin{proof}
$\pi=0$なら$D_{2j}(\pi)=X$なので結論する。
非零の場合、定理\ref{thm:pg-normal}を使い、
$N=n+2=3k+m$、$V=\bigoplus_{s=1}^kW_s\oplus U$、
$\dim W_s=3$とし、
\[
 F=\sum_{s=1}^kf_s(x_s,u),\qquad
 q=\sum_{s=1}^k
 \iota_{dF}
 (\partial_{x_{s1}}\wedge\partial_{x_{s2}}
  \wedge\partial_{x_{s3}})
\]
を取る。一般階数は$2k$である。
$j\ge k$なら$D_{2j}(\pi)=X$、また$2j<n$から
$\dim X=n\ge2j+1$なのでよい。

$0\le j<k$とし、$r=j+1$個の相異なるブロックを選ぶ。
その各々で$g_s(x_s):=f_s(x_s,0)$とする。
$\mathbb P(W_s)\cong\mathbb P^2$の零軌跡$V(g_s)$は、
$g_s\ne0$なら少なくとも一つの既約な射影曲線$C_s$を含む。
$g_s=0$の場合も任意の直線$C_s\subset\mathbb P(W_s)$を選べる。
従って両場合に
\[
 C_s\subset V(g_s)\subset\mathbb P(W_s),\qquad
 \dim C_s=1
\]
となる既約曲線を選べる。
それらの射影joinを
\[
 J_r=\operatorname{Join}(C_1,\ldots,C_r)
       \subset\mathbb P\Bigl(\bigoplus_{s=1}^r W_s\Bigr)
       \subset\mathbb P(V)
\]
と定める。
各$C_s$のアフィン錐は既約で次元2、また座標空間$W_s$は直和をなす。
従って錐の直積から得る$J_r$は既約かつ
\[
 \dim J_r=\sum_{s=1}^r(\dim C_s+1)-1=2r-1=2j+1.
\]
$J_r$の点では、選んだブロック上で各$g_s=0$、
他のブロックおよび$u$は零なので$F=0$である。
よって$J_r\subset X$。

次にランクを評価する。アフィン錐の任意の非零代表
$z=(x_1,\ldots,x_r,0,\ldots,0;u=0)$を取ると、
$q(z)$は選んだ$r$ブロックのJacobian bivector $q_s$の和で、
各$q_s$の階数は0または2である。
いずれかの$q_s(z)$が零なら
$\operatorname{rank}q(z)\le2r-2$であり、
射影化と$X$への制限でrankは増えないから
$\operatorname{rank}\pi_{[z]}\le2r-2$である。

残りの場合、選んだ各ブロックで$q_s(z)$は階数2であり、
その像は$\ker dg_s(x_s)\subset W_s$である。
$g_s(x_s)=0$と三次Euler等式から
$dg_s(x_s)(x_s)=3g_s(x_s)=0$となるので、
各ブロックのEulerベクトル$E_s=x_s$は$q_s(z)$の像に入る。
従って$E(z)=\sum_{s=1}^r E_s$は$q(z)$の像に入る。
$q(z)$は$F$をCasimirに持ち、接空間$T_z(F=0)$上の双ベクトルである。
射影接空間$T_{[z]}X=T_z(F=0)/\mathbb CE(z)$への射影により、
像の中の非零一次元$\mathbb CE(z)$を割ると、交代形式の階数は2減少する。
よって再び
\[
 \operatorname{rank}\pi_{[z]}\le2r-2=2j.
\]
以上から$J_r\subset D_{2j}(\pi)$を得る。
$D_{2j}(\pi)$はNoether閉集合なので、既約な$J_r$はある既約成分に含まれ、
その次元は少なくとも$2j+1$である。
\end{proof}

\begin{remark}[証明の適用範囲と既存の退化スキーム計算]
\label{rem:pg-bondal-scope}
本定理は退化軌跡の\emph{台}の次元下界であって、
各$D_{2j}$のscheme構造の被約性・primary分解を与えるものではない。
そのような詳細は前節の具体計算と区別する。
$C_s$が滑らかな平面三次曲線なら$J_r$はその曲線のjoinである。
$g_s=0$や特異な場合にも上の次元証明は成立する。
また既知のFano四次元までの結果と、
2026年の五次元Fanoについての一般定理
\cite{BondalDruelPereiraPymTouzet}との重複を明記する。
六次元以上のcubicにおける結論の先行研究網羅調査は未完了であり、
今回の証明の新規性は未確認である。
最終的な主張は正規形定理\ref{thm:pg-normal}に依存し、
その全証明を今回改めて独立検証したわけではない。
\end{remark}

\input{papers/sections/poisson_moduli_stacks}
\section{Fermat型の固定Poissonスキームと相対Poissonモジュライの既約成分}
\label{sec:pg-fermat-relative-components}
本節では、\emph{固定した滑らかなcubicのPoisson scheme}と、
\emph{cubic自身も変形させる相対Poisson scheme}を比較する。
いずれの「既約成分」もスキームの基礎位相空間（すなわち被約化）の既約成分を意味し、
その非被約構造や成分交叉のイデアルについての主張ではない。
Poisson双ベクトルのスカラーは割らず、相対側でのみ射影座標変換を割る。
本節の分類は定理\ref{thm:pg-normal}に依存し、その証明全体の独立再監査・
先行研究上の新規性確認は完了していない。

\subsection{対象と二つの成分数}
$N=n+2\ge5$、$V=\mathbb C^N$、$S=\operatorname{Sym}(V^*)$とする。
$F\in S_3$が滑らかな射影cubic $X_F\subset\mathbb P(V)$を定めるとき、
命題\ref{prop:pg-represent}により、複素点の集合として
\[
 \operatorname{Pois}(X_F)(\mathbb C)=
 \{P\in\Lambda^3V:\,[\iota_{dF}P,\iota_{dF}P]=0\}
\]
と同一視する（右辺には射影化の規約を適用する）。
Jacobi方程式は$P$について二次斉次である。
$F_N=\sum_{i=1}^N z_i^3$、$X_N=V(F_N)$について
$\mathfrak P_N^{\mathrm{Fer}}=\operatorname{Pois}(X_N)$と書く。

また滑らかな三次式の射影パラメータ開集合を
$U_N\subset\mathbb P(S_3)$とし、その上の相対Poissonスキームの非零部分を
$\mathfrak I_N^\times$とする。$[F]$の局所表示を選ぶと、
双ベクトルは$\pi=\iota_{dF}P$と書ける。
$F$のスカラー変更には$P$を逆スカラー倍して対応するので、
相対の定義は局所表示の選択に依存しない。
相対モジュライは商スタック
$\mathfrak M_n=[\mathfrak I_N^\times/\mathrm{PGL}(V)]$の基礎位相空間として扱う。

$\left\{\begin{smallmatrix}N\\k\end{smallmatrix}\right\}_{\ge3}$を、
$\{1,\ldots,N\}$を各部分集合の要素数が$3$以上である
$k$個の\emph{非順序}非空部分集合へ分割する方法の数とする。

\begin{theorem}[Fermat fiberと相対schemeの成分数]
\label{thm:pg-fermat-relative-count}
$n\ge3$、$N=n+2$とする。定理\ref{thm:pg-normal}の下で、
\begin{enumerate}
\item Fermat cubic $X_N$のアフィンPoissonスキームの被約既約成分は
\[
 \#\operatorname{Irr}(\mathfrak P_N^{\mathrm{Fer}})
 =\sum_{k=1}^{\lfloor N/3\rfloor}
  \left\{\begin{smallmatrix}N\\k\end{smallmatrix}\right\}_{\ge3}.
 \tag{F}\label{eq:pg-fermat-count}
\]
 一般Poisson階数$2k$の成分の個数は
 $\left\{\begin{smallmatrix}N\\k\end{smallmatrix}\right\}_{\ge3}$で、
 各成分のアフィン次元は$3N-8k$である。
\item cubic自身も動かす相対スキームでは、一般階数$2k$ごとに
 ただ一つの既約成分が存在し、従って
\[
 \#\operatorname{Irr}(\mathfrak M_n)=\left\lfloor\frac{N}{3}\right\rfloor.
 \tag{R}\label{eq:pg-relative-count}
\]
 その商スタック次元は$m=N-3k$として
\begin{equation}
 \label{eq:pg-relative-dim}
 d_{n,k}=k(3N-8)+\binom{m+2}{3}
 +k\left(10+6m+3\binom{m+1}{2}\right)-N^2.
\end{equation}
\end{enumerate}
$\pi=0$を追加しても双方の被約既約成分数は変わらない。
\end{theorem}

\subsection{固定Fermat fiberの成分を座標分割で記述する}
\begin{proof}[定理\ref{thm:pg-fermat-relative-count}の第一項]
まず$P\in\Lambda^3V$が$F_N$に関してPoissonであるとする。
FermatのHessianは
$\operatorname{Hess}_{F_N}(z)=6\operatorname{diag}(z_1,\ldots,z_N)$である。
補題\ref{lem:pg-linearize}のJacobi条件を対角rank一二次形式に適用すると
\begin{equation}
 \label{eq:pg-fermat-contractions}
 (\iota_{dz_i}P)\wedge(\iota_{dz_i}P)=0
 \quad(1\le i\le N)
\end{equation}
を得る（これは各縮約二形式のPl\"ucker方程式そのものである）。

定理\ref{thm:pg-normal}により、$P=\sum_{s=1}^k P_s$、
$P_s\in\Lambda^3 W_s\setminus\{0\}$、$\dim W_s=3$、
$W_s$は互いに独立と書ける。
各$\iota_{dz_i}P_s$は$\Lambda^2W_s$に属する二形式である。
もし同じ$i$に対して二つの異なる$s,t$の縮約が非零ならば、
それらの外積は$\Lambda^2W_s\otimes\Lambda^2W_t$の非零元であり、
他の組の寄与と相殺できないから\eqref{eq:pg-fermat-contractions}に反する。
従って、各座標線形形式$dz_i$は高々一つの$W_s$上でしか非零にならない。
$B_s=\{i:dz_i|_{W_s}\ne0\}$とすれば、$B_s$は互いに素で、
$|B_s|\ge3$かつ$W_s\subset V_{B_s}:=\langle e_i:i\in B_s\rangle$である。
未使用の座標添字をいずれかの$B_s$へ加えれば、
$\{1,\ldots,N\}$の大きさ$3$以上の$k$個のブロックへの分割$\mathcal B$を得る。

逆に、そのような分割$\mathcal B=\{B_1,\ldots,B_k\}$に対して
\[
 Z_{\mathcal B}=\bigoplus_{s=1}^k
 \widehat{\operatorname{Gr}}(3,V_{B_s})
 \ \subset\ \Lambda^3V
\]
を考える。ハットはPl\"ucker埋込のアフィン錐（零点も含む）を表す。
これは線形独立な$\Lambda^3V_{B_s}$の中の錐の直積であり、閉かつ既約である。
各$P_s$は分解可能で、$F_N=\sum_s\sum_{i\in B_s}z_i^3$は
それぞれ別の座標に依存するから、各Jacobian Poisson括弧はJacobi恒等式を満たし、
異なるブロック間の括弧も零である。従って
$Z_{\mathcal B}\subset\operatorname{Pois}(X_N)$である。
先の議論からこれらの$Z_{\mathcal B}$は被約集合全体を被覆する。

$|B_s|=b_s$とすると
$\dim\widehat{\operatorname{Gr}}(3,V_{B_s})=3(b_s-3)+1$なので
\[
 \dim Z_{\mathcal B}=\sum_s(3b_s-8)=3N-8k.
\]
異なる$k$については、縮約写像
$\Lambda^2V^*\to V$のrankが一般点で$3k$なので、
小さい$k$の成分の閉包は大きい$k$の一般点を含めない。
逆向きの包含も上の次元が$k$とともに厳密に減少するので不可能である。
同じ$k\ge2$の異なる分割については、各$W_s$がそのブロックの
全座標方向に非零射影を持つ一般点を取る。
このとき縮約二形式の平方零集合から$W_s$の組が順序を除いて
内在的に復元される（補題\ref{ne:blocks}）。
したがって各$W_s$の座標支持$B_s$も復元され、別の分割による
$Z_{\mathcal B'}$へその一般点は入らない。
$k=1$の分割はただ一つである。
ゆえに全$Z_{\mathcal B}$が相互に異なる既約成分で、
式\eqref{eq:pg-fermat-count}が得られる。
零点は全ての錐に含まれ、追加の成分にはならない。
\end{proof}

\subsection{動くcubicのrank層は一つずつになる}
\begin{proof}[定理\ref{thm:pg-fermat-relative-count}の第二項]
$P_k=e_{123}+e_{456}+\cdots+e_{3k-2,3k-1,3k}$と置く。
定理\ref{thm:pg-normal}によって、rank$\,2k$の相対Poisson対の三重ベクトル部分は
$P_k$の$\mathrm{GL}_N$軌道$\mathcal O_{N,k}$である。
この軌道は既約で、その次元は$k(3N-8)$である。
実際$N=3k+m$とすると、$P_k$の安定化子のLie次元は
$8k+3km+m^2$なので、$N^2-(8k+3km+m^2)=k(3N-8)$となる。
$P_k$を固定した三次式は
\[
 F=c(u)+\sum_{s=1}^k\bigl(f_{s,3}(x_s)+f_{s,2}(x_s,u)+f_{s,1}(x_s,u)\bigr)
\]
の形の線形空間であり、係数次元は
\[
 \ell_{N,k}=\binom{m+2}{3}
 +k\left(10+6m+3\binom{m+1}{2}\right).
\]
その射影化の滑らかなcubicの部分は非空既約開集合である
（Fermat三次式が例になる）。よってrank$\,2k$層は既約で
次元$k(3N-8)+\ell_{N,k}-1$を持ち、
$\dim\mathrm{PGL}_N=N^2-1$を引くと式\eqref{eq:pg-relative-dim}を得る。

異なるrank層の閉包が互いに包含されないことを確かめる。
小さい$k$の層の閉包では縮約写像のrankが高々$3k$という閉条件が残る。
一方、隣り合う層について直接計算すると
\[
 \bigl(k(3N-8)+\ell_{N,k}-1\bigr)
 -\bigl((k+1)(3N-8)+\ell_{N,k+1}-1\bigr)
 =3(3k-1)(N-3k)+8>0.
\]
従って大きい$k$の層の閉包が小さい$k$の層全体を含むこともない。
以上によりrankごとに一つずつの既約成分となる。
零切断はrank二の相対錐の閉包に含まれる。
\end{proof}

\subsection{数値の比較と意味}
分割数の指数母関数は
\begin{equation}
\label{eq:pg-fermat-egf}
 \sum_{N\ge0}\frac{t^N}{N!}
 \left\{\begin{smallmatrix}N\\k\end{smallmatrix}\right\}_{\ge3}
 =\frac1{k!}(e^t-1-t-t^2/2)^k.
\end{equation}
従って係数抽出だけで固定Fermatの各rank成分数を求められる。
\begin{center}
\begin{tabular}{r|rrrr|r|r}
\toprule
$n$ & Fermat rank 2 & rank 4 & rank 6 & rank 8
& Fermat計 & 相対計 \\
\midrule
3 & 1 & -- & -- & -- & 1 & 1 \\
4 & 1 & 10 & -- & -- & 11 & 2 \\
5 & 1 & 35 & -- & -- & 36 & 2 \\
6 & 1 & 91 & -- & -- & 92 & 2 \\
7 & 1 & 210 & 280 & -- & 491 & 3 \\
8 & 1 & 456 & 2100 & -- & 2557 & 3 \\
9 & 1 & 957 & 10395 & -- & 11353 & 3 \\
10 & 1 & 1969 & 42735 & 15400 & 60105 & 4 \\
\bottomrule
\end{tabular}
\end{center}
例えば$n=7$（$N=9$）では、rank四の$210$は
$(3,6)$型の分割$\binom93=84$と$(4,5)$型の分割$\binom94=126$の和、
rank六の$280$は$(3,3,3)$型の分割
$9!/((3!)^3\,3!)$であり、$1+210+280=491$となる。
一方、相対的にはこの三つのrankにつき一つずつの三成分だけが残る。
これは\emph{固定特殊ファイバーの既約成分が、全体の異なる既約成分と
一対一に対応する必要がない}ことを具体的に示している。

\begin{remark}[最大性・scheme構造・既知結果との区別]
固定Fermatの成分数が、全ての滑らかなcubicの中で最大になることは
上の議論からは従わない。特に高次元での最大性は別の問題である。
表は\emph{被約既約成分数}についてのもので、Jacobi方程式で定義される
Poisson schemeが被約であるとか、成分の交叉や零点での局所環を
完全に決定したものではない。
また$\mathfrak M_n$は商スタックである。
その粗モジュライの存在・構成や全ての幾何学的性質をここで主張しない。
式\eqref{eq:pg-fermat-count}--\eqref{eq:pg-relative-dim}の初出・新規性は
包括的には未確認であり、正規形の既知入力と本節の計数的帰結を区別する。
\end{remark}

\section{固定した滑らかなcubicの非楕円型Poisson成分}\label{ne:examples}
座標$(x,y,z,a,b,c,u)$に対応するベクトル基底を$(e_0,\ldots,e_6)$とする。
\begin{equation}\label{ne:F}
\begin{aligned}
 f_{\mathrm N}&=xyz+y^3+z^3,& f_{\mathrm C}&=xz^2+y^3,\\
 F_\epsilon&=f_\epsilon+a^3+b^3+c^3+u(x^2+a^2+b^2+c^2)+u^3,
 &X_\epsilon&=V(F_\epsilon)\subset\PP^6
 \quad(\epsilon\in\{\mathrm N,\mathrm C\}).
\end{aligned}
\end{equation}
$P_1=e_0\wedge e_1\wedge e_2$、$P_2=e_3\wedge e_4\wedge e_5$、
$L=\CC P_1\oplus\CC P_2$とおく。

\begin{theorem}[非楕円型を一般点に持つ既約成分]\label{ne:main}
$X_{\mathrm N}$、$X_{\mathrm C}$は滑らかな三次五次元多様体である。
$P\mapsto\pi_P=(\iota_{dF_\epsilon}P)|_{X_\epsilon}$の下で
$L\simeq\mathbb A^2$は$\operatorname{Pois}(X_\epsilon)$の一つの既約成分の被約構造であり、
$P_0=P_1+P_2$でPoissonスキームそのものは滑らかである。
各成分の一般Poisson階数は四である。

$\lambda\mu\ne0$の点$\lambda P_1+\mu P_2$の三変数ブロック分解は順序を除き内在的である。
一方のブロックの平面三次曲線は、$X_{\mathrm N}$では節点型、$X_{\mathrm C}$では尖点型であり、
他方は滑らかなFermat三次曲線である。
特異なブロックを座標変換で滑らかな楕円ブロックへ取り替えることはできない。
\end{theorem}
一般点での被約性も主張に含むが、軸$\lambda\mu=0$や原点での他成分との交わりは別問題である。

\subsection{全体のcubicの滑らかさ}
\begin{lemma}\label{ne:smooth}
\eqref{ne:F}の全偏微分の共通零点は原点だけである。
\end{lemma}
\begin{proof}
$u=0$では$F_a=F_b=F_c=0$より$a=b=c=0$。
節点型の勾配$(yz,xz+3y^2,xy+3z^2)$が零なら$y=z=0$であり、
尖点型の勾配$(z^2,3y^2,2xz)$が零でも同様である。
残る$F_u=x^2$から$x=0$となる。

$u\ne0$では$u=1$に規格化する。
$a,b,c$はそれぞれ$0$または$-2/3$であり、その非零個数を$r$とすると
$a^2+b^2+c^2=4r/9$（$0\le r\le3$）。
節点型の残る三式は
\[
 yz+2x=0,\qquad xz+3y^2=0,\qquad xy+3z^2=0.
\]
$yz=0$なら$x=y=z=0$。
$yz\ne0$なら後二式を掛けて$x^2=9yz=-18x$、従って$x=-18$。
よって$F_u=0$は、$3+4r/9=0$または$327+4r/9=0$を要求し、いずれも不可能である。
尖点型では$z^2+2x=0$、$3y^2=0$、$2xz=0$から$x=y=z=0$で、
$F_u=3+4r/9\ne0$となる。
\end{proof}

\subsection{全Poisson方程式の接空間}
命題\ref{prop:pg-represent}により全双ベクトル空間は$\Lambda^3\CC^7$、次元35である。
Jacobi延長とgradientの一次syzygy消滅の議論は固定した$F_\epsilon$に対する線形な単射・完全列を使うので、
双対数環へのテンソル後にも成立し、以下は全Poissonスキームの接空間を与える。
$Q=\sum_{i<j<k}q_{ijk}e_i\wedge e_j\wedge e_k$とする。
Jacobi条件を$P_0$で線形化した整数行列$M_\epsilon$は
\begin{equation}\label{ne:M}
\begin{split}
 (M_\epsilon Q)_{ijk\alpha\ell}=\sum_{r,b}\bigl(&Q^{ir\alpha}P_0^{jkb}+P_0^{ir\alpha}Q^{jkb}
  +Q^{jr\alpha}P_0^{kib}+P_0^{jr\alpha}Q^{kib}\\
 &+Q^{kr\alpha}P_0^{ijb}+P_0^{kr\alpha}Q^{ijb}\bigr)(F_\epsilon)_{br\ell}
\end{split}
\end{equation}
で与えられる。$i<j<k$、$\alpha,\ell\in\{0,\ldots,6\}$であり、
行$(i,j,k,\alpha,\ell)$と列$(i,j,k)$を辞書式順で並べる。
零行を除くと節点型408行、尖点型312行となる。

\begin{lemma}[接空間の整数証明書]\label{ne:rank}
両行列の階数は33、核は$L$である。
列$(012)$と$(345)$を除き、JSON証明書に指定した33行（独立組版では付録にも掲載）を取ると
\[
 \det M_{\mathrm N}'=2^{24}3^{12}=8916100448256,\qquad
 \det M_{\mathrm C}'=-2^{33}3^{12}=-4565043429507072.
\]
\end{lemma}
\begin{proof}
$\lambda P_1+\mu P_2$は、中心変数$u$を共有する二つの三変数Jacobian括弧を与える。
ブロック間の括弧は零で、各ブロックの係数は他方の変数に依存しない。
よってJacobi恒等式が成り立ち、$(012)$、$(345)$の二列は零。階数は高々33である。

\eqref{ne:M}と指定した添字による整数小行列について、付属コードは整数行列式を計算して
上の二つの非零値を得る。JSONには選択行・列だけでなく小行列の全成分を記録した。
従って階数は33である。
これは有限体での標本検査でなく、具体的な$F_\epsilon$の全接方程式に対する
有理数体上の厳密計算を証明書とした論証である。
さらにHessian表示を介さず$[\pi,d\pi]$の三次係数を直接生成した2940行の行列でも
階数33を検算している。
\end{proof}

\begin{proof}[定理\ref{ne:main}の成分性と階数]
$L$は二次元の閉既約部分スキームであり、その全点がPoissonである。
$P_0$での接空間は二次元なので、Poissonスキームの局所環はKrull次元二、
埋込み次元二の正則局所環となる。
同じ次元を持つ閉部分$L$はその唯一の局所既約成分であり、局所的にスキームとして一致する。
従って$L$は大域的な既約成分で、一般点で被約である。

点$[0:1:0:-1:0:0:0]$は両$X_\epsilon$に属する。
$\lambda\mu\ne0$のとき、そこでアフィン双ベクトルの像は
$\langle e_0,e_2,e_4,e_5\rangle$で階数四。
Eulerベクトル$e_1-e_3$はその像に含まれないため、射影化しても階数四である。
ブロックの内在性は次に示す。
\end{proof}

\subsection{ブロックの内在性}
\begin{lemma}[縮約二形式による復元]\label{ne:blocks}
$W=\bigoplus_{s=1}^k W_s$、$\dim W_s=3$、
$P=\sum_s P_s$、$0\ne P_s\in\Lambda^3 W_s$とする。
$W^*$内のスキーム
\[
 D_P=\{\alpha:(\iota_\alpha P)\wedge(\iota_\alpha P)=0\}
\]
は$k$個の三次元線形部分空間$W_s^*$の被約な和である。
従って$P$は順序を除き$\{W_1,\ldots,W_k\}$を決定する。
\end{lemma}
\begin{proof}
$\alpha=(\alpha_s)$とし、$\beta_s=\iota_{\alpha_s}P_s$と書く。
三次元なので$\beta_s^2=0$、異なる$s,t$の$\beta_s\wedge\beta_t$は
互いに独立な$\Lambda^2W_s\otimes\Lambda^2W_t$成分に属する。
$W_s^*\to\Lambda^2W_s$は同型なので、$D_P$のイデアルは異なるブロックの座標の積全てで生成される。
これは各座標部分空間のイデアルの交叉、従って根基的である。
その既約成分を復元すれば、他成分の和の零化空間として各$W_s$も復元できる。
\end{proof}
\begin{proof}[定理\ref{ne:main}の残り]
$\lambda\mu\ne0$ではsupportは$\langle e_0,e_1,e_2\rangle\oplus\langle e_3,e_4,e_5\rangle$。
補題\ref{ne:blocks}により二ブロックは一意である。
$F_\epsilon$の対応する射影平面への制限は$f_\epsilon$および$a^3+b^3+c^3$である。
$f_{\mathrm N}$は$[1:0:0]$近傍で$yz+y^3+z^3$となり節点、
$f_{\mathrm C}$は$z^2+y^3$となり尖点を持つ。
各式は$x$について一次で、その係数と定数項が互いに素なので既約である。
この特異性は内在的ブロックへの制限であり、座標の選択では除去できない。
\end{proof}

\subsection{同次量子ブロックの行列式曲線}
$u=0$の三変数ブロックを、$[y,z]=h\operatorname{sym}(f_x)$とその巡回置換で量子化する。
多重線形化した二次関係式の行列は
\[
 M_{\mathrm N}(h)=\begin{pmatrix}
 0&-(1+h/2)z&(1-h/2)y\\
 (1-h/2)z&-3hy&-(1+h/2)x\\
 -(1+h/2)y&(1-h/2)x&-3hz
 \end{pmatrix},
\]
\[
 M_{\mathrm C}(h)=\begin{pmatrix}0&-z&y-hz\\z&-3hy&-x\\-y-hz&x&-hx\end{pmatrix}.
\]
\begin{proposition}\label{ne:quantum}
\begin{align*}
 \det M_{\mathrm N}(h)&=-\frac h4\bigl((h^2+12)xyz+(12-3h^2)(y^3+z^3)\bigr),\\
 \det M_{\mathrm C}(h)&=-3h\bigl(xz^2+y^3-h^2yz^2\bigr).
\end{align*}
前者は$h\ne0$、$h^2\notin\{4,-12\}$で節点三次曲線、
後者は全ての$h\ne0$で尖点三次曲線である。
\end{proposition}
\begin{proof}
行列式を展開すると二公式を得る。
節点型は二係数が非零のとき$x$の定数倍で$f_{\mathrm N}$に移り、
尖点型は$x'=x-h^2y$により$x'z^2+y^3$となる。
\end{proof}
ここでの主張は多重線形化した関係式の\emph{行列式曲線}についてである。
全point moduleのmoduliとの一致、AS-regular性、全量子化の相対smoothnessは
この計算だけからは主張しない。
$h=0$で行列式は恒等的に零となり、$h$を除いた極限曲線と
可換三変数多項式環のpoint scheme $\PP^2$を区別する。

\subsection{全成分復元における役割}
この二成分は可換復元の反例ではない。
特異ブロックの点$[1:0:0]$では全体の$u$方向微分$F_u=x^2$が非零であり、
中心変数との混合が全体の滑らかさを保つ。
従って、各ブロックの楕円性を必要条件とする方法ではこれらの成分を扱えない。
全ポテンシャルの対角、中心方向の微分、その量子化との両立が課題となる。

特異曲線の一般化Jacobianにも注意が必要である。
節点三次曲線では$\Pic^0\simeq\mathbb G_m$、尖点型では$\Pic^0\simeq\mathbb G_a$。
前者は正規化$\PP^1$上の二点のfiberを非零スカラーで接着することで、
後者は$\CC[[t^2,t^3]]\subset\CC[[t]]$の単元商の一次係数$t$により得られる。
従って、滑らかな楕円Picardに対して得た有理写像の障害を、これらにそのまま適用してはならない。
一般化Jacobianの導入だけで圏比較が成立するとも主張しない。

\begin{remark}[固定底と動く底]
本定理の成分は固定した$X_\epsilon$の$\operatorname{Pois}(X_\epsilon)$の成分である。
cubic自体も動かす相対Poissonスキームでは、同じ点がより大きな成分の特殊点となる可能性がある。
固定cubic上の全成分と、全cubicを動かす相対成分のgeneric性は区別する。
\end{remark}
\input{papers/sections/quantization_foundations}
\input{papers/sections/fermat_mixed_universal_center}
\input{papers/sections/hesse_mixed_universal_center}
\input{papers/sections/universal_casimir_determinant}
\input{papers/sections/algebraic_quantization}
\input{papers/sections/general_dimension_universal_quantization}
\input{papers/sections/equivariant_central_lifts}
\input{papers/sections/intrinsic_rank_two_center}
\input{papers/sections/full_support_quantization}
\input{papers/sections/special_poisson_blocks}
\section{非楕円成分の一様な量子化と周囲の相対対角}
\label{nq:sec:quantization}
本節では前節の固定した二つのPoisson成分を、実在する二パラメータの代数族として量子化する。
形式級数の存在や有限次数のHilbert関数検査だけにとどめず、全次数のflatnessと
\emph{周囲の代数}の有限相対対角分解まで構成する。
残余dg圏の対角は別の対象であり、本節の結論をそのsaturationと読み替えない。

\subsection{一次資料との関係と主定理}
三変数の同次Jacobian代数のAS-regular性はMori--Smith
\cite[Theorem 4.2, Lemma 4.1]{MSNE12}の分類を用いる。
中央一次元による拡大にはLe Bruyn--Smith--Van den Bergh
\cite[Theorem 2.6, Corollary 2.7, Theorem 3.1.3]{LBSVNE12}、
低次ポテンシャルを加えるPBW変形とその双加群複体にはBerger--Taillefer
\cite[Theorems 3.1 and 3.6]{BTNE12}を用いる。
これらの一般定理を新規な成果とは主張しない。
本節では、前節の特定の節点・尖点成分に対して中心混合項を保った具体的な式を与え、
係数環、例外集合、底変更および対角の種類を明記する。
中心三次元の表示の初出は未確定である。

$r,s$を次数零、$x,y,z,a,b,c,u$を次数一とし、
\[
 R_{\mathrm C}=\CC[r,s],\qquad
 R_{\mathrm N}=\CC[r,s,(4-r^2)^{-1}]
\]
と置く。第一ブロックの代数$D_\epsilon$は$u$を中心として
\begin{align}
 \epsilon=\mathrm C:\quad
 [y,z]&=r(z^2+2ux),&[z,x]&=3ry^2,&[x,y]&=r(xz+zx),
 \label{nq:eq:cusp}\\
 \epsilon=\mathrm N:\quad
 [y,z]&=\tfrac r2(yz+zy)+2rux,&
 [z,x]&=\tfrac r2(zx+xz)+3ry^2,&
 [x,y]&=\tfrac r2(xy+yx)+3rz^2
 \label{nq:eq:node}
\end{align}
で定める。第二ブロック$E_s$も$u$を中心とし、
\begin{equation}\label{nq:eq:fermat}
 [a,b]=s(3c^2+2uc),\quad
 [b,c]=s(3a^2+2ua),\quad
 [c,a]=s(3b^2+2ub)
\end{equation}
で定める。異なるブロックの生成元は可換とする。したがって
\[
 \mathcal A_\epsilon=D_\epsilon\otimes_{R_\epsilon[u]}E_s
\]
は七生成元の次数付き代数である。

\begin{theorem}[二つの非楕円成分の代数的量子化]
\label{nq:thm:family}
以下の式で定める$\Omega_\epsilon\in\mathcal A_{\epsilon,3}$は中心的であり、
$\mathcal B_\epsilon=\mathcal A_\epsilon/(\Omega_\epsilon)$は次を満たす。
\begin{enumerate}
\item $\mathcal A_\epsilon$は次数ごとに有限局所自由で、その各幾何学的ファイバーは
noetherianなKoszul AS-regular domain、global dimension七、AS parameter七である。
\item $\Omega_\epsilon$は各幾何学的ファイバーで非零の中心正則元であり、
$\mathcal B_\epsilon$は次数ごとに有限局所自由で、形成は任意の底変更と両立する。
各ファイバーのHilbert級数とAS-Gorenstein型は
\[
 H_{\mathcal B}(t)=\frac{1-t^3}{(1-t)^7},\qquad
 \operatorname{injdim}\mathcal B=6,\qquad \text{AS parameter}=4.
\]
\item $(r,s)=(0,0)$で$\mathcal B_\epsilon$は前節の$X_\epsilon$の斉次座標環となる。
任意の$(\lambda,\mu)$に沿う$r=h\lambda,s=h\mu$への引戻しは、
$\lambda\pi_{P_1}+\mu\pi_{P_2}$を半古典極限に持つ。
\item $R_\epsilon[u]$上の$\mathcal A_\epsilon$の対角は、階数列
$1,6,15,20,15,6,1$の有限自由双加群分解を持つ。
この分解は任意の可換$R_\epsilon[u]$-代数への底変更後も完全である。
$R_\epsilon$上でも周囲の対角はperfectである。
\end{enumerate}
この定理は$\operatorname{qgr}\mathcal B_\epsilon$やその残余圏のsmoothnessを含まない。
\end{theorem}

\subsection{第一ブロックのmonic PBW表示}
この小節では$r=h$と書き、係数環を
$Q_{\mathrm C}=\CC[h]$、$Q_{\mathrm N}=\CC[h,(4-h^2)^{-1}]$とする。
$u$を係数に入れ、active語の長さ、次いで辞書式$x>y>z$の順序を用いる。
$\epsilon=\mathrm C$の関係式は
\begin{align}
 yz&\longmapsto zy+hz^2+2hux,\notag\\
 xz&\longmapsto zx-3hy^2,\notag\\
 xy&\longmapsto yx+2hzx-3h^2y^2
 \label{nq:eq:cusp-rewrite}
\end{align}
と等価である。節点型では$q=(2+h)/(2-h)$として
\begin{align}
 yz&\longmapsto qzy+\frac{4h}{2-h}ux,\notag\\
 xz&\longmapsto q^{-1}zx-\frac{6h}{2+h}y^2,\notag\\
 xy&\longmapsto qyx+\frac{6h}{2-h}z^2.
 \label{nq:eq:node-rewrite}
\end{align}
全ての右辺のactive語は左辺より小さい。重なりの曖昧さは$xyz$の一つだけである。
\begin{lemma}[全次数と任意の底変更に対するPBW性]
\label{nq:lem:pbw}
二つの還元系はいずれも合流的であり、
$\{z^iy^jx^k:i,j,k\ge0\}$は$D_\epsilon$の$Q_\epsilon[u]$-基底である。
この主張は任意の可換$Q_\epsilon[u]$-代数への底変更後も成立する。
\end{lemma}
\begin{proof}
節点型で二つの還元経路から得る共通の正規形は
\[
 qzyx+\frac{4h}{2-h}ux^2+\frac{6h}{2-h}(z^3-y^3).
\]
尖点型では両経路が
\begin{align*}
 zyx+3hz^2x+2hux^2-3hy^3-9h^2zy^2-6h^3z^2y
 -6h^4z^3\\
 {}-12h^3uyx-24h^4uzx+36h^5uy^2
\end{align*}
となる。よってDiamond lemma \cite[Theorem 1.2]{BergNE12}を適用できる。
係数は表示した環に属し、還元の先頭係数は全て一である。
同じ有限恒等式が任意の底変更後にも成り立つため、基底の主張も保存される。
\end{proof}

\begin{lemma}[幾何学的ファイバーの正則性]\label{nq:lem:regular4}
$D_\epsilon$の各幾何学的ファイバーは四次元noetherian Koszul AS-regular domainである。
\end{lemma}
\begin{proof}
$u=0$の三変数代数は、尖点型では全$h$、節点型では$h\ne\pm2$で
三次元Calabi--Yau AS-regularである\cite[Lemma 4.1, Theorem 4.2]{MSNE12}。
域性とnoetherian性は、次のOre表示からも確認できる。
節点型では量子平面$yz=qzy$上に
\[
 \sigma(y)=qy,\quad\sigma(z)=q^{-1}z,\quad
 \delta(y)=\frac{6h}{2-h}z^2,\quad
 \delta(z)=-\frac{6h}{2+h}y^2
\]
による$x$のOre拡大を取る。
尖点型ではJordan平面$yz-zy=hz^2$上に
\[
 \sigma(y)=y+2hz,\quad\sigma(z)=z,\quad
 \delta(y)=-3h^2y^2,\quad\delta(z)=-3hy^2
\]
を使う。いずれも$\sigma$は可逆、$\delta$は表示した関係を保つ$\sigma$-導分である。
補題\ref{nq:lem:pbw}より$u$は中心正則元である。
中央次数一拡大の定理\cite[Theorem 2.6, Corollary 2.7 and Remark 2.8]{LBSVNE12}
により四次元の結論を得る。ここで特異なのは三次曲線であって、周囲のAS代数ではない。
\end{proof}

\subsection{混合項を含む中心三次元の閉公式}
全ての積を、表示した順序の非可換積として読む。
\begin{proposition}[第一ブロックの中心元]\label{nq:prop:center}
次の元は$D_\epsilon$の中心に属し、$h=0$で$f_\epsilon+ux^2$となる。
\begin{align}
 \Gamma_{\mathrm C}(h)={}&
 z^2x+y^3-3hzy^2+ux^2
 -6h^2uyx-6h^3uzx+12h^4uy^2,
 \label{nq:eq:centerC}\\
 \Gamma_{\mathrm N}(h)={}&
 z^3+\left(\frac{2-h}{2+h}\right)^3y^3
 +\frac{h^2+12}{3(2+h)^2}zyx
 +\frac{h^2+12}{6(2+h)}ux^2\notag\\
 &-\frac{24h^2}{(2+h)^2}uzy
 -\frac{24h^3}{(2+h)^2}u^2x.
 \label{nq:eq:centerN}
\end{align}
\end{proposition}
\begin{proof}
補題\ref{nq:lem:pbw}の還元を両辺の積$\Gamma_\epsilon a$、$a\Gamma_\epsilon$
（$a=x,y,z$）へ適用すると、全てで同じ正規形となる。
付属の\texttt{nonelliptic\_quantization\_certificate.json}に、六つの共通正規形を全係数で収録した。
その係数は尖点で$\mathbb Q[h,u]$、節点で$\mathbb Q[h,u,(4-h^2)^{-1}]$に属し、
コードは差の分子が恒等的に零であることを検査する。
従ってこれは有限個の特殊化による観察ではなく、合流する還元による普遍恒等式の証明である。
$u$との可換性は定義による。$h=0$での式は直接代入で分かる。
\end{proof}
\begin{remark}[量子的な混合補正の必要性]
$\Gamma_\epsilon(h)|_{u=0}+ux^2$という素朴な持上げは一般には中心でない。
付属コードはその非零な交換子の正規形も記録する。
中心混合項を古典式のまま加えるだけでは、量子化の中心三次方程式にならない。
\end{remark}

第二ブロックでは
$\alpha=a+u/3$、$\beta=b+u/3$、$\gamma=c+u/3$、$t=3s$と置く。
関係式\eqref{nq:eq:fermat}は
$[\alpha,\beta]=t(\gamma^2-u^2/9)$とその巡回置換になる。
\begin{lemma}[Fermat側の混合中心元]\label{nq:lem:fermat-center}
\[
 \Gamma_{\mathrm F}(s)=
 \alpha^3+\beta^3+\gamma^3
 -3s^2\!\sum_{w\in\mathfrak S_3(\alpha\beta\gamma)}w
 -\frac{u^2}{3}(\alpha+\beta+\gamma)+\frac{2u^3}{9}
 \tag{F}\label{nq:eq:centerF}
\]
は$E_s$の中心元であり、$s=0$で$a^3+b^3+c^3+u(a^2+b^2+c^2)$となる。
\end{lemma}
\begin{proof}
独立な中心係数$\eta$を用い、
$r_1=xy-yx-t(z^2-\eta)$、$r_2=yz-zy-t(x^2-\eta)$、
$r_3=zx-xz-t(y^2-\eta)$とする。
$D=x^3+y^3+z^3$、$T=xyz+xzy+yxz+yzx+zxy+zyx$に対し、
自由代数$\mathbb Z[t,\eta]\langle x,y,z\rangle$で
\begin{align*}
 [3D-t^2T-9\eta(x+y+z),x]={}&
 -3r_1y^2+t^2r_1zx-t^2xr_1z-3yr_1y-t^2zr_1x\\
 &+t^2xzr_1-3y^2r_1+2t^2r_2x^2-3tr_2yz\\
 &-t^2xr_2x-t^2x^2r_2+3tyzr_2
 +3r_3z^2+3zr_3z+3z^2r_3\\
 &+9\eta r_1-9\eta r_3
\end{align*}
が成立する。語ごとの展開と巡回置換により三交換子を処理できる。
$t=3s,\eta=u^2/9$を代入し、$x,y,z$を$\alpha,\beta,\gamma$に置き換える。
中心スカラー$2u^3/9$を加えると表示式を得る。
\end{proof}

\subsection{相対ポテンシャル複体の実際の完全性}
$Q=Q_\epsilon$、$S=Q[u]$とし、$D=D_\epsilon$を$S$-代数とみる。
自由$S$-代数のcyclic potentialを
\[
 W=xyz-yxz-h\operatorname{sym}(f_\epsilon+ux^2),
 \qquad r_i=\partial^{\mathrm{cyc}}_{x_i}W,\quad(x_1,x_2,x_3)=(x,y,z)
\]
とする。定義関係式は$r_1=r_2=r_3=0$である。

$\partial^{\mathrm{nc}}_j$は語の各$x_j$を取り除き、その左右の部分を分ける演算とする：
\[
 \partial^{\mathrm{nc}}_j(a_1\cdots a_d)
 =\sum_{a_\ell=x_j} a_1\cdots a_{\ell-1}\otimes a_{\ell+1}\cdots a_d.
\]
係数環$S$は微分しない。
\begin{proposition}[底変更と両立する第一ブロックの相対対角]
\label{nq:prop:relative-diagonal}
$D^e_S=D\otimes_S D^{\mathrm{op}}$とする。有限自由双加群の複体
\[
 0\longrightarrow D^e_S
 \xrightarrow{d_3}(D^e_S)^3
 \xrightarrow{d_2}(D^e_S)^3
 \xrightarrow{d_1}D^e_S
 \xrightarrow{m}D\longrightarrow0
 \tag{K}\label{nq:eq:potential-resolution}
\]
を、基底$e_i$（生成元）、$\varepsilon_i$（関係式）について
\begin{align*}
 d_1(e_i)&=x_i\otimes1-1\otimes x_i,\\
 d_2(\varepsilon_i)&=\sum_j\sum_w\ \sum_{w=p x_j q} c_{iw}p e_j q,\\
 d_3(1)&=\sum_i(x_i\varepsilon_i-\varepsilon_i x_i)
\end{align*}
で定める。ここで$r_i=\sum_w c_{iw}w$と展開している。
この複体は完全で、任意の可換$S$-代数への底変更後も完全である。
\end{proposition}
\begin{proof}
$d_1d_2=0$は$r_i=0$から、$d_2d_3=0$はcyclic potentialの恒等式
$\sum_i[x_i,r_i]=0$を非可換微分することから従う。
active次数に関して、複体の自由基底のfiltration次数を$0,1,2,3$と取る。
$ux^2$はactive次数二なので、随伴次数付き複体は$u=0$の同次ポテンシャルの
Koszul双加群分解を$Q[u]$へ拡大したものである。

この同次複体が$Q$上で完全であることを確認する。
補題\ref{nq:lem:pbw}の$u=0$で、各次数成分は有限自由$Q$-加群である。
各複素点に特殊化すると、Mori--Smithの三次元Calabi--Yau代数に対する
同次ポテンシャル分解となり完全である
\cite[Theorem 4.2]{MSNE12}\cite[\S3.2]{BTNE12}。
各次数における複体は有限個の有限自由$Q$-加群からなり、
全閉点のfiberで完全なので$Q$上で完全である。
さらに最後の項も自由であるから、各次数の完全列は$Q$-加群として分裂する。
従って任意の底変更後にも完全性が保たれる。

filtrationは下に有界で、各元のactive次数は有限である。
随伴複体でcycleの最高次数を消し、次数を下げる帰納法により
元の複体の完全性を得る。底変更後も同じmonic PBW表示と同じ随伴分解を持つため、
同じ証明が適用できる。これは既知のPBWポテンシャル分解
\cite[Theorem 3.6の証明]{BTNE12}を、今回の係数環と底変更について明示したものである。
\end{proof}

\begin{remark}[中心方向の通常微分は使えない]\label{nq:rem:no-derivation}
$h\ne0$のファイバーでは、$\partial_u(u)=1$を満たす次数$-1$の
代数導分$\partial_u:D_\epsilon\to D_\epsilon$は存在しない。
実際、次数から$\partial_u(x),\partial_u(y),\partial_u(z)$はスカラーとなる。
尖点型の$[y,z]-h(z^2+2ux)=0$を微分して$u=0$に落とすと、
$-2h(x+\partial_u(z)z)=0$という一次の関係が必要になるが、PBW性に反する。
節点型でも$-h(2x+\partial_u(y)z+\partial_u(z)y)=0$となって同様である。
相対対角\eqref{nq:eq:potential-resolution}は$u$を係数としており、
この不存在とは矛盾しない。中心三次商の残余圏を調べる際に、
可換多項式の$\partial_u$をそのまま作用させることはできない。
\end{remark}

\subsection{七生成元への組立てと中心三次商}
まず$E_s$を$\CC[s]$上で考える。$u=0$のFermat型三次元Jacobian代数は
全ての$s$で三次元Calabi--Yau AS-regularである
\cite[Theorem 4.2(1)]{MSNE12}。
$u$を加えた関係式の一次混合係数は対角行列であり、
cyclic potentialのsyzygyにより
\cite[Theorem 3.1.3]{LBSVNE12}の条件を$\gamma=0$で満たす。
従って各$s$で$E_s$は中心正則$u$を持つ四次元Koszul AS-regular domainである。
各次数の$\CC[s]$-加群は有限生成で全fiberの次元が一定なので自由となる。
$u$による完全列を次数ごとに分裂させると、$E_s$は$\CC[s,u]$上自由で、
active filtrationの随伴次数付き環は同次Fermat代数の係数拡大となる。
そのポテンシャル複体も命題\ref{nq:prop:relative-diagonal}と同じ議論で、
任意の底変更と両立する長さ三の相対自由対角分解を与える。

\begin{proof}[定理\ref{nq:thm:family}の証明]
中心三次元を
\[
 \Omega_\epsilon=\Gamma_\epsilon(r)+\Gamma_{\mathrm F}(s)+u^3
 \tag{Q}\label{nq:eq:total-center}
\]
と定める。中心性と原点への特殊化は既に示した。
両ブロックが$R_\epsilon[u]$上自由なので、そのテンソル積も自由である。
全体の$u=0$での商は、二つの三次元Koszul AS-regular代数のテンソル積であり、
Hilbert級数$(1-t)^{-6}$、AS型$(6,6)$を持つ。
このテンソル積がnoetherian domainであることには注意が要る。
一般のnoetherian代数のテンソル積についての無条件な主張は用いず、
第一ブロックの上記Ore表示を第二のFermat代数上へ係数拡大する。
これにより二変数量子平面またはJordan平面と、もう一つの可逆Ore拡大の反復として得られる。

$u$は全体でも中心正則元である。
中心正則次数一拡大に対するKoszul性の持上げ
\cite[Theorem 2.6, Remark 2.8]{LBSVNE12}と、
$H_{\mathcal A}(t)=(1-t)^{-7}$から、Koszul双対の級数は$(1+t)^7$となりglobal dimensionは七である。
AS-Gorenstein型は中心正則元に対するchange-of-ringsの
\[
 \operatorname{Ext}_{\mathcal A}^{i}(k,\mathcal A)
 \simeq\operatorname{Ext}_{\mathcal A/(u)}^{i-1}
             (k,\mathcal A/(u))(1)
\]
および左側の同じ式から$(7,7)$となる。
noetherian性と域性は$u$-進次数に関する随伴環
$(\mathcal A/(u))[u]$から従う。正次数の$u$を用いるので、固定した全次数における
最高項の消去は有限回で停止する。

各ファイバーで$\Omega_\epsilon\bmod u$の第一ブロック次数$(3,0)$の項は
$\Gamma_\epsilon(r)|_{u=0}$であり、そのPBW正規形は非零である。
第二ブロック次数$(0,3)$の項とは相殺できないため、$\Omega_\epsilon\ne0$である。
周囲がdomainなので左右正則である。
各次数で$\Omega_\epsilon:\mathcal A_{d-3}\to\mathcal A_d$は
全ての幾何学的fiberで単射だから、局所的に分裂した単射であり、
余核は有限局所自由で任意の底変更と両立する。
中心正則次数三商のchange-of-ringsからAS型$(6,4)$を得る。
半古典極限は関係式を$h$で割って$r=h\lambda,s=h\mu$の原点へ制限すればよい。

最後に、$R_\epsilon[u]$上で二つの長さ三の相対ポテンシャル分解をテンソルしtotalizeする。
その自由階数列は$(1+t)^3(1+t)^3=(1+t)^6$であり、完全性と任意の底変更が保たれる。
$R_\epsilon$上では包絡代数から相対包絡代数への商は中心元
$u\otimes1-1\otimes u$による商であり、この元は正則である。
その長さ一のKoszul分解と相対分解から、周囲の対角は絶対的にもperfect（射影次元高々七）となる。
\end{proof}

\subsection{相対残余圏について残る仕事}
本節は二つの固定Poisson成分の稠密開を同時に扱う量子化族を与える。
一つの方向だけを固定した形式弧ではなく、$(r,s)$の両方向を保持する。
周囲の代数の有限対角分解は実在し、その底変更も証明した。
しかし、中心三次商のaffine coneは可換点でも原点で特異である。
\emph{周囲の対角を中心三次商でテンソルするだけ}では、その商や残余圏の対角にはならない。
次に必要なのは、中心三次商の相対局所cohomology・射影化とOrlov windowを含むproperなdg模型、
およびその模型の対角のperfect性である。
この過程で中心方向の混合を除かず、注意\ref{nq:rem:no-derivation}の障害を扱わなければならない。
一般次元・全Poisson成分への可換復元は、引き続きこの未解決段階を含む。
\appendix
\section{Hesse型Poisson量子化の明示的な三次中心元}
\label{sec:explicit-hesse-center}
本節では量子化の定理が与える三次中心元を、滑らかなHesse型
三変数ブロックについて閉公式で指定する。
この公式はSklyanin代数の既知の中心三次式
\cite[Theorem 1およびSection 2]{DeLaetLeBruyn}
（原式は\cite[p.~35, Remark (1)]{SmithTate}）
のPoisson量子化パラメータへの代入である。
従って中心元の存在という既知の結果自体を新規な主張とはしない。
本研究で有用なのは、$0$を含む明示的なパラメータ開集合と
$\hbar=0$での留数を正規化した公式を同時に得ることである。

\subsection{公式と適用開集合}
$t\in\CC$、$t^3\ne-1$とし、
\[
 f_t(x,y,z)=\frac{x^3+y^3+z^3}{3}+txyz
\]
と置く。$\hbar$を次数零のパラメータとして、
\[
 [y,z]=\hbar\left(x^2+\frac{t}{2}(yz+zy)\right),
 \quad
 [z,x]=\hbar\left(y^2+\frac{t}{2}(zx+xz)\right),
 \quad
 [x,y]=\hbar\left(z^2+\frac{t}{2}(xy+yx)\right)
 \tag{H}\label{eq:hesse-commutators}
\]
による次数付き代数を$S_{\hbar,t}$と書く。
ここで
\[
 S=x^3+y^3+z^3,\quad
 C=xyz+yzx+zxy,\quad
 D=yxz+xzy+zyx,\quad T=(C+D)/6.
\]
次の$\hbar$の多項式を定める：
\[
 U(\hbar,t)=1+\frac{\hbar^2t^2}{2}
  +\frac{\hbar^4t}{6}-\frac{\hbar^4t^4}{48},
 \qquad
 W(\hbar,t)=
 \frac{(4-\hbar^2t^2)
 (12t+\hbar^2(t^3-8))}{48}.
 \tag{HU}\label{eq:hesse-UW}
\]
また
\[
 \Delta_{\hbar,t}=
 \left(t+\frac{\hbar^2}{3}
       \left(1+\frac{t^3}{4}\right)\right)^3
  +\left(1-\frac{\hbar^2t^2}{4}\right)^3.
 \tag{HD}\label{eq:hesse-discriminant}
\]
と置く。$\Delta_{0,t}=1+t^3\ne0$である。

\begin{theorem}[Hesse中心三次式の閉公式]
\label{thm:explicit-hesse-center}
$t^3\ne-1$に対し、
\[
 U_t=\{\hbar\in\mathbb A^1_{\CC}:
   U(\hbar,t)(1-\hbar^2t^2/4)\Delta_{\hbar,t}\ne0\}
\]
とする。$0\in U_t$であり、
各$\hbar\in U_t$に対して
$S_{\hbar,t}$は三生成元のKoszul AS-regular domainで、
Hilbert級数$(1-z)^{-3}$を持つ。
元
\[
 \boxed{\quad
 \Omega_{\hbar,t}=
   \frac{U(\hbar,t)}3(x^3+y^3+z^3)
   +W(\hbar,t)\operatorname{sym}(xyz)
 \quad}
 \tag{HC}\label{eq:hesse-central-element}
\]
は$S_{\hbar,t}$の次数三の中心正則元であり、
$\Omega_{0,t}=f_t$である。
特に$t=0$では
$\Delta_{\hbar,0}=1+\hbar^6/27$で、
\[
 \Omega_{\hbar,0}
 =\frac{x^3+y^3+z^3}{3}
   -\frac{2\hbar^2}{3}\operatorname{sym}(xyz)
\]
を得る。
\end{theorem}
\begin{proof}
$q=\hbar t/2$、
$a=1-q$、$b=-(1+q)$、$c=-\hbar$と置く。
\eqref{eq:hesse-commutators}は
\[
 ayz+bzy+cx^2=0,\quad
 azx+bxz+cy^2=0,\quad
 axy+byx+cz^2=0
\]
となる。$\hbar\ne0$かつ
$(1-\hbar^2t^2/4)\Delta_{\hbar,t}\ne0$なら、
$abc\ne0$で、標準Sklyaninの
\[
 (a^3+b^3+c^3)^3\ne(3abc)^3
\]
が成り立つ。
実際、左辺の差は
$-27\hbar^3\Delta_{\hbar,t}$である。
従ってこの開部分では点曲線が滑らかな楕円曲線であり、
$S_{\hbar,t}$は三次元Sklyanin AS-regular domainである
\cite{SmithTate,DeLaetLeBruyn}。
$\hbar=0$の代数は可換多項式環である。

既知の中心元の公式は
\[
 G=a(b^3-c^3)C+
   b(c^3-a^3)D+
   c(a^3-b^3)S
 \tag{G}\label{eq:hesse-sklyanin-G}
\]
である\cite[Section 2]{DeLaetLeBruyn}。
一方、三つの定義関係式にそれぞれ$x,y,z$を左から掛けて足すと
\[
 C-D=\hbar S+\frac{\hbar t}{2}(C+D)
 \tag{R}\label{eq:hesse-CD}
\]
を得る。この等式で$G$を$S,C+D$で表し、
$a,b,c$を代入して係数を整理すると
\[
 G=-9\hbar\left(
       \frac{U(\hbar,t)}3S+
       \frac{W(\hbar,t)}6(C+D)\right)
   =-9\hbar\Omega_{\hbar,t}.
 \tag{I}\label{eq:hesse-G-identity}
\]
これは$\hbar$で割る操作をせずに成立する多項式恒等式である。
$\hbar\ne0$の滑らかな楕円Sklyanin開部分では
$G$が中心なので$\Omega_{\hbar,t}$も中心である。
各次数成分が$\hbar$について有限自由であることは
三次元SklyaninのHilbert級数と零ファイバーにおける
同じ級数、および元論文のPBW構成から従う。
従って各交換子$[\Omega_{\hbar,t},x]$等は
有限自由加群の切断として稠密開集合上で零であり、
$U_t$の全体で零となる。
$\Omega_{0,t}=f_t$は$T_0=xyz$と
$U(0,t)=1$、$W(0,t)=t$から分かる。
既知のSklyanin中心元$G$は楕円開部分で非零であり、
代数はdomainだから正則である。
零ファイバーでも$f_t$は非零因子である。
\end{proof}


\begin{corollary}[任意の滑らかな三変数三次式]
\label{cor:arbitrary-smooth-ternary-cubic}
$g(x_1,x_2,x_3)$が滑らかな平面三次曲線を定める
斉次三次式であるとする。$g$のJacobian Poisson括弧
\[
 [x_2,x_3]=\hbar\operatorname{sym}(\partial_{x_1}g),
 \quad [x_3,x_1]=\hbar\operatorname{sym}(\partial_{x_2}g),
 \quad [x_1,x_2]=\hbar\operatorname{sym}(\partial_{x_3}g)
\]
から得られる量子化は、$0$を含むZariski開集合上で
三次中心元の閉公式を持つ。
実際、Hesse標準形\cite[Lemma 1]{ArtebaniDolgachev}により
ある$M\in\mathrm{GL}_3(\CC)$、$\lambda\in\CC^\times$、
$t^3\ne-1$が存在して
\[
 g(x)=\lambda f_t(Mx)
\]
となる。$y=Mx$、$\hbar'=\lambda\det(M)\hbar$とおけば、
\[
 \boxed{\ \Omega^g_\hbar(x)=
  \lambda\Omega_{\hbar',t}(y)\big|_{y=Mx}\ }
 \tag{HN}\label{eq:arbitrary-ternary-central}
\]
は三次中心元で、$\Omega^g_0=g$である。
\end{corollary}
\begin{proof}
Poissonの定数三重ベクトル
$\partial_{x_1}\wedge\partial_{x_2}\wedge\partial_{x_3}$
は座標変換$y=Mx$の下で
$\det(M)\,
 \partial_{y_1}\wedge\partial_{y_2}\wedge\partial_{y_3}$
に移る。従って$y$の交換子は
\[
 [y_2,y_3]
 =\lambda\det(M)\hbar\,
   \operatorname{sym}(\partial_{y_1}f_t),
\]
およびその巡回置換となる。
対称化は線形座標変換と可換なので、量子化代数は
$S_{\hbar',t}$への次数付き同型を持つ。
その同型によって定理\ref{thm:explicit-hesse-center}の
中心元を戻し、留数が$g$になるよう$\lambda$倍する。
\end{proof}

\begin{corollary}[中心変数と独立ブロックの閉公式]
\label{cor:hesse-block-center}
相異なる三変数ブロック$x_s=(x_s,y_s,z_s)$と
中心変数$u$について
\[
 F=\sum_{s=1}^k
 \left(\frac{x_s^3+y_s^3+z_s^3}{3}
       +t_sx_sy_sz_s\right)+P(u),
 \qquad t_s^3\ne-1,
\]
$P$は任意の斉次三次式とする。
各ブロックに\eqref{eq:hesse-commutators}を課し、
異なるブロック間と$u$を可換に取る。
各$U_{t_s}$の共通開部分で
\[
 \Omega_\hbar
   =\sum_s\Omega_{\hbar,t_s}+P(u)
 \tag{HB}\label{eq:hesse-multiblock}
\]
は次数三の中心元で、$\Omega_0=F$である。
各$u$の一次形式$\ell_{si}(u)$を選んで
$x_{si}$を$x_{si}+\ell_{si}(u)$に置き換えた
非自明な混合項を含むPoissonブロックでも、
同じ置換による明示公式が成立する。
\end{corollary}
\begin{proof}
各$S_{\hbar,t_s}$の中心元は他のブロックと可換であり、
$P(u)$も中心なので和は中心である。
中心線形シフトは生成元の次数付き同型を定め、
同型によって中心元が中心元へ移る。
\end{proof}


\subsection{可換復元の候補となる点曲線}
\begin{proposition}[楕円点曲線のHesseパラメータ]
\label{prop:hesse-point-cubic}
定理\ref{thm:explicit-hesse-center}の条件で$\hbar\ne0$とする。
$S_{\hbar,t}$のpoint schemeを定める滑らかな平面三次曲線は
\[
 E_{\hbar,t}:\quad
 x^3+y^3+z^3+3\,\vartheta(\hbar,t)xyz=0,
\quad
 \boxed{\quad
 \vartheta(\hbar,t)=
 \frac{t+\frac{\hbar^2}{3}(1+t^3/4)}
 {1-\hbar^2t^2/4}
 \quad}
 \tag{EP}\label{eq:hesse-point-parameter}
\]
である。$\vartheta(0,t)=t$と形式的に延長され、
\[
 \Delta_{\hbar,t}
 =\left(1-\frac{\hbar^2t^2}{4}\right)^3
  \left(1+\vartheta(\hbar,t)^3\right).
 \tag{ED}\label{eq:hesse-point-delta}
\]
特に$\hbar\in U_t$なら、右辺のHesse三次曲線は滑らかである。
\end{proposition}
\begin{proof}
滑らかな楕円型三次元Sklyanin代数のpoint schemeは
\[
 abc(x^3+y^3+z^3)-(a^3+b^3+c^3)xyz=0
\]
である\cite{SmithTate,DeLaetLeBruyn}。
定理の証明で与えた
$a=1-\hbar t/2$、$b=-(1+\hbar t/2)$、$c=-\hbar$
について
\[
 abc=\hbar(1-\hbar^2t^2/4),\qquad
 a^3+b^3+c^3
 =-3\hbar t-\hbar^3(1+t^3/4)
\]
を代入し、共通係数$\hbar$で割ると表示を得る。
最後の恒等式は\eqref{eq:hesse-discriminant}の定義から従う。
\end{proof}
\begin{remark}[幾何学的比較との関係]
この楕円点曲線と中心三次元から得られる非可換cubicの
rotation付き可換復元は論文Cで扱う。
具体的な一楕円ブロックの比較については
同論文の「一ブロックの明示的可換復元」を参照されたい。
ここでは中心三次量子化の存在・公式と圏同値を区別する。
\end{remark}

\begin{remark}[適用範囲と新規性]
\label{rem:hesse-scope}
一般の$f_s(x_s,u)$には$u$とactive変数の本質的な
混合項があり、必ずしも上の中心シフト形にはならない。
その場合の中心元は依然として形式的存在と
有限自由加群間の線形方程式による代数化で定まる。
また、\eqref{eq:hesse-central-element}は
周囲の量子化と中心元の明示式であり、
残余dg圏の相対saturationや
rotation付き可換復元を結論するものではない。
中心元の基本公式は古典的Sklyanin理論に由来し、
本節のHesseパラメータへの翻訳の新規性は未確認である。
\end{remark}

\input{papers/sections/ambient_recovery_general}
\section{射影完全交叉と斉次多様体のPoisson構造}
\label{sec:general-poisson-existence}
本節は三次超曲面のPoisson構造を一般の射影多様体の中に位置づける。
完全交叉に対する非存在の部分はBr\"uckmannの古典的消滅定理
（Fu--Hwang \cite[Lemma 4.1]{FuHwangCI} に引用）による。
本節の既知結果を新規な分類とみなさない。

\subsection{滑らかな射影完全交叉の存在判定}
\begin{theorem}\label{thm:ci-poisson-existence}
複素数体上の滑らかな通常の射影完全交叉
\[
 X=X_{d_1,\ldots,d_c}\subset\PP^{n+c},
 \qquad n\ge3,\quad d_j\ge2
\]
について、以下は同値である。
\begin{enumerate}[label=(\arabic*)]
\item $H^0(X,\bigwedge^2T_X)\ne0$。
\item $X$は非零の正則Poisson構造を持つ。
\item 多重次数は空列、$(2)$、$(3)$、$(2,2)$のいずれかである。
\end{enumerate}
空列は$\PP^n$を表す。一次式の方程式はあらかじめ除いている。
\end{theorem}
\begin{proof}
$s=\sum_{j=1}^c(d_j-1)$とおく。随伴公式から
\[
 \omega_X\simeq\OO_X(-(n+1-s)),\qquad
 \bigwedge^2T_X\simeq\Omega_X^{n-2}(n+1-s)
 \tag{2.1}\label{eq:ci-forms}
\]
である。Br\"uckmannの消滅\cite[Lemma 4.1]{FuHwangCI}は
\[
 H^0(X,\Omega_X^p(k))=0
 \qquad(1\le p\le n-1,\ k\le p)
 \tag{2.2}\label{eq:ci-vanishing}
\]
を与える。$n\ge3$なので$p=n-2\ge1$を取り、
$s\ge3$なら$n+1-s\le n-2$ゆえ
$H^0(X,\bigwedge^2T_X)=0$となる。
したがって(1)なら$s\le2$であり、$d_j\ge2$より列挙した四種類に限る。

逆に$s\le2$の場合、一般の独立な一次形式
$L_0,\ldots,L_{n-2}$を選び、射影空間のEuler場
$E=\sum_i z_i\partial_{z_i}$を用いて
\[
 \alpha=\iota_E(dL_0\wedge\cdots\wedge dL_{n-2})
       \in H^0(X,\Omega_X^{n-2}(n-1))
\]
を定める。一般の選択では$\alpha|_X$は非零である。
$L_0\ne0$上では
\[
 \alpha=L_0^{n-1}\,
 d(L_1/L_0)\wedge\cdots\wedge d(L_{n-2}/L_0)
\]
となり、一般点でその核は有理写像
$X\dashrightarrow\PP^{n-2}$のファイバーに接する
階数二の可積分分布である。
$s=2$なら\eqref{eq:ci-forms}によって$\alpha$を
二重ベクトルに対応させる。
$s=1$なら非零の一次切断、$s=0$なら非零の二次切断を掛け、
同じ同型によって対応させる。
いずれも一般階数二で非零であり、
可積分分布を定める分解可能二重ベクトルなので
Schouten条件$[\pi,\pi]=0$を満たす。
これで(3)から(2)、(2)から(1)が従う。
\end{proof}

\begin{remark}[曲線と曲面]\label{rem:ci-small-dimension}
曲線では$\bigwedge^2T_X=0$。
曲面では$\bigwedge^2T_X=\omega_X^{-1}$で
Schouten条件は自動的に成立する。
滑らかな完全交叉曲面の非零Poisson構造の存在条件は
$s\le3$であり、上記に加えて$(4)$、$(2,3)$、
$(2,2,2)$のK3完全交叉が現れる。
従って$n\ge3$は本質的である。
特異完全交叉には上の消滅をそのまま適用しない。
\end{remark}

\subsection{トーリック多様体}
\begin{proposition}\label{prop:toric-poisson}
次元二以上の複素正規トーリック多様体には
非零の代数的Poisson括弧がある。
滑らかなら$\pi\in H^0(X,\bigwedge^2T_X)$で与えられる。
\end{proposition}
\begin{proof}
稠密トーラス$T\simeq(\CC^*)^n$のLie代数の
一次独立な二元が定める大域導分を$D_1,D_2$とする。
$T$が可換なので$[D_1,D_2]=0$であり、
$\pi=D_1\wedge D_2$はJacobi条件を満たす。
正規トーリックの各不変アフィンチャートで
$\{f,g\}=D_1(f)D_2(g)-D_2(f)D_1(g)$
は正則で、貼り合わさる。
稠密トーラス上では二導分は独立なので非零である。
\end{proof}

\subsection{滑らかな射影斉次多様体}
\begin{theorem}\label{thm:homogeneous-poisson}
$\CC$上の滑らかな連結射影斉次多様体$X$で
$\dim X\ge2$ならば、$X$は非零の正則Poisson構造を持つ。
\end{theorem}
\begin{proof}
Borel--Remmertの定理\cite{BorelRemmert}により
$X\simeq A\times G/P$と分解できる。
$A$はアーベル多様体、$G/P$は半単純群の旗多様体の直積である。
$A$の並進ベクトル場は可換であり、
$G/P$の最大トーラスから来る基本ベクトル場も可換である。
後者で非自明な単純因子の階数が二以上ならば、
$P$に含まれない単純根$\alpha$と、
それに隣接する単純根$\beta$に対し、
$\alpha$および$\alpha+\beta$は接空間の
一次独立な二つのトーラスウェイトを与える。
従って最大トーラスの一般軌道は二次元以上である。
階数一の非自明因子は$\PP^1$であり、
二つ以上のそのような因子があれば同様に二方向がある。
残る場合、$\dim X\ge2$から$A$の並進方向を
必要な数だけ加えられる。
以上より$X$上に一般点で一次独立な二つの
可換な大域ベクトル場$D_1,D_2$が存在する。
その外積は非零で
$[D_1\wedge D_2,D_1\wedge D_2]=0$を満たす。
\end{proof}
\begin{remark}
非射影の任意の斉次空間に対して同じことを
示したわけではない。
次元一では二重ベクトルは常に零となる。
\end{remark}

\subsection{Poisson量子化の明示性}
\begin{proposition}\label{prop:poisson-quantization-explicit}
本論文のPoisson正規形
$F=\sum_{s=1}^k f_s(x_s,u)$に対し、
量子化を定める二次関係式は明示的に
\[
 [x_{s2},x_{s3}]
 =\hbar\,\operatorname{sym}_R(\partial_{x_{s1}}f_s),
 \quad
 [x_{s3},x_{s1}]
 =\hbar\,\operatorname{sym}_R(\partial_{x_{s2}}f_s),
 \quad
 [x_{s1},x_{s2}]
 =\hbar\,\operatorname{sym}_R(\partial_{x_{s3}}f_s)
\]
と書ける。異なるブロックと中心変数の間の交換子は零である。
三次中心元$\Omega_\hbar$は零のZariski近傍で存在し、
有限自由加群間の中心性方程式
$z\mapsto([z,x_{si}])_{s,i}$の核を解いて選べる。
しかし任意の入力に一様な閉公式が得られたとは主張しない。
\end{proposition}
\begin{proof}
二次関係式は各ブロックのcyclic potential
$W_s=x_{s1}x_{s2}x_{s3}-x_{s2}x_{s1}x_{s3}
-\hbar\,\operatorname{sym}_R(f_s)$のcyclic derivativeである。
中心元の存在と、核の計算を用いた代数化は
本論文の量子化節で証明する。
\end{proof}
\begin{example}[Fermat型]
$[x,y]=\hbar z^2$、$[y,z]=\hbar x^2$、
$[z,x]=\hbar y^2$なら
\[
 \Omega=x^3+y^3+z^3
      -2\hbar^2\operatorname{sym}(xyz)
\]
は実際に中心元となる。直接の交換子計算は
本論文のFermat型の補題にある。
\end{example}

% 2026-11-12: article A only. Do not include B/C-only references.
% Main manuscript: papers/poisson_cubic_quantization.tex
\begin{thebibliography}{99}
\bibitem{PHartshorne}
R.~Hartshorne,
\emph{Algebraic Geometry}, Graduate Texts in Mathematics 52,
Springer, 1977.
\href{https://doi.org/10.1007/978-1-4757-3849-0}{doi:10.1007/978-1-4757-3849-0}.

\bibitem{PEtingofLie}
P.~Etingof,
\emph{Lie Groups and Lie Algebras}, MIT lecture notes, 2024 course version.
\href{https://live.ocw.mit.edu/courses/18-755-lie-groups-and-lie-algebras-ii-spring-2024/mit18_755_s24_lec_full.pdf}{MIT OpenCourseWare 公開原稿}.
本文では Proposition~16.9、\S19.4、Theorem~25.17、\S32.2 を参照する。

\bibitem{PMayanskiy}
E.~Mayanskiy,
Poisson cohomology of two Fano threefolds,
\emph{Journal of Algebra} \textbf{424} (2015), 21--45.
\href{https://doi.org/10.1016/j.jalgebra.2014.08.049}{doi:10.1016/j.jalgebra.2014.08.049};
\href{https://arxiv.org/abs/1303.5916}{arXiv:1303.5916}.
本文で用いる Grassmann 記述は Theorem~1。

\bibitem{PMoriSmith}
I.~Mori and S.~P.~Smith,
The classification of 3-Calabi--Yau algebras with 3 generators and 3 quadratic relations,
\href{https://arxiv.org/abs/1502.07403v2}{arXiv:1502.07403v2} (2016).
Theorems~1.3, 1.4, 4.2 and Corollaries~1.6, 2.7, 2.9.
点曲線のHessian公式はLemma~7.1, Theorem~7.2, pp.~23--25。

\bibitem{PBergerTaillefer}
R.~Berger and R.~Taillefer,
Poincar\'e--Birkhoff--Witt deformations of Calabi--Yau algebras,
\emph{Journal of Noncommutative Geometry} \textbf{1} (2007), 241--270.
\href{https://arxiv.org/abs/math/0610112}{arXiv:math/0610112}.
本文で用いる potential による PBW 変形は Theorem~1.1(i)。

\bibitem{PKontsevich}
M.~Kontsevich,
Deformation quantization of Poisson manifolds,
\emph{Letters in Mathematical Physics} \textbf{66} (2003), 157--216.
\href{https://doi.org/10.1023/B:MATH.0000027508.00421.bf}{doi:10.1023/B:MATH.0000027508.00421.bf};
原稿名は \emph{Deformation quantization of Poisson manifolds, I},
\href{https://arxiv.org/abs/q-alg/9709040}{arXiv:q-alg/9709040} (1997).
参照箇所は原稿のSections~4.4--4.6（Maurer--Cartan変形functorとformality）、
6.4（$L_\infty$関係）、8.2および8.3.3.1（中心のtangent mapとwheel）。

\bibitem{PDolgushev}
V.~Dolgushev,
The Van den Bergh duality and the modular symmetry of a Poisson variety,
\href{https://arxiv.org/abs/math/0612288v3}{arXiv:math/0612288v3} (2007).
Theorem~3 および Proposition~1。

\bibitem{BondalGualtieriPym}
M.~Gualtieri and B.~Pym,
\emph{Poisson modules and degeneracy loci},
Proc. Lond. Math. Soc. \textbf{107} (2013), 627--654.
\href{https://doi.org/10.1112/plms/pds090}{doi:10.1112/plms/pds090};
\href{https://arxiv.org/abs/1203.4293}{arXiv:1203.4293}.
Bondal予想の定式化、四次元Fanoの場合の証明。

\bibitem{BondalDruelPereiraPymTouzet}
S.~Druel, J.~V.~Pereira, B.~Pym and F.~Touzet,
\emph{Bondal's conjecture in dimension five},
\href{https://arxiv.org/abs/2606.03424}{arXiv:2606.03424} (2026).
Fano五次元の全場合、一般の奇数次元の部分結果。

\bibitem{FuHwangCI}
B. Fu and J.-M. Hwang,
\emph{Isotrivial VMRT-structures of complete intersection type},
Asian Journal of Mathematics \textbf{22} (2018), 333--354.
Lemma 4.1 (Br\"uckmannの消滅定理を引用).
\url{https://doi.org/10.4310/AJM.2018.v22.n2.a10}.

\bibitem{BorelRemmert}
A. Borel and R. Remmert,
\emph{\"Uber kompakte homogene K\"ahlersche Mannigfaltigkeiten},
Mathematische Annalen \textbf{145} (1962), 429--439, Theorem I.
\url{https://wglobke.github.io/publications/BorelRemmert1962.pdf}.

\bibitem{EG}
P. Etingof and V. Ginzburg,
\emph{Noncommutative del Pezzo surfaces and Calabi--Yau algebras},
Journal of the European Mathematical Society \textbf{12} (2010), 1371--1416.
Sections 3.5--3.6, Theorem 3.6.2, Proposition 8.1.1(i) およびその証明。
\url{https://doi.org/10.4171/JEMS/235}.
\href{https://arxiv.org/abs/0709.3593v4}{arXiv:0709.3593v4}.

\bibitem{LBSVNE12} L. Le Bruyn, S. P. Smith and M. Van den Bergh,
\emph{Central extensions of three dimensional Artin--Schelter regular algebras},
Math. Z. \textbf{222} (1996), 171--212.
Theorem 2.6, Corollary 2.7, Remark 2.8, Theorem 3.1.3.
\url{https://sites.math.washington.edu/~smith/Research/Central-extns.pdf}.

\bibitem{MSNE12} I. Mori and S. P. Smith,
\emph{The classification of 3-Calabi--Yau algebras with 3 generators and 3 quadratic relations},
Math. Z. \textbf{287} (2017), 215--241.
Lemma 4.1, Theorem 4.2, Corollary 4.3; arXiv v2, pp. 15--17.
\url{https://arxiv.org/abs/1502.07403}; doi:10.1007/s00209-016-1824-5.

\bibitem{BergNE12} G. M. Bergman,
\emph{The diamond lemma for ring theory}, Adv. Math. \textbf{29} (1978), 178--218, Theorem 1.2.
\url{https://doi.org/10.1016/0001-8708(78)90010-5}.

\bibitem{BTNE12} R. Berger and R. Taillefer,
\emph{Poincar\'e--Birkhoff--Witt deformations of Calabi--Yau algebras},
J. Noncommut. Geom. \textbf{1} (2007), 241--270.
Theorems 3.1 and 3.6, especially the potential bimodule complex in the proof of Theorem 3.6.
\url{https://arxiv.org/abs/math/0610112}; doi:10.4171/JNCG/6.

\bibitem{SmithTate}
S. P. Smith and J. Tate,
\emph{The center of the 3-dimensional and 4-dimensional Sklyanin algebras},
K-Theory \textbf{8} (1994), 19--63.
Corollary 2.3 (p.~26); Theorem 3.7(c) (p.~38).
\url{https://doi.org/10.1007/BF00962089};
\url{https://sites.math.washington.edu/~smith/Research/tate.pdf}.

\bibitem{DeLaetLeBruyn}
K. De Laet and L. Le Bruyn,
\emph{The geometry of representations of 3-dimensional Sklyanin algebras},
Algebras and Representation Theory \textbf{18} (2015), 761--776.
arXiv:1405.1158v1, Theorem 1 and Section 2 (canonical central cubic formula);
the result is a reformulation of the Artin--Schelter and Smith--Tate formula.
\url{https://arxiv.org/abs/1405.1158}.

\bibitem{ArtebaniDolgachev}
M. Artebani and I. Dolgachev,
\emph{The Hesse pencil of plane cubic curves},
L'Enseignement Math\'ematique (2) \textbf{55} (2009), no.~3--4, 235--273,
Lemma 1 (every smooth plane cubic admits a Hesse canonical form).
\url{https://doi.org/10.4171/LEM/55-3-3};
\url{https://arxiv.org/abs/math/0611590}.

\bibitem{LPTPoissonCubic}
F.~Loray, J.~V.~Pereira and F.~Touzet,
\emph{Foliations with trivial canonical bundle on Fano 3-folds},
Math. Nachr. \textbf{286} (2013), 921--940.
\href{https://arxiv.org/abs/1112.3871}{arXiv:1112.3871}.

\end{thebibliography}

\end{document}
